% !TEX program = xelatex
% 编译方式：在本目录执行两遍 xelatex（详见 docs/_manual_common/README.md）。
% 章节源文件位于 chapters/，由本文件 \input 引入，各章不含导言区。
% 源码依据（行号以 2026-08-17 检出为准）：
%   include/hacdcpf/reliability/reliability_assessment.hpp
%   include/hacdcpf/reliability/failure_mode.hpp
%   include/hacdcpf/analysis/three_stage_reliability.hpp
%   src/reliability/{reliability_assessment,failure_mode,reliability_data,three_stage_reliability}.cpp
\documentclass[UTF8,fontset=fandol,a4paper,12pt]{ctexart}
\usepackage{../../_manual_common/hysim_manual}
\linespread{1.18}
\setlength{\parskip}{0.22em}

\title{\textbf{交直流混合高弹性能源电力系统仿真平台}\\[0.5em]
       {\Large 可靠性评估技术手册}\\[0.3em]
       {\large 数学模型、评估方法、接口参数与验证校核}}
\author{HySim-XJTU-HRPES（西安交通大学 HRPES 团队）}
\date{\today}

\begin{document}
\maketitle
\begin{abstract}
\noindent 本手册依据 HySim-XJTU-HRPES 平台可靠性评估模块
（\srcpath{src/reliability/} 与 \srcpath{include/hacdcpf/reliability/}，
以及位于 \srcpath{include/hacdcpf/analysis/three\_stage\_reliability.hpp} 的三阶段
恢复评估器）整理，覆盖统一可靠性参数解析器、非序贯/序贯蒙特卡洛、元件级 FMEA 与
故障模式级 FMEA、三阶段故障恢复 MILP、频率-持续时间解析法、重要抽样、精确状态灵敏度、
保护配合、智能检测/隔离/恢复、信息物理事件树、序贯信息事件队列与在线 Mass-Matrix DAE 闭环，
以及尾部风险与配电可靠性指标。每个方法给出功能定位、源码等价数学模型、
全部选项与结果字段的单位制和默认值、验证依据及已知局限。文档风格与《潮流计算技术
手册》《最优潮流技术手册》一致，全部公式逐条注明源码出处（文件:函数名锚点，禁止行号）。

\smallskip
\noindent\textbf{核心口径}：蒙特卡洛与 FMEA 在纯交流系统上采用
\emph{AC-only DC-OPF}（\fld{model\_scope}=\texttt{"ac-only-dcopf"}）；含直流网络时采用
\texttt{"hybrid-acdc-network-lp"}，联合计入 AC/DC 有功网络、直流负荷削减及换流器传输。
三阶段模型为耦合交直流 LinDistFlow 恢复 MILP。上述稳态模型均不认证非线性交流无功/电压可行性；
全部近似、拒绝边界与模型覆盖不足均写入结果的 \fld{model\_scope}、
\fld{model\_limitations} 与 \fld{ValidityFlags}。
\end{abstract}

\tableofcontents
\newpage

% 第 1 章：架构总览与方法族
\section{架构总览与方法族}\label{sec:rel-overview}

\subsection{功能定位}
可靠性评估模块回答“系统在元件随机故障下少供多少电、失负荷多频繁多久”。全部方法
遵循同一建模主线（\srcpath{docs/theory/reliability\_mathematical\_models\_and\_intelligent\_cyber\_physical\_assessment.md}
第 1 节的中心原理）：
\begin{equation}
\text{可靠性}=\text{事件频率}\times\text{类条件后果}\times\text{持续时间}.
\end{equation}
命名空间为 \texttt{hacdcpf::analysis}。

\subsection{统一数据流}
\begin{figure}[H]\centering
\begin{tikzpicture}[font=\footnotesize,
  box/.style={draw,rounded corners,minimum height=1.0cm,text width=2.7cm,align=center,inner sep=4pt},
  arr/.style={->,thick}]
  \node[box] (rich) at (0,0) {富语义模型\\\texttt{HybridPowerSystem}};
  \node[box] (res) at (3.7,0) {参数解析\\$\lambda,r,U$ 规范化};
  \node[box] (enum) at (7.4,0) {状态/故障枚举\\抽样·N-1·目录};
  \node[box] (patch) at (11.1,0) {后果算子\\$G_m{=}\Phi_m(G_0,x_m)$};
  \node[box] (shed) at (3.7,-2.5) {切负荷引擎\\DC-OPF·LP·MILP};
  \node[box] (metric) at (7.4,-2.5) {指标聚合\\EENS·LOLE·LOLF};
  \node[box] (audit) at (11.1,-2.5) {口径审计\\\fld{model\_scope}};
  \draw[arr] (rich)--(res); \draw[arr] (res)--(enum); \draw[arr] (enum)--(patch);
  \draw[arr] (patch.south) |- (shed.west);
  \draw[arr] (shed)--(metric); \draw[arr] (metric)--(audit);
\end{tikzpicture}
\caption{可靠性评估统一数据流：解析$\to$枚举$\to$后果映射$\to$切负荷求解$\to$指标$\to$口径审计。}
\end{figure}

\subsection{公共入口族}
全部入口位于命名空间 \texttt{hacdcpf::analysis}；头文件除三阶段模型在
\srcpath{include/hacdcpf/analysis/} 外，均在 \srcpath{include/hacdcpf/reliability/}，下表“头文件”列仅
记基名，函数名与“入口”列同名：
\begin{center}\footnotesize
\begin{tabular}{@{}L{4.9cm}L{4.9cm}L{3.8cm}@{}}
\toprule
入口（函数） & 头文件 & 方法族\\
\midrule
\srcpath{resolve\_reliability\_params} & \srcpath{reliability\_assessment.hpp} & 参数解析（全方法共用）\\
\srcpath{run\_nonsequential\_mc} & \srcpath{reliability\_assessment.hpp} & 非序贯 MC\\
\srcpath{run\_sequential\_mc} & \srcpath{reliability\_assessment.hpp} & 序贯 MC（含 LOLF）\\
\srcpath{run\_distribution\_fmea} & \srcpath{reliability\_assessment.hpp} & 元件级 FMEA\\
\srcpath{run\_failure\_mode\_fmea} & \srcpath{failure\_mode.hpp} & 故障模式级 FMEA（N-2）\\
\srcpath{run\_three\_stage\_reliability} & \srcpath{three\_stage\_reliability.hpp} & 三阶段恢复 MILP\\
\srcpath{run\_frequency\_duration\_analysis} & \srcpath{reliability\_assessment.hpp} & 频率-持续时间解析\\
\srcpath{build\_failure\_mode\_catalog} & \srcpath{failure\_mode.hpp} & 故障模式目录\\
\srcpath{evaluate\_failed\_network\_state} & \srcpath{reliability\_assessment.hpp} & 后果引擎入口\\
\srcpath{compute\_distribution\_indices} & \srcpath{reliability\_assessment.hpp} & SAIFI/SAIDI/ASAI\\
\bottomrule
\end{tabular}
\end{center}

\subsection{物理后果引擎阶梯}
不同方法调用不同保真度的切负荷引擎（\srcpath{docs/theory/reliability\_mathematical\_models\_and\_intelligent\_cyber\_physical\_assessment.md}
第 4.3 节）：
\begin{center}\footnotesize
\begin{tabular}{@{}L{3.6cm}L{4.2cm}L{6.4cm}@{}}
\toprule
引擎 & 使用者 & 有效边界\\
\midrule
AC DC-OPF & 纯交流 MC/FMEA & 无电压幅值/无功可行性（$P=B\theta$）\\
混合 AC/DC 网络 LP & 混合 MC、元件/故障模式 FMEA & DC 传输线性稳态；无非线性 AC 校验\\
耦合 LinDistFlow 恢复 MILP & 三阶段模型 & 线性化损耗；拒绝 LCC 与多端口能量路由器\\
\bottomrule
\end{tabular}
\end{center}

\subsection{诚实结果口径}
每个结果对象都暴露 \fld{model\_scope}、\fld{model\_limitations} 与 \fld{ValidityFlags}，
是数学结果的组成部分，下游 API 必须保留：
\begin{itemize}
  \item MC/FMEA 在纯交流上为 \texttt{"ac-only-dcopf"}，
        \fld{validity.\allowbreak dc\_load\_curtailment\_included}=false——EENS/LOLE 不含 DC 负荷中断；
  \item 混合系统 FMEA 升级为 \texttt{"hybrid-acdc-network-lp"}；
  \item 三阶段模型为
        \texttt{"coupled-acdc-\allowbreak lindistflow-\allowbreak restoration-milp"}，
        含 LCC/能量路由器
        时清 \fld{ok} 并置物理有效性标志为假；
  \item \fld{critical\_components.importance} 是共现归因（可相加超 100\%），\textbf{非} Birnbaum
        边际重要度（\srcpath{reliability\_assessment.hpp} \fld{ReliabilityResult} 注释，:ReliabilityResult）。
\end{itemize}

% ── 通用理论层（修订版 2 新增，工程参考级）──
% theory_reliability_foundations.tex —— 可靠性评估通用理论：元件模型、网络约简、充裕度与蒙特卡洛
% 本文件为理论层章节，遵循 docs/modules/THEORY_WRITING_GUIDE.md。

\section{电力系统可靠性评估的数学基础}\label{sec:rel-th}

\begin{theorynote}
本章为通用数学理论，按电力系统可靠性工程的标准体系（Billinton--Allan~\cite{rel:billinton-allan}）撰写，
覆盖：两状态元件的连续时间马尔可夫（CTMC）可用度模型、串并联网络的频率—持续时间（F\&D）约简与
最小割集、发电充裕度的停运容量概率表（COPT）递推与 LOLP/LOLE/LOLF/LOLD 指标、非序贯与序贯蒙特卡洛
估计量及其无偏性与变异系数（CoV）收敛率、IEEE Std 1366 配电可靠性指标（SAIFI/SAIDI/CAIDI/ASAI）、
以及尾部风险 VaR/CVaR。全部结论以平台实现（\srcpath{src/reliability/}）为落点，规模上限、保真度与
适用范围通过结果范围和有效性标志声明。纯交流系统采用交流网络 DC-OPF 削减口径
（\texttt{model\_scope="ac-only-dcopf"}）；含直流网络的系统采用
\texttt{hybrid-acdc-network-lp}，联合计入 AC/DC 有功平衡、直流负荷削减及 VSC/DC-DC 有界传输。
两种口径均不认证非线性交流无功与电压可行性。
\end{theorynote}

\subsection{记号与约定}\label{sec:rel-th-notation}
\begin{itemize}
  \item $\lambda$：元件故障率（占/年，occ/yr）；$\mu=1/r$：修复率；$r$：平均修复时间 MTTR（yr 或 hr）；
        $m=1/\lambda$：平均无故障时间 MTTF；$T=m+r$：平均循环时间（MTBF）；
  \item $A$：可用度（稳态在运行概率）；$U=1-A$：强迫停运率 / 不可用度（FOR）；$f$：状态频率（occ/yr）；
  \item $H=8760$（或 $8736$）：年小时数；$L$：负荷（MW）；$C$：可用发电容量（MW）；
  \item EENS：期望缺供电量（MWh/yr）；EDNS：期望缺供功率（MW）；
  \item LOLP/LOLE/LOLF/LOLD：失负荷概率/期望（hr/yr）/频率（occ/yr）/持续（hr/occ）；
  \item 主动型元件：$p_d$ 每次需求失效概率、$\nu_d$ 需求频率（次/年），等效年频率 $\lambda_a=\nu_d p_d$。
\end{itemize}
除特别声明，所有元件相互独立，故障/修复时间服从指数分布（马尔可夫假设）。

\subsection{两状态元件模型}\label{sec:rel-th-twostate}\label{sec:rel-param}
可修元件在“运行”（up）与“故障”（down）间交替，构成两状态连续时间马尔可夫链，生成元矩阵
\begin{equation}
 Q=\begin{bmatrix}-\lambda & \lambda\\ \mu & -\mu\end{bmatrix}.
 \label{eq:rel-th-Q}
\end{equation}

\begin{theorem}[稳态可用度与频率]\label{thm:rel-th-avail}
两状态元件~\eqref{eq:rel-th-Q} 的稳态概率、频率为
\begin{equation}
 A=\frac{\mu}{\lambda+\mu}=\frac{m}{m+r},\qquad
 U=\frac{\lambda}{\lambda+\mu}=\frac{\lambda r}{1+\lambda r},\qquad
 f=\lambda A=\mu U=\frac{1}{m+r}.
 \label{eq:rel-th-AUf}
\end{equation}
\end{theorem}

\begin{proof}
稳态分布 $\pi=[A,U]$ 满足 $\pi Q=0$ 与 $A+U=1$。$\pi Q=0$ 的首列给出平衡方程
$-\lambda A+\mu U=0$，即 $\lambda A=\mu U$；代入 $A+U=1$ 解得 $A=\mu/(\lambda+\mu)$、
$U=\lambda/(\lambda+\mu)$。以 $\mu=1/r$ 代入得 $U=\lambda/(\lambda+1/r)=\lambda r/(1+\lambda r)$。
状态频率为“离开该状态的稳态流率”$f=\lambda A=\mu U$；又 $\lambda A=\lambda\mu/(\lambda+\mu)
\allowbreak=1/(1/\lambda+1/\mu)\allowbreak=1/(m+r)$。
\end{proof}

平台以单一\emph{解析器}把异构输入（$\lambda$、MTTR、MTBF、FOR、$p_d/\nu_d$）折算为
式~\eqref{eq:rel-th-AUf} 的规范量 $(\lambda,r,U)$，供所有方法共享，并记录 MTBF 语义与数据来源
（\srcpath{include/hacdcpf/reliability/reliability\_assessment.hpp:resolve\_reliability\_params}）。

\subsection{网络约简：串并联与最小割集}\label{sec:rel-th-network}
独立两状态 CTMC 的串并联子网可用稳态可用度、边界穿越频率和平均故障持续时间逐级合并。下式同时列出
工程中常用的小不可用度（$U_i\ll1$）近似；平台实现保留精确稳态乘积与边界穿越频率，而非只计算一阶式。

\begin{proposition}[串联合并]\label{prop:rel-th-series}
$n$ 个串联元件（任一故障即系统故障），系统三元组为
\begin{equation}
 U_s=1-\prod_i(1-U_i)\approx\sum_iU_i,\qquad
 \lambda_s=\sum_i\lambda_i,\qquad
 r_s=\frac{\sum_i\lambda_ir_i}{\sum_i\lambda_i}.
 \label{eq:rel-th-series}
\end{equation}
\end{proposition}

\begin{proof}
系统可用要求全部元件可用，$A_s=\prod_iA_i$，故 $U_s=1-\prod_i(1-U_i)$，一阶展开得
$\sum_iU_i$。系统故障频率为各元件“单独触发系统故障”的频率之和
$f_s=\sum_i f_i\prod_{j\ne i}A_j\approx\sum_i\lambda_i$（$A_j\approx1$，$f_i\approx\lambda_i$）；
由 $\lambda_s=f_s/A_s\approx f_s$ 得 $\lambda_s=\sum_i\lambda_i$。等效修复时间由 $U_s=\lambda_s r_s$
（小量下 $U\approx\lambda r$）与 $U_s\approx\sum\lambda_ir_i$ 联立得 $r_s=\sum\lambda_ir_i/\sum\lambda_i$。
\end{proof}

\begin{proposition}[双元件并联合并]\label{prop:rel-th-parallel}
两个并联元件（同时故障才系统故障），在 $\lambda_ir_i\ll1$ 下
\begin{equation}
 U_p=U_1U_2\approx\lambda_1r_1\lambda_2r_2,\qquad
 \lambda_p\approx\lambda_1\lambda_2(r_1+r_2),\qquad
 r_p=\frac{r_1r_2}{r_1+r_2}.
 \label{eq:rel-th-parallel}
\end{equation}
\end{proposition}

\begin{proof}
独立性给 $U_p=U_1U_2$。系统进入故障态需一元件先故障（频率 $\approx\lambda_1$ 或 $\lambda_2$）而另一
元件恰处于故障态（概率 $U_2$ 或 $U_1$）：$f_p\approx\lambda_1U_2+\lambda_2U_1
=\lambda_1\lambda_2 r_2+\lambda_2\lambda_1 r_1=\lambda_1\lambda_2(r_1+r_2)$。等效
$\lambda_p=f_p/A_p\approx f_p$。$r_p$ 由 $U_p=\lambda_p r_p$ 得
$r_p=U_1U_2/[\lambda_1\lambda_2(r_1+r_2)]=r_1r_2/(r_1+r_2)$。
\end{proof}

一般（含网络约束、非串并联）物理拓扑用\emph{最小割集}表述：系统故障事件
$=\bigcup_k(\text{割集 }K_k\text{ 全故障})$，一阶近似 $U_{\mathrm{sys}}\approx\sum_k\prod_{i\in K_k}U_i$，
主导项为低阶（一阶、二阶）割集。平台对一般\emph{物理网络}仍以蒙特卡洛状态抽样统计削减（见
\S\ref{sec:rel-th-mc}）；对显式信息服务成功路径则已实现共享元件保持、QoS 筛选、最小割集和精确容斥
（\S\ref{sec:rel-icp-avail}），两者不可混称。

\subsection{发电充裕度：停运容量概率表与 LOLE/LOLF}\label{sec:rel-th-copt}
发电充裕度评估把机组停运卷积成\emph{停运容量概率表}（COPT）。

\begin{proposition}[COPT 递推]\label{prop:rel-th-copt}
设已叠加若干机组得累积停运概率 $P(X)=\Pr\{\text{停运容量}\ge X\}$。再叠加一台容量 $C$、强迫停运率
$U$ 的机组，则
\begin{equation}
 P'(X)=(1-U)\,P(X)+U\,P(X-C).
 \label{eq:rel-th-copt}
\end{equation}
频率表按同样卷积递推（附加 $\pm$ 机组故障/修复率的频率增量）。
\end{proposition}

\begin{proof}
新机组以概率 $1-U$ 在运（停运容量不变，贡献 $P(X)$），以概率 $U$ 停运（停运容量增加 $C$，
把原表右移 $C$，贡献 $P(X-C)$）。两互斥事件相加即~\eqref{eq:rel-th-copt}。
\end{proof}

由 COPT 与负荷 $L$，失负荷概率与期望为
\begin{equation}
 \mathrm{LOLP}=\Pr\{C<L\}=P(\,\text{停运}>C_{\mathrm{inst}}-L\,),\qquad
 \mathrm{LOLE}=\sum_{t}\mathrm{LOLP}(L_t)\,\Delta t,
 \label{eq:rel-th-lole}
\end{equation}
$C_{\mathrm{inst}}$ 为装机容量。F\&D 法进一步给出失负荷频率 $\mathrm{LOLF}$（越过 $L$ 的净向上频率）与
持续 $\mathrm{LOLD}=\mathrm{LOLE}/\mathrm{LOLF}$。平台以 F\&D 结构体
\srcpath{FrequencyDurationResult}（\fld{capacity\_outage\_levels}、\fld{cumulative\_probability}、
\fld{cumulative\_frequency}、\fld{lolp}、\fld{lole\_fd}、\fld{lolf\_fd}、\fld{lold}）承载 COPT 与指标。

\begin{example}[两机组 COPT 与充裕度指标]\label{ex:rel-th-copt}
两台 $60$~MW 机组，强迫停运率 $q=0.05$（$p=0.95$），$\mathrm{MTTF}=1900$~h、$r=100$~h
（校验 $U=r/(\mathrm{MTTF}{+}r)=0.05$，$\lambda=8760/1900=4.61/\text{yr}$、$\mu=8760/100=87.6/\text{yr}$）。
停运容量 $X$ 的概率表为 $P(0){=}p^2{=}0.9025$、$P(60){=}2pq{=}0.095$、$P(120){=}q^2{=}0.0025$。
设峰荷 $L=50$~MW、总容量 $120$~MW，则失负荷当可用容量 $120-X<50$，即 $X\ge120$（两机全停）：
\begin{equation*}
 \mathrm{LOLP}=P(X{\ge}120)=0.0025,\qquad \mathrm{LOLE}=8760\times0.0025=21.9~\text{h/yr}.
\end{equation*}
两机全停态的离开率为 $2\mu=175.2/\text{yr}$，故 $\mathrm{LOLF}=0.0025\times175.2=0.438~\text{occ/yr}$、
$\mathrm{LOLD}=\mathrm{LOLE}/\mathrm{LOLF}=50$~h（$=1/(2\mu)$，两机中任一先修复即结束）。若峰荷升至 $70$~MW，
则 $X{\ge}60$ 即失负荷，$\mathrm{LOLP}=0.095{+}0.0025=0.0975$——凸显充裕度对备用裕度的强敏感性。
\end{example}

\subsection{公共后果模型：突变算子、最小切负荷与引擎阶梯}\label{sec:rel-th-conseq}
失效模式\emph{不}直接等于元件停运，而经\emph{后果算子}映射为网络 / 控制\emph{突变}：
\begin{equation}
 G_m=\Phi_m(G_0,x_m),
\end{equation}
$G_0$ 为健康富系统、$x_m$ 为模式严重度、$\Phi_m$ 产生一或多个突变：拓扑（强迫停运 / 卡开 / 强迫切负荷）、
容量（按残余容量因子降额）、控制（定值冻结 / 失控 / 失去构网）、保护（拒动后区域扩展）、观测（量测 /
通信退化）、恢复（禁止某动作，如拒合）。硬停运在同一目标上压制降额与失控；不兼容的强迫开 / 合突变报为
冲突；所选物理引擎无法表示的突变标为\emph{不支持}而非人为置零。对失效态 $x$，各方法最终求\emph{最小切负荷}
\begin{equation}
 S(x)=\min_{y\in\mathcal F(G_x)}\sum_iw_ip_i^{sh},
 \label{eq:rel-th-minshed}
\end{equation}
$\mathcal F(G_x)$ 为所选物理模型的可行运行 / 恢复域。实现把 VOLL 置于可行发电边际成本之上，使源成本仅在
最小切负荷解间打破简并；可靠性只读 $S$ 与 $s_i$，不读经济目标值。三级\emph{物理后果引擎阶梯}为：
\begin{center}\footnotesize
\begin{tabular}{@{}L{2.9cm}L{3.4cm}L{7.3cm}@{}}
\toprule
引擎 & 使用者 & 主物理 / 有效边界\\
\midrule
交流 DC-OPF & 纯交流 MC/FMEA & $P{=}B\theta$、节点平衡、源 / 支路限、切负荷；无电压幅值 / 无功可行性\\
混合 AC/DC 网络 LP & 混合 MC、元件 / 故障模式 FMEA & AC 角—流、DC 平衡、有界 VSC/DC-DC 传输、储能 / DER；DC 传输线性稳态，无非线性 AC 校验\\
耦合 LinDistFlow 恢复 MILP & 三阶段模型 & AC 有 / 无功 $+$ DC 有功平衡、平方电压、辐射森林、开关、双向换流、DER / 储能时序；线性化损耗，拒 LCC 与多端口路由器\\
\bottomrule
\end{tabular}
\end{center}
依据：\srcpath{src/reliability/reliability\_assessment.cpp:evaluate\_failed\_network\_state}。

\subsection{蒙特卡洛可靠性评估}\label{sec:rel-th-mc}
当网络约束、多重故障与调度重分布使解析枚举不可行时，蒙特卡洛（MC）以状态抽样估计指标。

\paragraph{非序贯（状态抽样）。}
每次抽样对各元件独立取 Bernoulli 状态 $x_i\sim\mathrm{Bern}(U_i)$（$1$=故障），对该系统状态解
DC-OPF 得最小削减 $C(\mathbf x)$（MW），估计
\begin{equation}
 \widehat{\mathrm{EENS}}=\frac1N\sum_{k=1}^N C(\mathbf x^{(k)})\cdot H,\qquad
 \widehat{\mathrm{EDNS}}=\frac1N\sum_{k=1}^N C(\mathbf x^{(k)}).
 \label{eq:rel-th-nsq}
\end{equation}

\begin{theorem}[无偏性与 CoV 收敛]\label{thm:rel-th-cov}
估计量~\eqref{eq:rel-th-nsq} 无偏：$\mathbb E[\widehat{\mathrm{EDNS}}]=\mathbb E[C(\mathbf x)]=\mathrm{EDNS}$；
样本均值的变异系数
\begin{equation}
 \mathrm{CoV}=\frac{\sigma_{\widehat{}}}{\widehat\mu}=\frac{\sigma_C}{\mu_C\sqrt N}=O\!\left(N^{-1/2}\right),
 \label{eq:rel-th-cov}
\end{equation}
$\sigma_C,\mu_C$ 为单样本削减的标准差与均值。停机准则为 $\mathrm{CoV}\le\varepsilon$。
\end{theorem}

\begin{proof}
$\mathbf x^{(k)}$ 独立同分布，$\mathbb E[\tfrac1N\sum_kC(\mathbf x^{(k)})]=\mathbb E[C(\mathbf x)]$，
无偏。样本均值方差 $\mathrm{Var}(\widehat\mu)=\sigma_C^2/N$，标准差 $\sigma_C/\sqrt N$；
除以均值 $\mu_C$ 即 $\mathrm{CoV}=\sigma_C/(\mu_C\sqrt N)$，随 $N$ 以 $N^{-1/2}$ 衰减。
\end{proof}

CoV 的 $O(N^{-1/2})$ 率与\emph{维数无关}，是 MC 相对枚举的根本优势；但稀有事件（$\mu_C$ 小、
$\sigma_C/\mu_C$ 大）收敛慢，需重要抽样等方差缩减。依据：\srcpath{src/reliability/reliability\_assessment.cpp:run\_nonsequential\_mc}
（\fld{cov\_threshold}、\fld{max\_iterations}、\fld{final\_cov}）。

\paragraph{序贯（时序模拟）。}
按时间推进：各元件抽指数分布的运行/修复时长（率 $\lambda_i$/$\mu_i$）生成时序状态，逐时对
含负荷曲线的系统解 DC-OPF，累计年度 EENS 与\emph{越限次数}得 LOLF。序贯法额外给出频率与持续
（LOLF、LOLD）及其年际分布，代价是每样本为完整一年。依据：
\srcpath{src/reliability/reliability\_assessment.cpp:run\_sequential\_mc}、
\srcpath{build\_ieee\_rts24\_load\_profile}。

\subsection{配电可靠性指标（IEEE Std 1366）}\label{sec:rel-th-1366}
以负荷点 $i$ 的故障率 $\lambda_i$、年停电时间 $U_i$（hr/yr）与客户数 $N_i$ 定义系统指标：
\begin{align}
 \mathrm{SAIFI}&=\frac{\sum_i\lambda_iN_i}{\sum_iN_i},&
 \mathrm{SAIDI}&=\frac{\sum_iU_iN_i}{\sum_iN_i},&
 \mathrm{CAIDI}&=\frac{\mathrm{SAIDI}}{\mathrm{SAIFI}},\\
 \mathrm{ASAI}&=1-\frac{\mathrm{SAIDI}}{H},&
 \mathrm{ASUI}&=1-\mathrm{ASAI}.&&
\end{align}
SAIFI 为户均停电次数、SAIDI 为户均停电时长、CAIDI 为次均停电时长、ASAI 为供电可用率。平台以
\srcpath{DistributionIndices}（\fld{saifi}/\fld{saidi}/\fld{caidi}/\fld{asai}/\fld{asui} 及逐节点
\fld{nodal\_cif}/\fld{nodal\_cid}）承载，需负荷点客户数据方计算（\fld{compute\_distribution\_indices}）。

\subsection{尾部风险：VaR 与 CVaR}\label{sec:rel-th-var}
可靠性关注\emph{尾部}的高损失场景。对年损失 $E$（如年 EENS）与置信 $\alpha$，
$\mathrm{VaR}_\alpha(E)=\inf\{e:\Pr[E\le e]\ge\alpha\}$，严格的期望损失（Expected Shortfall）为
\begin{equation}
 \mathrm{ES}_\alpha(E)=\frac1{1-\alpha}\int_\alpha^1\mathrm{VaR}_u(E)\,\mathrm du
 =\min_{\eta\in\mathbb R}\Big\{\eta+\frac1{1-\alpha}\mathbb E[(E-\eta)^+]\Big\},
 \label{eq:rel-th-cvar}
\end{equation}
右式为 Rockafellar--Uryasev 表征~\cite{rel:rockafellar}；连续分布下等于
$\mathbb E[E\mid E\ge\mathrm{VaR}_\alpha]$。CVaR 是\emph{相干}风险度量（次可加、单调、平移不变、正齐次），
优于 VaR。\textbf{实现注}：平台按经验分布对最坏 $1-\alpha$ 概率质量积分；当 VaR 处有并列样本时只取补足
尾概率所需的分数质量，因此与上式的经验 ES 一致，而不是把全部 VaR 原子重复计入。序贯年样本为首选输入。
承载于 \srcpath{TailRiskMetrics}
（\fld{eens\_var}/\fld{eens\_cvar}/\fld{eens\_percentiles}）。

\subsection{采样不确定性与稀有事件}\label{sec:rel-th-uncert}
对 $n$ 个独立状态样本，缺供功率估计的标准误与渐近 $1-\alpha$ 区间为
\begin{equation}
 \mathrm{SE}(\widehat{\mathrm{EDNS}})=\frac{\widehat\sigma_S}{\sqrt n},\qquad
 \bar S\pm z_{1-\alpha/2}\frac{\widehat\sigma_S}{\sqrt n},\qquad
 \mathrm{SE}(\widehat{\mathrm{PLC}})=\sqrt{\frac{\widehat p(1-\widehat p)}{n}}.
\end{equation}
稀有失负荷宜用 Wilson 或精确二项区间；序贯年样本须以\emph{年级}方差 / 区间而非把小时当独立。极小
PLC/LOLE 下朴素 MC 统计低效。平台已实现优势比扭曲重要抽样、似然比无偏修正和 ESS 诊断；交叉熵、子集模拟与
按故障阶分层不属于该接口范围，完整推导见 \S\ref{sec:rel-exact}。这与 CoV 停机准则（定理~\ref{thm:rel-th-cov}）
互补：CoV 度量\emph{相对}收敛，SE / 区间度量\emph{绝对}不确定性。

\subsection{薄弱环节辨识（后处理）}\label{sec:rel-th-weakpoint}
评估后按元件对\emph{失负荷态}的贡献排序，定位薄弱环节。平台的元件重要度定义为
\srcpath{ReliabilityResult::ComponentImportance}，是\emph{共现 / 归因}度量，而非 Birnbaum 边际重要度
$\partial\mathrm{EENS}/\partial U_c$：以“元件 $c$ 停运时观测到的失负荷削减份额”定义
\begin{align}
 \text{loss\_weighted\_risk}_c&=\frac{\sum_k\lambda_k\,\mathbf 1[c\in\text{down}(k)]\,S_k}{\sum_k\lambda_k S_k},\\
 \Pr(c\text{ 停}\mid\text{系统失负荷})&=\frac{\sum_k\lambda_k\,\mathbf 1[c\in\text{down}(k)]\,\mathbf 1[S_k>\varepsilon]}{\sum_k\lambda_k\,\mathbf 1[S_k>\varepsilon]},
\end{align}
并给出共停关联 EENS $\text{associated\_eens}_c$。\textbf{诚实口径}：多重失效态把整段削减\emph{归因于每个}
停运元件，故各元件贡献之和可\emph{超过} 100\%——应读作\emph{排序}而非边际风险。精确状态接口另行计算
$\partial\mathrm{EENS}/\partial U_c$、Birnbaum 与 Fussell--Vesely，见 \S\ref{sec:rel-exact}。依据：\srcpath{src/reliability/reliability\_assessment.cpp}
（\fld{critical\_components}）。

\subsection{数值算例与基准（IEEE RTS-24）}\label{sec:rel-th-numexp}\label{sec:rel-th-state-consequence}
本节先规定 IEEE RTS-24 网络约束可靠性评估应求解的数学对象。该模型是后续代码修订与基准验收的
判据；在全部一致性条件通过前，不以某次蒙特卡洛汇总值反推模型正确。

\subsubsection{物理状态与唯一表示}
令 $\mathcal N$、$\mathcal E$、$\mathcal G$ 分别为 AC 母线、物理支路和机组集合，
$a_e(x),a_g(x)\in\{0,1\}$ 为状态 $x$ 下的可用标志。求解器使用的带电边集必须为
\begin{equation}
 \mathcal E(x)=\{e\in\mathcal E:a_e(x)=1\}.
 \label{eq:rel-rts-edge-set}
\end{equation}
若导入器为物理支路 $e$ 生成参数、控制或结果归因元数据 $t$，并以映射
$\rho(t)=e$ 关联，则 $t$ 不得再产生独立图边或独立随机变量：
\begin{equation}
 a_t(x)\equiv a_{\rho(t)}(x),\qquad
 \mathcal E_{\rm topo}(x)=\mathcal E_{\rm flow}(x)=\mathcal E(x).
 \label{eq:rel-rts-representation-invariant}
\end{equation}
式~\eqref{eq:rel-rts-representation-invariant} 是\emph{表示不变性}，不是数值近似。违反它会使拓扑筛选
与功率平衡看到两个不同系统，任何 EENS 结果均无物理解释。

由 $\mathcal E(x)$ 得连通分量 $\mathcal C(x)=\{C_1,\ldots,C_K\}$。按 MATPOWER DC 模型~\cite{rel:zimmerman}，对每个分量任取一个母线
$r_k\in C_k$ 固定 $\theta_{r_k}=0$，即可消除 DC 节点电纳矩阵的一维零空间：
\begin{equation}
 \operatorname{rank} B_k=|C_k|-1,\qquad \theta_{r_k}=0.
 \label{eq:rel-rts-angle-reference}
\end{equation}
因此 \texttt{SLACK} 最多提供一个可复用的角度规范，不是供电能力；\texttt{NoSlack} 也不等于失电岛。
对式~\eqref{eq:rel-rts-dcopf}，无需判定输入数据中“有没有平衡节点”：每个 $C_k$ 任取一个
$r_k$ 即可。一个没有导入 slack 标签、但含在线可调机组的分量仍可求解；一个保留 slack 标签、但没有
可用注入能力的分量仍会缺供。这里必须区分两类问题：普通 DC 潮流需要指定承担功率不平衡的 slack
注入，而 DC-OPF 已在所有母线同时执行式~\eqref{eq:rel-rts-balance}，功率不平衡由优化变量分配，
角度参考不承担额外注入。

\subsubsection{同一状态上的最小切负荷 DC 模型}
对时段 $h$ 的毛需求 $d_{ih}\ge0$，定义负荷削减 $s_{ih}\in[0,d_{ih}]$、机组出力 $p_{gh}$、
支路潮流 $f_{eh}$ 和相角 $\theta_{ih}$。静态电源、可再生电源、光伏、放电储能以及负的作者负荷
统一聚合为状态相关固定正注入 $\bar p^{\rm fix}_{ih}(x)\ge0$，并定义固定电源削减
$p^{gc}_{ih}$。网络状态后果模型为
\begin{subequations}\label{eq:rel-rts-dcopf}
\begin{align}
 S(x,h)=\min\quad & \sum_{i\in\mathcal N}s_{ih}, \label{eq:rel-rts-dcopf-obj}\\
 \text{s.t.}\quad
&\sum_{g\in\mathcal G_i}p_{gh}+\bar p^{\rm fix}_{ih}(x)-p^{gc}_{ih}
  +s_{ih}-d_{ih}
   =\sum_{e\in\delta(i)}K_{ie}f_{eh}, &&i\in\mathcal N,\label{eq:rel-rts-balance}\\
&f_{eh}=a_e(x)b_e(\theta_{o(e)h}-\theta_{d(e)h}), &&e\in\mathcal E,\label{eq:rel-rts-flow}\\
&-a_e(x)\bar f_e\le f_{eh}\le a_e(x)\bar f_e, &&e\in\mathcal E,\label{eq:rel-rts-limit}\\
&0\le s_{ih}\le d_{ih},\quad
  0\le p^{gc}_{ih}\le\bar p^{\rm fix}_{ih}(x), &&i\in\mathcal N,\label{eq:rel-rts-shed}\\
 &\theta_{r_kh}=0, &&C_k\in\mathcal C(x).\label{eq:rel-rts-ref}
\end{align}
\end{subequations}
其中 $K$ 为支路—节点关联矩阵，$b_e$ 为 DC 电纳。当前可靠性后果模型采用等权 MW 削减，且与
平台 DC-OPF 一致，不计支路移相角。若分量 $C_k$ 没有任何可用供给，对式~\eqref{eq:rel-rts-balance} 在 $C_k$ 上求和，边界潮流
相消，立即得到
\begin{equation}
 \sum_{i\in C_k}s_{ih}=\sum_{i\in C_k}d_{ih}.
 \label{eq:rel-rts-dead-island}
\end{equation}
所以无源岛的“全负荷切除”应由同一个优化模型自然产生，不需要以 slack 标签或设备类型作启发式回退。
更强地说，在采用 HL-II 再调度语义~\eqref{eq:rel-rts-hlii-redispatch} 时，
\begin{equation}
 p=0,\qquad f=0,\qquad \theta=0,\qquad s=d,\qquad
 p^{gc}=\bar p^{\rm fix}
 \label{eq:rel-rts-constructive-point}
\end{equation}
总是一个可行点；即使整个系统没有在线机组，或某孤岛固定注入大于负荷，模型也应返回可行最优解，而不是以
$n_g=0$、供给过剩或 \texttt{NoSlack} 返回求解失败。固定注入与毛负荷必须分开聚合；若先合成净负荷再令
$s_i\le\max(d_i-\bar p_i^{\rm fix},0)$，式~\eqref{eq:rel-rts-constructive-point} 将被错误破坏。

目标~\eqref{eq:rel-rts-dcopf-obj} 按字典序实现。令 $S^\star$ 为第一级最优值，第二级求解
\begin{equation}
 \min\ \sum_g C_g(p_g)+\sum_i\gamma_i p^{gc}_i
 \quad\text{s.t. 式~\eqref{eq:rel-rts-dcopf} 且 }\sum_i s_i=S^\star,
 \label{eq:rel-rts-lexicographic-second}
\end{equation}
其中 $\gamma_i\ge0$ 为固定电源削减代价。单一有限 VOLL 加权只有在给出目标界和求解容差时才能证明等价；否则
“VOLL 很大”只是工程近似，必须在结果有效性中声明。

\subsubsection{机组下界的三种运行语义}
机组约束不能在未声明评估层级时统一改写。至少有以下三个互不等价的模型：

\paragraph{A. HL-I 发电充裕度。}
忽略网络，仅比较可用容量与负荷：
\begin{equation}
 0\le p_{gh}\le a_g(x)\bar P_g.
 \label{eq:rel-rts-hli-dispatch}
\end{equation}
这与 COPT 的“在线容量可从 0 使用到 $\bar P_g$”一致，不消费 MATPOWER 的经济调度 $P_g^{\min}$。

\paragraph{B. HL-II 网络约束充裕度。}
保留式~\eqref{eq:rel-rts-dcopf} 的网络约束，并仍采用
\begin{equation}
 0\le p_{gh}\le a_g(x)\bar P_g.
 \label{eq:rel-rts-hlii-redispatch}
\end{equation}
该口径假设故障后允许充裕度再调度，但不模拟开停机、最小开停时间与爬坡。它适合把网络约束增量与
经典 RTS COPT 比较，但不能冒充生产模拟。

\paragraph{C. 固定/时序开机组合的运行可靠性。}
若开机状态 $y_{gh}\in\{0,1\}$ 来自 UC 或生产模拟，则
\begin{equation}
 a_g(x)y_{gh}P_g^{\min}\le p_{gh}\le a_g(x)y_{gh}\bar P_g.
 \label{eq:rel-rts-commitment}
\end{equation}
此时低负荷岛可能满足 $\sum_gP_g^{\min}>\sum_i d_i$。模型必须明确选择以下机制之一：消费 UC 的
紧急停机决策、加入有物理含义的弃电/外送变量、或把该状态报告为\emph{运行约束不可行}；不得把这种
供给过剩转换成负荷不足。没有 $y_{gh}$、爬坡、最小开停时间和弃电语义时，直接保留正 $P_g^{\min}$
或直接把全部 $P_g^{\min}$ 清零都不能称为通用正确模型。

\begin{center}\footnotesize
\begin{tabular}{@{}L{3.2cm}L{3.7cm}L{3.7cm}L{3.7cm}@{}}
\toprule
口径 & 机组下界 & 可与经典 RTS 比较 & 必须报告的限制\\
\midrule
HL-I COPT & $0$ & 是，发电充裕度 & 无网络\\
HL-II 再调度 & $0$ & 可比较网络约束增量 & 无 UC/爬坡/电压无功\\
固定/时序组合 & $a_gyP^{\min}$ & 仅与同一 UC 轨迹比较 & 需弃电或紧急停机语义\\
\bottomrule
\end{tabular}
\end{center}

\subsubsection{物理后果、基线缺额与数值失败必须分离}
对求解成功状态，定义原始物理削减 $S^{\rm raw}(x,h)$。同负荷下健康状态为 $x_0$，则
\begin{equation}
 S^0(h)=S^{\rm raw}(x_0,h),\qquad
 S^{\rm inc}(x,h)=\max\{0,S^{\rm raw}(x,h)-S^0(h)\}.
 \label{eq:rel-rts-raw-incremental}
\end{equation}
标准 EENS 使用 $S^{\rm raw}$；$S^{\rm inc}$ 是故障增量诊断。若 $S^0(h)>0$，应判定基准系统、负荷
曲线或运行语义未闭环，并同时报告原始量与增量量，不能静默用增量量替换标准指标。

令 $q(x,h)=1$ 表示优化器取得满足可行性容差的解。数值失败 $q=0$ 不是物理状态，故主估计量只能对
已认证状态定义：
\begin{equation}
 \widehat{\mathrm{EENS}}_{\rm cert}
 =\frac{H}{N}\sum_{k:q_k=1}S^{\rm raw}(x^{(k)},h_k),\qquad
 p_{\rm fail}=\frac1N\sum_k\mathbf1[q_k=0].
 \label{eq:rel-rts-certified-eens}
\end{equation}
若为失败状态临时取 $S^{\rm UB}=\sum_i d_i$，它只能形成保守上界
\begin{equation}
 \widehat{\mathrm{EENS}}^{\rm UB}
 =\widehat{\mathrm{EENS}}_{\rm cert}
 +\frac{H}{N}\sum_{k:q_k=0}S^{\rm UB}_k,
 \label{eq:rel-rts-eens-upper-bound}
\end{equation}
不得把式~\eqref{eq:rel-rts-eens-upper-bound} 标成普通 EENS。每个失败状态至少记录：状态键、连通分量、
在线源与上下界、求解器状态、原始/对偶残差及后端链。只有证明失败源于模型物理不可行，且模型已包含
式~\eqref{eq:rel-rts-shed} 所需的松弛后，才能把它解释为系统后果。

\begin{theorynote}
本节契约已在 AC-only HL-II 与混合 AC/DC 有功网络两条代码路径闭环实现。AC-only 入口强制开启
毛负荷削减、固定正注入削减、$P_g^{\min}=0$ 和
两阶段字典序目标：第一阶段最小化 $\sum_i s_i$，第二阶段加入 $\sum_i s_i=S^\star$ 等式后
最小化 PWL 发电成本与 $p^{gc}$ 代价。混合 LP 同样执行两阶段目标，VSC/DC--DC 紧急再调度区间包含
零传输。两条路径均使用唯一物理边集和每个 AC 连通分量一个角度参考。DC-OPF 在规范投影坐标中对
等式、不等式、盒界和有限性作后验认证；若后端失败或证书超差，评估立即抛出含残差与后端链的错误，
该状态不写入 EENS。序贯 MC 的 \fld{annual\_eens} 与 \fld{eens\_mwh\_yr} 逐小时累加
$S^{\rm raw}$，\fld{baseline\_eens\_mwh\_yr} 单列 $S^0$，
\fld{incremental\_eens\_mwh\_yr} 单列 $S^{\rm inc}$；LOLE/LOLF 也按原始失负荷标志计算。
因此健康状态已有缺额时，不会再被静默扣除出标准指标。
\end{theorynote}

\subsubsection{诊断数据与方法对比}
在 IEEE 可靠性测试系统 RTS-24（$24$ 母线、$33$ 机组、$38$ 支路；\srcpath{data/case24\_ieee\_rts.m}
$+$ \srcpath{apply\_ieee24\_reliability\_data}）上运行两类 MC（\srcpath{tools/pf\_doc\_benchmark.cpp}
模式 \texttt{rel}，\texttt{ac-only-dcopf} 口径）：
可靠性数据按 \texttt{case24\_ieee\_rts.m} 的 33 行机组和 38 行支路逐行映射，并以
母线号、$P_g^{\max}$ 和支路端点签名校验。第 15 行 $P_g^{\max}=0$ 是同步调相机，仍保留其
记录身份但不提供有功容量。旧实现按母线给同一母线全部机组复制一个 FOR/MTTR，并把并联支路按
端点合并匹配；这会改变 COPT 和网络状态概率，是结果偏高/偏低均可能发生的数据构模错误。
\begin{center}
\resizebox{\textwidth}{!}{%
\begin{tabular}{@{}lccccc@{}}
\toprule
层级/方法 & 负荷口径 & CoV/年数 & EENS/(MWh$\cdot$yr$^{-1}$) & LOLE/(hr$\cdot$yr$^{-1}$) & 诊断\\
\midrule
HL-I 当前机组映射（COPT） & 峰值 $2850$~MW & 精确卷积 & 仅容量缺额 & 见 FD & 不含网络故障\\
非序贯 MC（逐台数据，诊断） & 固定峰值、8760 h & $0.0314/2\times10^4$ 态 & $123\,156.8$ & $738.5$ & 未达到停止条件\\
序贯 MC（逐台数据，诊断） & 8736 h 曲线 & $0.1152/400$ 年 & $1\,222.6$ & $10.54$ & 未达到停止条件\\
确定性 HL-II 扫描 & N-0/N-1/N-2 & $2557/2557$ 态 & -- & -- & 失败 0、残差 0\\
\bottomrule
\end{tabular}
}
\end{center}
两类 MC 数值不作为模型可行性证明：非序贯使用固定峰值，序贯使用时序曲线，且本次运行的 CoV
尚未达到停止阈值。当前模型先用确定性扫描证明每个 N-0/N-1/N-2 状态均可认证，再把 MC
作为概率估计。修正后每小时采用
\begin{equation}
 S^{\mathrm{inc}}_{h}=\max\{0,S(x_h,L_h)-S(x_0,L_h)\},\qquad
 \mathrm{EENS}^{\mathrm{inc}}=\frac1{N_y}\sum_y\sum_h S^{\mathrm{inc}}_{y,h}.
\end{equation}
其中 $x_0$ 是同一负荷倍率下的 N-0 状态。该量按式~\eqref{eq:rel-rts-raw-incremental} 只能称为
故障增量 EENS；标准原始 EENS 仍须另报。上表两次 MC 均为诊断运行，不能当作最终收敛基准；
最终基准必须在预设 CoV 停止条件满足且失败计数为零后发布。
这里的 MC \texttt{converged} 仅表示统计误差停止条件；它不能覆盖状态级求解证书。后者若为
\texttt{converged=false} 或后验可行性失败，属于构模/数值求解缺陷，绝不作为“尚未抽够样本”处理，
也不允许继续产生可靠性指标。
\begin{boundarynote}
经典 RTS-24 的 EENS $\sim1.2\times10^3$~MWh/yr、LOLE $\sim9.4$~hr/yr 通常是 HL-I 发电充裕度口径，仅对机组 COPT 卷积，不抽样输电支路。发电与网络风险一般重叠，不能直接作相互独立的加法分解。应在相同状态与负荷上计算反事实增量
\begin{equation}
 \Delta\mathrm{EENS}_{\rm network}
 =H\,\mathbb E\!\left[S_{\rm HL\text{-}II}(x,h)-S_{\rm HL\text{-}I}(x,h)\right],
 \label{eq:rel-rts-network-increment}
\end{equation}
两模型采用相同机组可用状态和运行语义时，HL-II 可行域是 HL-I 可行域加网络限制后的子集，故该增量应非负。
无源岛削减与其余连通区域的 OPF 削减可按\emph{互斥母线集合}相加，但它们已经包含在
$S_{\rm HL\text{-}II}$ 内，不能再次叠加。式~\eqref{eq:rel-rts-eens-upper-bound} 只保留为历史
诊断定义；当前执行契约不接受任何数值失败状态，不能把失败状态的上界写入 EENS，也不能混入
物理网络风险。结果结构现已报告
\texttt{opf\_failed\_probability}、\texttt{dead\_island\_eens\_mwh\_yr} 和
\texttt{opf\_failed\_eens\_mwh\_yr}。当前闭环模型要求两个 \texttt{opf\_failed} 字段严格为零；
\texttt{dead\_island\_eens} 是求解成功后的物理归因，可以非零。历史失败回退值只保留为缺陷定位记录，
不能称为 RTS 物理风险；若状态求解返回 \texttt{converged=false} 或后验证书失败，入口立即抛出包含
状态键、求解器状态和残差的错误，调用方必须修复构模或数值算法后重新运行。

产生回退的具体原因是 MATPOWER 变压器同时存在两种表示：带变比支路保留在 \texttt{ac.branches}，并额外生成带
\texttt{source\_branch\_idx} 的 \texttt{Transformer2W} 元数据。旧拓扑图把两者都当成连通边，而 DC-OPF 只把 \texttt{ac.branches}
装入支路约束。因此旧实现中“变压器支路停运、元数据边仍在”的状态会被图判为有源连通岛，随后
OPF 面对实际断开的 $B$ 矩阵而失败；旧失败分支又把存活负荷全部计为削减。现行拓扑、AC-only DC-OPF、
混合 AC/DC LP 和随机元件集合均跳过链接元数据边，且失败状态直接抛错，不再生成 EENS 数值。

该故障链可写成同一状态向量上的两个不同边集：
\begin{equation}
 E_{\mathrm{graph}}=E_{\mathrm{branch}}\cup E_{\mathrm{transformer\ metadata}},\qquad
 E_{\mathrm{opf}}=E_{\mathrm{branch}} .
\end{equation}
若对应原始支路状态为 $x_b=1$ 而元数据仍为在线，旧实现得到
$G(E_{\mathrm{graph}})$ 连通、但 $B(E_{\mathrm{opf}})$ 已断开；拓扑预筛因此不切除该负荷，
而 OPF 无法为断开的实际分量建立可行功率平衡，触发“存活负荷全部削减”的回退。反向状态
（元数据停运而原始支路在线）则会采到一个不对应任何物理故障的随机事件。两种情况分别造成
假阴性和假阳性，不能用“保守模型”解释。现行图、OPF 和 MC 元件集合均跳过
\texttt{source\_branch\_idx>0} 的链接元数据，因此该双重表示不再产生独立随机状态。

表示层修订采用的不变量是：对任意链接元数据行 $t$，只要 $t.\texttt{source\_branch\_idx}>0$，它只能参与
结果归因和映射，不能参与图边、孤岛边或独立蒙特卡洛元件集合；其在线状态完全由对应
\texttt{ACBranch} 决定。该不变量由
\srcpath{tests/test\_reliability\_resolver.cpp} 中的两母线回归测试执行验证：健康状态拓扑连通，
停运原始支路后恰为两个 AC 岛、负荷岛状态为 \texttt{NoSlack}，且 DC-OPF 返回 10 MW 的可审计切负荷，
不再借助元数据边维持连通。自适应孤岛子系统也采用同一过滤规则（\srcpath{src/power\_flow/island\_detector.cpp}）。

同类但独立的错误来自把母线 \texttt{SLACK} 类型误作在线发电能力。进一步按
式~\eqref{eq:rel-rts-balance}--\eqref{eq:rel-rts-ref} 推导可知，修复并不是把判据改成“岛内是否有
在线源”，而是完全取消 DC-OPF 的平衡节点/在线源可行性预筛：无 slack 有源岛由局部角度参考求解，
有 slack 无源岛和无 slack 无源岛都由 $s=d$ 内生得到全切负荷。现行实现已经删除该在线源预筛，
并以无机组、固定注入过剩、断岛和链接变压器四类回归用例约束这一行为。
\end{boundarynote}

\subsection{与实现的对应}\label{sec:rel-th-impl}
\begin{center}\footnotesize
\begin{longtable}{@{}L{5.0cm}L{9.4cm}@{}}
\toprule
\textbf{理论对象} & \textbf{平台实现状态}\\
\midrule
\endfirsthead
\toprule
\textbf{理论对象} & \textbf{平台实现状态}\\
\midrule
\endhead
\bottomrule
\endfoot
两状态可用度~\eqref{eq:rel-th-AUf}、参数解析 &
  \implfull{src/reliability/reliability\_assessment.cpp:resolve\_reliability\_params}\\
非序贯（状态抽样）MC~\eqref{eq:rel-th-nsq}、CoV 停机 &
  \implfull{src/reliability/reliability\_assessment.cpp:run\_nonsequential\_mc}\\
序贯（时序）MC、LOLF/年际分布 &
  \implfull{src/reliability/reliability\_assessment.cpp:run\_sequential\_mc}\\
COPT 递推~\eqref{eq:rel-th-copt}、LOLP/LOLE/LOLF/LOLD &
  \implfull{src/reliability/reliability\_assessment.cpp:run\_frequency\_duration\_analysis}（独立两状态机组、恒定峰荷；容量不分箱，缺 MTTR 时不伪造 LOLF/LOLD）\\
串并联解析约简~\eqref{eq:rel-th-series}--\eqref{eq:rel-th-parallel} &
  \implfull{src/reliability/reliability\_assessment.cpp:reduce\_series\_frequency\_duration / reduce\_parallel\_frequency\_duration}（独立两状态 CTMC）\\
一般物理网络最小割集解析枚举 &
  \implfull{src/reliability/reliability\_assessment.cpp:evaluate\_physical\_network\_reliability}（成功路径容斥精确可用率、极小击中集、稳定 ID 输出）\\
IEEE 1366 配电指标 &
  \implfull{include/hacdcpf/reliability/reliability\_assessment.hpp:DistributionIndices}（需客户数据）\\
VaR/CVaR 尾部风险 &
  \implfull{include/hacdcpf/reliability/reliability\_assessment.hpp:TailRiskMetrics}（经验 ES 对 VaR 原子作分数权重）\\
混合 AC/DC 有功网络与直流侧削减 &
  \implfull{src/reliability/reliability\_assessment.cpp:evaluate\_hybrid\_fmea\_network\_lp}（\texttt{hybrid-acdc-network-lp}）\\
网络约束状态后果模型~\eqref{eq:rel-rts-edge-set}--\eqref{eq:rel-rts-eens-upper-bound} &
  \implfull{src/optimal\_power\_flow/dc\_opf.cpp:solve\_dc\_opf 与 src/reliability/reliability\_assessment.cpp:evaluate\_state / evaluate\_hybrid\_fmea\_network\_lp}（唯一物理边、逐分量角度参考、$P_g^{\min}=0$、$p^{gc}$、两阶段字典序及规范坐标后验可行性证书）\\
\end{longtable}
\end{center}

\subsection{参考文献}
\begin{thebibliography}{9}\small
\bibitem{rel:billinton-allan} R.~Billinton and R.~N. Allan, \emph{Reliability Evaluation of Power
  Systems}, 2nd ed. New York: Plenum Press, 1996.
\bibitem{rel:billinton-li} R.~Billinton and W.~Li, \emph{Reliability Assessment of Electric Power
  Systems Using Monte Carlo Methods}. New York: Plenum Press, 1994.
\bibitem{rel:rts79} IEEE APM Subcommittee, ``IEEE Reliability Test System,'' \emph{IEEE Trans. Power
  App. Syst.}, vol.~PAS-98, no.~6, pp.~2047--2054, 1979.
\bibitem{rel:rts96} C.~Grigg \emph{et al.}, ``The IEEE Reliability Test System--1996,'' \emph{IEEE
  Trans. Power Syst.}, vol.~14, no.~3, pp.~1010--1020, 1999.
\bibitem{rel:ieee1366} IEEE Std 1366-2012, \emph{IEEE Guide for Electric Power Distribution
  Reliability Indices}. New York: IEEE, 2012.
\bibitem{rel:rockafellar} R.~T. Rockafellar and S.~Uryasev, ``Optimization of conditional
  value-at-risk,'' \emph{Journal of Risk}, vol.~2, no.~3, pp.~21--41, 2000.
\bibitem{rel:zimmerman} R.~D. Zimmerman, C.~E. Murillo-S\'anchez, and R.~J. Thomas,
  ``MATPOWER: Steady-state operations, planning, and analysis tools for power systems
  research and education,'' \emph{IEEE Trans. Power Syst.}, vol.~26, no.~1, pp.~12--19, 2011.
\end{thebibliography}

% physical_consequence_models.tex -- 物理后果模型独立审计章
% 理论先于实现：本章给出主网、配网、运行拓扑重构和三阶段恢复的统一可行域。

\section{物理后果模型与可行性审计}
\label{sec:rel-physical-consequence}

\begin{theorynote}
本章把“物理后果”从可靠性指标聚合中分离出来。故障事件只改变设备状态、保护隔离状态、可用源、
开关候选和阶段时钟；后果算子在给定状态上求解一个带有\emph{显式负荷削减}的最小后果问题。
主网采用网络约束 HL-II DC-OPF，含直流设备时采用耦合 AC/DC 有功网络 LP；配网采用带电状态、
LinDistFlow、电压、热限、辐射森林和运行拓扑重构的三阶段 MILP。所有后果模型都必须包含
$0\le s_i\le d_i$，所以“全切负荷”是一个构造性可行点，而不是数值求解失败的替代值。
任何求解器失败、后验残差超差、整数性失败或结果非有限，均属于构模/数值缺陷：该状态立即中止，
不得把全切负荷写入 EENS、LOLE、SAIDI 或 SAIFI。
\end{theorynote}

\subsection{统一的状态—后果算子}
设原始设备模型为 $G_0$，故障/保护/信息状态为 $x$，阶段为 $s$，外生负荷和可再生可用性为 $h$。
状态投影先生成唯一的物理边集、源集和开关候选：
\begin{equation}
 (\mathcal E_x,\mathcal G_x,\mathcal L_x,\mathcal W_x)
 =\Pi(G_0,x,h,s).
 \label{eq:physical-projection}
\end{equation}
后果算子定义为
\begin{equation}
 S_{x,s}(h)=\min_{u\in\mathcal F(G_0,x,h,s)}\sum_{i\in\mathcal N}w_i s_i,
 \qquad 0\le s_i\le d_i(h).
 \label{eq:physical-consequence-operator}
\end{equation}
其中 $u$ 包含发电、潮流、电压、储能、换流器、带电状态和开关状态。$S_{x,s}$ 是物理后果；
求解器状态、迭代次数和最优间隙是认证元数据，不能被当成后果变量。

对每个状态至少需要四个不变量：
\begin{enumerate}
  \item 拓扑边集与潮流边集相同；链接元数据只能做映射和归因，不生成第二条电气边；
  \item AC、DC 母线分别在各自域内编号，稳定组件 ID、vector position 和图索引不混用；
  \item 故障元件在规定的隔离/修复边界前保持停运，不能被 Stage 3 的重构悄悄重新合闸；
  \item 所有正负荷都存在有界削减变量，固定正注入存在有界弃电变量或等价的固定注入削减。
\end{enumerate}

\subsection{主网：HL-II DC-OPF 后果模型}
令 $a_e(x),a_g(x)\in\{0,1\}$ 表示支路和机组可用状态，$d_i$ 为毛负荷，$\bar p_i^{fix}$ 为固定正注入。
对每个 AC 连通分量 $C_k$ 任取角度参考 $r_k$，不以输入数据中的 \texttt{SLACK} 标签作为可行性判据：
\begin{subequations}\label{eq:physical-main-dcopf}
\begin{align}
 \min\quad & \sum_i s_i,\\
 \sum_{g\in\mathcal G_i}p_g+\bar p_i^{fix}-p_i^{gc}+s_i-d_i
 &=\sum_{e\in\delta(i)}K_{ie}f_e, &&i\in\mathcal N,\\
 f_e&=a_e b_e(\theta_{o(e)}-\theta_{d(e)}),&&e\in\mathcal E,\\
 -a_e\bar f_e\le f_e&\le a_e\bar f_e,&&e\in\mathcal E,\\
 0\le p_g&\le a_g\bar p_g,\quad 0\le s_i\le d_i,\quad
 0\le p_i^{gc}\le\bar p_i^{fix},\\
 \theta_{r_k}&=0,&&C_k\in\mathcal C(x).
\end{align}
\end{subequations}
在线机组下界为零，表示 HL-II 故障后允许再调度；这不是 UC 运行语义。第一阶段最小化总削减得到
$S^\star$，第二阶段加入 $\sum_i s_i=S^\star$，再最小化发电成本和固定电源削减代价。
因此在线源为零、孤岛无 slack、固定注入过剩和断开支路都具有显式可行点：
\begin{equation}
 p=0,\ f=0,\ \theta=0,\ s=d,\ p^{gc}=\bar p^{fix}.
 \label{eq:physical-main-feasible-point}
\end{equation}
若该模型返回 \texttt{converged=false}，不是物理后果，而是模型或数值求解器缺陷。

\subsection{混合交直流主网后果}
混合系统的 AC 支路采用同一角度—潮流方程，DC 网络采用节点有功平衡、平方电压降和支路限额；VSC、
DC--DC 变换器用有界双向传输变量耦合两个域。对 VSC $v$：
\begin{equation}
 P_v^{ac}=-p_v^++\eta_vp_v^-,\qquad
 P_v^{dc}=\eta_vp_v^+-p_v^-,\qquad 0\le p_v^+,p_v^-\le\bar p_v.
 \label{eq:physical-vsc-transfer}
\end{equation}
DC 负荷同样有 $0\le s_i^{dc}\le d_i^{dc}$，DC 固定源有界削减。混合 LP 的可行性构造为所有可调源、
潮流、换流器传输取零，AC/DC 负荷全部削减，固定注入全部削减；因此增加 DC 域不会把一个原本可行的
负荷削减模型变成“无源即不可行”。LCC 和多端口能量路由器若未提供相应换相/端口方程，入口必须拒绝，
不得返回零影响结果。

\subsection{配网：单阶段运行后果}
配网把“可带电”与“可注入”分开。$y_i\in\{0,1\}$ 表示母线带电，$z_{ij}\in\{0,1\}$ 表示线路闭合，
$p_i^{sh},q_i^{sh}$ 为显式的有功、无功削减：
\begin{align}
 0\le p_i^{sh}&\le P_i^d,\qquad
 \min(0,Q_i^d)\le q_i^{sh}&\le\max(0,Q_i^d),\\
 q_i^{sh}&=\frac{Q_i^d}{P_i^d}p_i^{sh}\quad(P_i^d>0),\\
 0\le p_g&\le \bar P_g y_{i(g)},\qquad
 0\le p_b^{st}\le\bar P_b^{st}y_{i(b)},\\
 P_i^d-p_i^{sh}&\le P_i^d y_i,\\
 q_i^{sh}&\ge Q_i^d(1-y_i)\quad(P_i^d=0,Q_i^d>0),\\
 q_i^{sh}&\le Q_i^d(1-y_i)\quad(P_i^d=0,Q_i^d<0).
 \end{align}
纯无功负荷不能被丢弃或强行套用未定义的 $Q_i^d/P_i^d$。当前目标在最小有功削减后，以不超过
$10^{-8}$ 的归一化项最小化 $P_i^d=0$ 母线的 $|q_i^{sh}|$，只消除无功削减退化，不改变有功 EENS 口径。
交流 LinDistFlow 约束为
\begin{align}
 \sum_{j:(j,i)}P_{ji}-\sum_{j:(i,j)}P_{ij}+p_i^g+p_i^{st}+p_i^{sh}&=P_i^d,\\
 \sum_{j:(j,i)}Q_{ji}-\sum_{j:(i,j)}Q_{ij}+q_i^g+q_i^{sh}&=Q_i^d,\\
 -M(1-z_{ij})\le v_j-v_i+2(r_{ij}P_{ij}+x_{ij}Q_{ij})&\le M(1-z_{ij}),\\
 \underline V_i^2\le v_i&\le\overline V_i^2,\\
 |P_{ij}\cos\phi+Q_{ij}\sin\phi|&\le z_{ij}\bar S_{ij}\cos(\pi/J).
 \end{align}
视在功率约束使用额定圆的内接 $J$ 边形，因此模型可行点不会超热限；$J=16$ 是当前默认值。

\subsection{配网三阶段与运行拓扑重构}
对故障 $k$，阶段时长为 $\Delta t_1=\tau_{iso}$、$\Delta t_2=\tau_{sw}$、$\Delta t_3=\tau_{rep}$。
故障元件状态满足
\begin{equation}
 z_{k,s}=0\quad(s=1,2,3),\qquad z_{\ell,2}\in\{0,1\},
 \qquad z_{\ell,3}=z_{\ell,2}\quad(\ell\text{ 为已接受联络}).
 \label{eq:physical-three-stage-topology}
\end{equation}
Stage 1 固定保护隔离后的拓扑；Stage 2 在可操作联络和开关预算内重构；Stage 3 在修复窗保持 Stage-2
拓扑，故障元件在 $\tau_{rep}$ 结束前不恢复。若储能剩余能量、阶段时长或其他物理边界变化，Stage 3
固定 Stage-2 开关变量后重新求解；若这些数据与 Stage 2 完全相同，则 Stage-2 的已认证最优解及最优值可被
精确复用。这是同一优化问题的复用，不是未校核地复制一个后果值。

运行拓扑重构的辐射森林采用有向参考下可正可负的商品流 $F_{ij}$ 和根集合 $\mathcal R$。记
$\kappa$ 为候选健康图中含根的连通分量数；当前实现的约束为：
\begin{align}
 \sum_{j:(j,i)}F_{ji}-\sum_{j:(i,j)}F_{ij}&=y_i, &&i\notin\mathcal R,\\
 \sum_{j:(j,r)}F_{jr}-\sum_{j:(r,j)}F_{rj}&\le 0, &&r\in\mathcal R,\\
 |F_{ij}|&\le (|\mathcal N|-1)z_{ij},\\
 z_{ij}&\le y_i,\quad z_{ij}\le y_j,\\
 y_i+y_j-z_{ij}&\le 1, &&(i,j)\text{ 为健康常闭边},\\
 \sum_{(i,j)}z_{ij}&\le\sum_i y_i-\kappa.
 \end{align}
商品流保证每个带电非根母线由某个根可达，边数不等式排除多余环；健康常闭边在两端均带电时不得被
无记录地打开，正常开断的联络边才由 Stage 2 决定是否闭合。VSC/DC--DC 是跨域能量耦合，不计入 AC
或 DC 辐射树基数；当前完整物理口径对 AC、DC 两域均执行辐射森林约束。

储能在阶段间按能量守恒推进：
\begin{equation}
 E_{b,s+1}=E_{b,s}-p_{b,s}^{st}\Delta t_s,\qquad
 0\le p_{b,s}^{st}\le\min(\bar P_b, E_{b,s}/\Delta t_s).
 \label{eq:physical-storage-chronology}
\end{equation}
健康 N-0 反事实从同一初始能量独立演化；故障路径与健康路径不能共享已耗尽的电池状态。

\subsection{任意故障的构造性可行性定理}
\begin{theorynote}
\textbf{定理（负荷削减可行性）}：若每个模型包含有限的有功/无功负荷削减上界、固定正注入削减上界，所有电压
允许区间非空且包含所选参考值，Big-M 足以在开边时解耦端电压；对配网还要求强制带电根及不可操作的
健康常闭边能够形成根森林（或已提供可开断分段设备使其可形成根森林）。若故障只通过将设备/边/源的
上界置零或把开关固定为开断来改变可行域，则任意单故障、多故障、孤岛和保护隔离状态至少存在一个可行点。
\end{theorynote}
证明取所有可调发电、储能、换流器传输和支路潮流为零；所有母线角度/电压取参考允许值；每个负荷取
$s_i=d_i$，固定正注入取 $p_i^{gc}=\bar p_i^{fix}$。主网节点平衡逐点成立，断开的支路流为零。
配网若无源，则令全部 $y_i=z_{ij}=0$；若存在被强制带电的根，则取上述根森林中的母线和边带电，
被删除的故障边和不采用的联络边取零，并沿树从根向每个非根发送一单位商品。也可只保留满足不可操作
常闭边约束所需的最小根分量，其余母线失电。所有负荷仍取
$p_i^{sh}=P_i^d,q_i^{sh}=Q_i^d$，故有功、无功潮流与源出力均可取零；
储能不放电，能量约束成立。
固定拓扑快捷预解若强制保留健康闭环，则可能与生成森林冲突，但通用 MILP 可选择允许的辐射拓扑。
因此在声明的有效输入域内，故障不会破坏通用模型可行性；若实现报告不可行，必须检查是否错误固定了
正出力、遗漏削减变量、电压输入区间非法、Big-M 不足、重复拓扑边、不可分段健康闭环、错误固定闭环
或数值证书失效。对缺少分段设备元数据的健康闭环，模型会把为满足辐射性产生的失电计入 N-0 原始削减；
这不是故障增量，但提示输入拓扑不在“全负荷可辐射运行”的条件内。

\subsection{审计矩阵与实现对应}
\begin{center}\footnotesize
\begin{tabular}{@{}L{3.5cm}L{4.2cm}L{4.8cm}L{2.3cm}@{}}
\toprule
层级 & 必须存在的可行性松弛 & 主要实现 & 认证结果\\
\midrule
主网 HL-II & $s,d,p^{gc}$；$P_g^{min}=0$；逐岛角度参考 & \srcpath{src/optimal\_power\_flow/dc\_opf.cpp} & 证书/失败即抛错\\
混合 AC/DC & AC/DC 削减、VSC/DC--DC 零传输 & \srcpath{src/reliability/reliability\_assessment.cpp} & 两阶段 LP 证书\\
配网 Stage 1 & $p^{sh},q^{sh}$；保护后固定开断 & \srcpath{src/reliability/three\_stage\_reliability.cpp} & MILP 后验认证\\
配网 Stage 2 & 联络开关、商品流、开关预算 & 同上；运行拓扑重构 & 拓扑/潮流校核\\
配网 Stage 3 & 保持 Stage-2，故障仍开断，储能能量递推 & 同上 & 阶段 3 重解/证书\\
\bottomrule
\end{tabular}
\end{center}

当前执行契约明确区分三种情况：\textbf{物理削减}（求解成功且 $S>0$）、\textbf{保护不准入}（保护逻辑
阻止恢复动作，但可在安全拓扑上求解削减）和\textbf{构模/数值失败}（无证书或非有限）。前两种的已认证
削减按各自持续时间进入 EENS；保护不准入返回 ``\texttt{success (protection interlock)}'' 并设置恢复未准入
标志，不得伪装成 solver failed。第三种立即返回错误，不能以全切负荷伪装成保守结果。

\subsection{验证要求}
每次修改后必须执行：
\begin{enumerate}
  \item 无机组、无 slack、有源孤岛、无源孤岛、固定注入过剩、并联支路和链接变压器元数据故障；
  \item 主网 N-0/N-1/N-2 全状态有限、可行、削减在 $[0,\sum d]$；
  \item 配网每个故障的三个阶段均有状态、证书、拓扑和热限结果；Stage 2/3 联络集合一致；
  \item 任意阶段的求解失败都能在调用边界看到阶段号、状态和残差，且 EENS/SAIDI 不增加伪造项；
  \item 结果中的模型范围、线性化边界、拒绝设备和求解证书与本章逐项一致。
\end{enumerate}

固定拓扑预解的不可行不属于上述第三类失败：它只是用来快速复用上一阶段拓扑的候选解，闭环网络在
“所有健康边均固定闭合”时可能违反辐射森林约束。实现将此情况标为可重试状态，随后调用允许运行拓扑
重构的通用 MILP；只有通用 MILP 也无法取得证书时，才按构模/数值错误中止。这样既不放松最终辐射性，
也不把合法的 Stage-1 重构误判为不可行。

\section{稀有事件估计、精确状态枚举与灵敏度}\label{sec:rel-exact}

\subsection{问题定义与三种计算口径}
设可靠性元件集合为 $\mathcal C$，状态 $x_c=1$ 表示停运，不可用率为 $U_c$。给定状态向量
$x$，后果求解器返回切负荷 $S(x)$。非序贯蒙特卡洛、精确状态枚举与一阶 FMEA 都计算同一个后果函数，
差别只在概率积分方式：
\begin{align}
 \mathrm{EDNS}&=\sum_x p(x)S(x),\qquad
 \mathrm{EENS}=H\,\mathrm{EDNS},\\
 \mathrm{PLC}&=\sum_xp(x)\mathbf 1[S(x)>\varepsilon],
 \qquad \mathrm{LOLE}=H\,\mathrm{PLC}.
\end{align}
精确枚举保留全部状态，普通蒙特卡洛按 $p(x)$ 抽样，重要抽样按扭曲分布 $q(x)$ 抽样并乘似然比。
FMEA 则只取预先定义的单元件或故障模式事件，不能代替多重停运概率积分。

对独立二态元件，状态概率为
\begin{equation}
 p(x)=\prod_{c\in\mathcal C}U_c^{x_c}(1-U_c)^{1-x_c}.
 \label{eq:rel-exact-state-probability}
\end{equation}
平台只对 $0<U_c<1$ 的非退化元件展开二进制维度。$U_c=0$ 固定运行，$U_c=1$ 固定停运；二者仍保留在
传给后果求解器的完整状态向量中。这样可避免确定状态把复杂度从 $2^m$ 虚增到 $2^n$，又不会丢掉确定停运后果。

\subsection{不可用率的统一解析}
若给出运行态故障强度 $\lambda_c$（次/年）和平均修复时间 $r_c$（小时），指数修复二态模型给出
\begin{equation}
 \mu_c=H/r_c,\qquad U_c=\frac{\lambda_c}{\lambda_c+\mu_c}.
\end{equation}
若算例直接给出强迫停运率，则解析器保留该值；若给出 MTBF，必须先依据
\srcpath{MtbfConvention} 判断它是平均失效前时间还是完整周期。所有方法经同一解析器得到
\srcpath{ReliabilityParams}，否则“方法对比”会被输入语义差异污染。

\subsection{普通非序贯蒙特卡洛}
普通抽样以 $X_c\sim\operatorname{Bernoulli}(U_c)$ 生成状态。$N$ 个样本的无偏估计为
\begin{equation}
 \widehat{\mathrm{EDNS}}=\frac1N\sum_{n=1}^NS(X^{(n)}),\qquad
 \widehat{\mathrm{PLC}}=\frac1N\sum_{n=1}^N\mathbf 1[S(X^{(n)})>\varepsilon].
\end{equation}
对切负荷样本均值，平台报告变异系数
\begin{equation}
 \mathrm{CoV}=\frac{\widehat\sigma_S/\sqrt N}{\widehat\mu_S}.
\end{equation}
当均值接近零时，相对误差会不稳定，因此同时保留固定最大样本数。固定随机种子只保证伪随机序列可复现，
不保证不同编译器、不同并行分块或后果求解器版本得到逐样本相同结果；工程复现应同时冻结构建与算例版本。

\subsection{优势比扭曲的重要抽样}
平台把每个元件的停运优势比乘以正因子 $\kappa$：
\begin{equation}
 \frac{\widetilde U_c}{1-\widetilde U_c}
 =\kappa\frac{U_c}{1-U_c},\qquad
 \widetilde U_c=\frac{\kappa U_c}{1-U_c+\kappa U_c}.
 \label{eq:rel-is-twist}
\end{equation}
$\kappa>1$ 增加故障状态出现概率，$\kappa=1$ 与普通抽样逐样本同分布。对从 $q$ 抽到的状态，似然比为
\begin{equation}
 w(x)=\frac{p(x)}{q(x)}
 =\prod_c\left(\frac{U_c}{\widetilde U_c}\right)^{x_c}
 \left(\frac{1-U_c}{1-\widetilde U_c}\right)^{1-x_c}.
\end{equation}
于是
\begin{equation}
 \widehat{\mathrm{EDNS}}_{IS}=\frac1N\sum_{n=1}^Nw_nS_n,qquad
 \widehat{\mathrm{PLC}}_{IS}=\frac1N\sum_{n=1}^Nw_n\mathbf 1[S_n>\varepsilon]
\end{equation}
保持无偏。实现使用对数似然累加再指数化，避免许多元件概率连乘的下溢；非有限权重直接报错，不以零替代。

\subsection{有效样本量与扭曲诊断}
权重退化由
\begin{equation}
 N_{eff}=\frac{(\sum_nw_n)^2}{\sum_nw_n^2}
\end{equation}
度量，满足 $0<N_{eff}\le N$。均值似然比
$\overline w=N^{-1}\sum_nw_n$ 理论上趋近 $1$，但有限样本会偏离；它是覆盖诊断，不是强制归一化因子。
平台返回 \fld{importance\_sampling\_used}、\fld{importance\_twisting\_factor}、
\fld{importance\_mean\_likelihood\_ratio} 与 \fld{importance\_effective\_sample\_size}。

不能仅凭“故障样本更多”断言方差下降。若 $\kappa$ 过大，绝大多数样本集中在极高阶停运，而真实风险由一、二阶
状态主导，权重会强烈退化。工程上应扫描多个 $\kappa$，同时比较估计标准误和 $N_{eff}/N$；不能事后挑选最有利
的单次结果。序贯蒙特卡洛目前明确拒绝该选项，因为其路径似然需包含每次跃迁与持留时间密度，不能复用静态状态权重。

\subsection{精确状态枚举}
设非退化元件数为 $m$，平台枚举 $2^m$ 个状态。对每个状态调用与非序贯 MC 相同的最小切负荷后果引擎，累加
式~\eqref{eq:rel-exact-state-probability}。时间复杂度为
\begin{equation}
 T_{exact}=O(2^mT_{consequence}),\qquad M_{exact}=O(n+m),
\end{equation}
其中实现逐状态流式累加，不保存全部后果，因此内存不随 $2^m$ 增长。默认最多二十个非退化元件；超限明确拒绝，
不静默截断元件，也不改用近似方法。

精确结果返回状态概率和、枚举状态数、基准 EENS/EDNS/PLC/LOLE，以及每个元件的稳定 ID、容器位置和内部状态位。
三种身份不得混用：\fld{global\_state\_index} 是完整状态向量位，\fld{component\_position} 是对应类型容器位置，
\fld{component\_index} 是对外稳定 \fld{.index}。HTTP 响应同时给出 \fld{stable\_id}，便于与 GUI 资产一致定位。

\subsection{Birnbaum 重要度的离散推导}
固定除元件 $c$ 外的状态 $x_{-c}$，分别计算 $c$ 运行和停运的系统后果。EENS 对 $U_c$ 的精确导数为
\begin{align}
 B_c^{EENS}
 &=\frac{\partial\mathrm{EENS}}{\partial U_c}\\
 &=H\sum_{x_{-c}}p(x_{-c})
 \left[S(x_c{=}1,x_{-c})-S(x_c{=}0,x_{-c})\right].
 \label{eq:rel-birnbaum-eens}
\end{align}
该式不是有限差分近似，因为独立二态模型对单个 $U_c$ 是仿射的。相同方法给出 PLC 的 Birnbaum 导数：
\begin{equation}
 B_c^{PLC}=\sum_{x_{-c}}p(x_{-c})
 \left[\mathbf 1(S_1>\varepsilon)-\mathbf 1(S_0>\varepsilon)\right].
\end{equation}
若元件停运在某些状态反而降低切负荷，导数可为负；实现不钳零，因为负值可能揭示调度、拓扑或后果模型的非单调性。

\subsection{EENS 导数与参数扰动}
\fld{eens\_derivative\_mwh\_yr\_per\_unit\_unavailability} 与
Birnbaum EENS 值相同，单位是“每单位不可用率变化的 MWh/年”。小扰动满足
\begin{equation}
 \Delta\mathrm{EENS}\approx B_c^{EENS}\Delta U_c.
\end{equation}
若输入参数是故障率或修复时间，还需链式法则。由 $U=\lambda/(\lambda+H/r)$ 得
\begin{align}
 \frac{\partial U}{\partial\lambda}
 &=\frac{H/r}{(\lambda+H/r)^2},\\
 \frac{\partial U}{\partial r}
 &=\frac{\lambda H/r^2}{(\lambda+H/r)^2}.
\end{align}
因此设备改造的边际价值应结合实际可改变的参数，而不能把“每单位不可用率”的导数直接解释为“每减少一次故障”。

\subsection{Fussell--Vesely 归因}
平台按元件停运状态中包含的风险定义
\begin{equation}
 FV_c=\frac{\sum_{x:x_c=1}p(x)H S(x)}{\mathrm{EENS}}.
 \label{eq:rel-fv}
\end{equation}
当总 EENS 为零时取零并声明边界。多元件同时停运的后果会同时进入多个 $FV_c$，所以各元件 FV 之和可以超过一；
它是“含该元件失效的风险占比”，不是互斥责任分摊。Birnbaum 衡量边际变化，FV 衡量现状归因，两者排序不同是正常的。

\subsection{与共现重要度的区别}
普通 MC 的 \fld{critical\_components} 以采样中“元件停运且系统失负荷”的共现和损失权重排序。该量受抽样噪声、
联合停运与扭曲权重影响，不等于式~\eqref{eq:rel-birnbaum-eens}。精确灵敏度接口单独返回 Birnbaum、EENS 导数和 FV，
从数据结构上阻止 GUI 把共现分数误标为边际重要度。

\subsection{五兆瓦直流负荷解析算例}
考虑一个不可用率 $U=0.1$ 的 VSC，其停运时使 $5$ MW 直流负荷全年不可供。两状态后果为
$S_0=0$、$S_1=5$ MW，故
\begin{align}
 \mathrm{EENS}&=0.9\times0+0.1\times5\times8760=4380\ \mathrm{MWh}/\text{年},\\
 B^{EENS}&=8760(5-0)=43800\ \mathrm{MWh}/(\text{年}\cdot\text{单位不可用率}),\\
 FV&=4380/4380=1.
\end{align}
测试把 VSC 的稳定 ID 特意设为 $77$，而容器位置仍为 $0$，从而防止实现因偶然数值相等而混淆身份空间。

\subsection{数值验收关系}
精确枚举用于校验蒙特卡洛而不是取代大系统抽样。小系统上必须满足：
\begin{enumerate}
 \item 状态概率和在 $10^{-12}$ 容差内为一；
 \item $\kappa=1$ 时重要抽样逐样本权重为一、$N_{eff}=N$；
 \item 普通 MC 与重要抽样置信区间覆盖精确 EENS；
 \item Birnbaum 与对 $U_c$ 的中心有限差分一致；
 \item 稳定 ID、容器位置和状态位可分别回溯；
 \item $U=0/1$ 元件不扩大枚举维数，但后果状态正确保留。
\end{enumerate}
实现落点为 \srcpath{compute\_exact\_reliability\_sensitivity} 与
\srcpath{run\_nonsequential\_mc}；注册验证位于 \srcpath{test\_reliability\_resolver}，HTTP 契约由
统一 \srcpath{/api/session/run\_reliability} 入口的 \fld{exact\_sensitivity} 和 NSQ 分支覆盖。

\section{精确法、抽样法与故障分析法的统一对照算例}\label{sec:rel-estimation-casebook}

保护和信息模型解决“一个状态产生什么后果”，统计估计层解决“状态以什么权重进入年度指标”。本章在同一二状态假设下对照
精确枚举、普通非序贯蒙特卡洛、重要抽样、频率--持续时间和 N-2 故障分析，说明哪些结果应一致、哪些指标因模型口径不同
不能直接比较。

\subsection{二状态元件的统一参数化}
元件 $i$ 的失效率为 $\lambda_i$ 次/年，平均修复时间为 $r_i$ h。采用常失效率、常修复率两状态 CTMC 时，修复率
$\mu_i=8760/r_i$ 次/年，稳态不可用率
\begin{equation}\label{eq:rel-case-unavailability}
 U_i=\frac{\lambda_i}{\lambda_i+\mu_i}.
\end{equation}
若输入直接给 FOR，则精确枚举和非序贯 MC 可使用 FOR 作为 $U_i$；但没有 MTTR 时不能反推出状态转移频率，F\&D 仍可
给 LOLP/LOLE，却不能伪造 LOLF 和 LOLD。

所有方法必须使用同一可靠性数据策略。\verb|case_data_only| 只消费案例字段；\verb|missing_only| 仅在字段缺失时使用默认库；
\verb|template_override| 按模板覆盖。若一行使用案例 FOR、另一行使用模板失效率，即使网络相同也不是方法对照。

\subsection{精确独立状态枚举}
对 $n$ 个非退化独立元件，状态 $x\in\{0,1\}^n$，约定 $x_i=1$ 表示停运。状态概率为
\begin{equation}
 p(x)=\prod_iU_i^{x_i}(1-U_i)^{1-x_i}.
\end{equation}
每个状态调用同一最小切负荷后果函数，得到增量损失 $L(x)$。精确年度指标为
\begin{align}
 \mathrm{EENS}_{exact}&=H\sum_xp(x)L(x),\\
 \mathrm{LOLE}_{exact}&=H\sum_xp(x)\mathbf1[L(x)>\varepsilon].
\end{align}
退化元件 $U_i=0$ 或 1 不扩张位空间，而是固定运行或停运。实现上限为 20 个非退化元件；超过上限明确拒绝，因为
$2^{20}=1,048,576$ 已是单次交互分析的可见成本，不能静默退回无误差声明的近似。

状态概率和
\begin{equation}
 r_p=\left|1-\sum_xp(x)\right|
\end{equation}
必须在浮点容差内为零。后果评估失败的状态不能以 $L=0$ 继续；否则最困难的停运状态会被系统性解释为无损失。

\subsection{Birnbaum 与 EENS 导数}
对元件 $i$，将其强制停运和强制运行，其余元件仍按原分布，定义
\begin{align}
 E_i^-&=H\,\mathbb E[L\mid x_i=1],\\
 E_i^+&=H\,\mathbb E[L\mid x_i=0].
\end{align}
Birnbaum 能量重要度为
\begin{equation}\label{eq:rel-case-birnbaum}
 B_i=E_i^--E_i^+.
\end{equation}
因为总期望可写成 $E=U_iE_i^-+(1-U_i)E_i^+$，故
\begin{equation}
 \frac{\partial E}{\partial U_i}=B_i.
\end{equation}
专项测试同时用中心有限差分修改 $U_i$，要求导数与式~\eqref{eq:rel-case-birnbaum} 一致。若把内部向量位置当成稳定组件 ID，
有限差分会修改错误元件，这项测试会立即失败。

Fussell--Vesely 风险占比定义为
\begin{equation}
 FV_i=\frac{E-E_i^+}{E},\qquad E>0.
\end{equation}
它回答“消除元件 $i$ 的随机失效可消除多少现有风险”，而 Birnbaum 回答“不可用率微增的边际代价”。二者排序可以不同，
结果表必须同时给单位：$B_i$ 为 MWh/年/单位不可用率，$FV_i$ 无量纲。

\subsection{普通非序贯蒙特卡洛}
普通 MC 从目标分布 $p(x)$ 独立抽取 $N$ 个状态，估计量
\begin{equation}
 \widehat E_{MC}=\frac{H}{N}\sum_{k=1}^NL(x^{(k)}).
\end{equation}
其方差
\begin{equation}
 \operatorname{Var}(\widehat E_{MC})=\frac{H^2}{N}
 \operatorname{Var}_p[L(X)].
\end{equation}
低概率高损失状态会使大部分样本为零，有限样本中关键状态的命中次数可能仍为零。程序返回样本数、EENS 历史和最终
变异系数；固定随机种子只保证复现，不把抽样结果变成精确值。

非序贯 MC 的 LOLE 是 $H$ 乘失负荷概率；它没有事件时间顺序，不能给真实 LOLF。响应对 LOLF 返回不可用原因，而不是
数值零。序贯 MC 才通过状态转移统计年度停电次数和持续时间。

\subsection{优势比扭曲重要抽样}
重要抽样把元件停运优势比乘扭曲因子 $\kappa>0$。目标不可用率 $U_i$ 转为提议分布
\begin{equation}\label{eq:rel-case-twisted-u}
 Q_i=\frac{\kappa U_i}{1-U_i+\kappa U_i}.
\end{equation}
从 $q(x)$ 抽样后，似然比
\begin{equation}
 w(x)=\frac{p(x)}{q(x)}=prod_i
 \left(\frac{U_i}{Q_i}\right)^{x_i}
 \left(\frac{1-U_i}{1-Q_i}\right)^{1-x_i}.
\end{equation}
无偏估计量为
\begin{equation}
 \widehat E_{IS}=\frac{H}{N}\sum_{k=1}^Nw(x^{(k)})L(x^{(k)}).
\end{equation}
实现使用未自归一化似然比保持无偏；把分母改为 $\sum_kw_k$ 会产生自归一化估计量，有限样本下有偏，不能与当前契约混用。

当 $\kappa=1$ 时，式~\eqref{eq:rel-case-twisted-u} 给 $Q_i=U_i$，每个样本 $w=1$，重要抽样必须逐样本退化为普通抽样。
专项测试固定相同随机种子，要求平均似然比精确为 1、有效样本量精确为 $N$。这是似然比方向最强的单元检查。

\subsection{有效样本量与扭曲强度}
归一化权重 $\widetilde w_k=w_k/\sum_jw_j$，有效样本量
\begin{equation}
 N_{eff}=\frac{1}{\sum_k\widetilde w_k^2}
 =\frac{(\sum_kw_k)^2}{\sum_kw_k^2}.
\end{equation}
$N_{eff}\le N$。扭曲因子增大时关键停运状态出现更多，但权重离散也增大，ESS 可能下降。最优扭曲不是“越大越好”，应在
关键状态覆盖和权重退化之间平衡。

平均似然比 $\overline w=N^{-1}\sum_kw_k$ 的理论期望为 1。有限样本明显偏离 1 可能表示样本量不足或扭曲过强；长期系统性
偏离则提示似然比公式、退化状态处理或概率下溢错误。响应同时返回 $\kappa$、ESS 和 $\overline w$，不能只显示一个“已启用
重要抽样”标志。

\subsection{精确法与抽样法的量化比较}
对小系统先计算 $E_{exact}$，再用不同种子运行普通 MC 与 IS。定义相对误差
\begin{equation}
 e_r=\frac{|\widehat E-E_{exact}|}{\max(E_{exact},\epsilon)}.
\end{equation}
单次 IS 误差不保证小于单次 MC；正确比较应在多个独立种子上比较均方误差和置信区间。重要抽样的价值是对选定稀有事件
降低方差，不是改变期望。

建议使用三组实验：$\kappa=1$ 验证恒等性；$\kappa=2\sim4$ 检查关键停运覆盖和 ESS；过大的 $\kappa$ 用于展示权重退化。
若 $\kappa=4$ 时 ESS 接近 $N$ 且关键状态更多，扭曲有效；若 ESS 只有个位数，应降低扭曲或增加样本。

\subsection{序贯蒙特卡洛与非序贯均值}
序贯 MC 对每个元件交替抽取运行时间和修复时间，生成逐时状态 $X(t)$。年度损失
\begin{equation}
 E_y=\int_0^HL(X_y(t))dt
\end{equation}
，多年度均值估计 EENS。若负荷恒定、元件独立两状态、每年从稳态初始化且无电池记忆或事件重叠，序贯与非序贯 EENS 应在
抽样误差内一致。加入逐时负荷、共因持续过程、电池耗尽和恢复状态后，两者不再同口径。

序贯结果的 LOLE 对停电区间并集积分，LOLF 统计从无损失到有损失的进入次数。平均持续时间
$\mathrm{LOLD}=\mathrm{LOLE}/\mathrm{LOLF}$ 仅在 LOLF 正且时序有效时返回。用非序贯失负荷样本数代替 LOLF 会把样本频数
误当物理事件频率。

\subsection{精确容量停运表}
频率--持续时间法对发电机容量停运状态做卷积。机组 $i$ 容量为 $C_i$、不可用率 $U_i$，从初始状态
$(0,1)$ 递推
\begin{align}
 P_{new}(c)&\mathrel{+}=P(c)(1-U_i),\\
 P_{new}(c+C_i)&\mathrel{+}=P(c)U_i.
\end{align}
容量值不分箱，只有浮点容差内相同的停运容量合并。因此
\fld{exact\_capacity\_states=true} 表示未用人为 MW 桶近似。

给定总装机 $C_T$ 和峰荷 $D$，容量不足状态满足 $C_T-c<D$。由这些状态汇总
\begin{align}
 \mathrm{LOLP}&=\sum_{c:C_T-c<D}P(c),\\
 \mathrm{LOLE}&=H\,\mathrm{LOLP}.
\end{align}
若每台机组有修复率，还可由状态进入/离开频率计算累计频率和 LOLF；缺 MTTR 时
\fld{frequency\_valid=false}，LOLP/LOLE 仍可用。

F\&D 只比较总可用发电容量与恒定峰荷，不求网络潮流、DC 负荷、保护和恢复。它与网络 MC 的 EENS 数值不能直接比较，但
在忽略网络且后果只由容量缺额决定的特制小系统中，COPT 状态概率应与精确枚举一致。

\subsection{串并联系统频率--持续时间约化}
两个独立两状态元件串联，系统工作仅当二者都工作，稳态可用率
\begin{equation}
 A_s=A_1A_2.
\end{equation}
等效失效频率可由从工作态流出的频率求得
\begin{equation}
 f_s=A_1A_2(\lambda_1+\lambda_2).
\end{equation}
等效不可用率 $U_s=1-A_s$，若 $U_s>0$，平均停运时间为 $r_s=8760U_s/f_s$。

并联系统失效仅当二者都失效，$U_p=U_1U_2$。进入双失效态的频率为
\begin{equation}
 f_p=U_1A_2\lambda_2+A_1U_2\lambda_1.
\end{equation}
实现直接使用 CTMC 稳态流量，不用“失效率相乘”这一量纲错误公式。专项测试以状态流守恒和极限情形检查：任一串联元件
完全不可用时系统不可用；任一并联元件完全可用时系统不可用率为零。

\subsection{低需求保护功能的 PFD}
对失效率 $\lambda$、完备试验周期 $T$、试验后恢复如新的低需求保护功能，周期平均失效概率
\begin{equation}
 \mathrm{PFD}_{avg}=1-\frac{1-e^{-\lambda T}}{\lambda T}.
\end{equation}
当 $\lambda T\ll1$ 时
\begin{equation}
 \mathrm{PFD}_{avg}=\frac{\lambda T}{2}
 -\frac{(\lambda T)^2}{6}+O((\lambda T)^3).
\end{equation}
实现对小参数使用稳定数值形式，避免 $1-e^{-x}$ 相消。PFD 是按需失效概率，不是年故障频率；调用方需乘需求频率后才能
形成保护拒动事件频率。

\subsection{N-2 二阶交互修正}
设正常态切负荷 $S_0$，单停运后果 $S_i,S_j$，双停运后果 $S_{ij}$。二阶交互量
\begin{equation}\label{eq:rel-case-n2-interaction}
 \Delta S_{ij}=S_{ij}-S_i-S_j+S_0.
\end{equation}
独立小不可用率下，二阶期望近似
\begin{equation}
 \mathbb E[S]\approx S_0+sum_iU_i(S_i-S_0)
 +\sum_{i<j}U_iU_j\Delta S_{ij}.
\end{equation}
若后果可加，$S_{ij}=S_i+S_j-S_0$，交互修正为零；若双停运触发孤岛或容量瓶颈，$\Delta S_{ij}>0$；若两个停运后果
重叠且不会重复切负荷，交互可为负。

故障模式对偶行保留原始联合暴露用于排序，系统聚合使用式~\eqref{eq:rel-case-n2-interaction} 修正，避免把 N-1 已计风险再加
一次。结果报告第一阶 EENS、二阶交互 EENS、跳过对数和二阶展开是否完整。超过对偶评估上限或低于联合不可用率门槛时，
完整性为假，报告不能称为完整 N-2 总量。

\subsection{N-2 加性与非加性算例}
加性算例取 $S_0=0,S_i=2,S_j=3,S_{ij}=5$ MW，则 $\Delta S_{ij}=0$，二阶修正不改变 N-1 总量。若双停运
使 10 MW 岛失电，$S_{ij}=10$，则 $\Delta S_{ij}=5$ MW，二阶风险增加
$HU_iU_j\times5$。若两次单停运都切同一 3 MW 负荷，$S_i=S_j=S_{ij}=3$，则交互为 $-3$ MW，修正删除重复计数。

专项测试至少覆盖加性零交互和网络协同非零交互，并检查对偶稳定 ID 顺序规范化，防止 $(i,j)$ 与 $(j,i)$ 重复。

\subsection{经验 VaR 与严格期望短缺}
对年度损失样本 $E_1,\ldots,E_N$ 排序，经验 VaR 是满足累计概率达到 $\alpha$ 的最小值。离散分布在 VaR 点可能有概率原子，
严格期望短缺需只取该原子中进入最坏 $1-\alpha$ 尾部的部分。设大于 VaR 的样本权重为 $p_>$，VaR 原子权重为 $p_=$，则
取用比例
\begin{equation}
 \gamma=\frac{1-\alpha-p_>}{p_=}
\end{equation}
并计算
\begin{equation}
 \mathrm{ES}_{\alpha}=\frac{
 \mathbb E[E\mathbf1(E>\mathrm{VaR})]
 +\gamma\,\mathrm{VaR}\,P(E=\mathrm{VaR})}{1-\alpha}.
\end{equation}
简单平均所有 $E\ge\mathrm{VaR}$ 会把整个 VaR 原子纳入，离散样本下可能严重偏差。实现以分数权重处理原子，测试用大量重复
损失值与闭式尾部均值比较。

\subsection{数据质量对结果的影响}
可靠性参数可能来自案例、模板、默认补全或非法输入修正。结果记录来源计数、缺失数、默认使用数和警告。默认库能让探索性
分析运行，但不能把默认补全结果解释为资产实测风险。比较算法时，数据策略和系统指纹必须相同；比较资产时，还应报告每个
关键元件的参数来源。

MTBF 约定也会改变失效率。若 MTBF 表示运行时间均值，$\lambda=8760/MTBF$；若表示包含修复的日历周期间隔，则需要结合
MTTR 解析。平台要求显式约定并在解析器中统一转换，避免 MC、FMEA 和 F\&D 各自解释同一字段。

\subsection{方法一致性矩阵}
在独立二状态、恒定负荷、相同后果函数的条件下，精确枚举是无抽样基准，普通 MC 与 $\kappa=1$ 的 IS 应收敛到它；适度
扭曲 IS 仍收敛到同一均值但方差可降低；序贯 MC 的长期 EENS 也应一致，并额外给事件频率。F\&D 仅在忽略网络、后果只由
容量缺额决定时与精确枚举一致。N-1/N-2 FMEA 按故障频率和持续时间展开，不等同于完整稳态状态概率，除非稀有事件近似和
持续时间定义严格匹配。

因此“六种算法结果不同”本身不证明哪一个错误。先检查：状态空间是否相同、负荷是否恒定、是否有时序记忆、是否求网络
后果、是否包含 N-2、是否使用同一数据策略、是否报告增量还是含 N-0 的总量。只有这些条件一致后，误差才可归因于抽样或
近似。

\subsection{程序级验收顺序}
估计层推荐按以下不变量验收：解析后的 $U_i$ 在 $[0,1]$；精确状态概率和为一；$\kappa=1$ 时权重全为一且 ESS 为 $N$；
扭曲后 ESS 在 $(0,N]$ 且平均似然比有限；Birnbaum 与有限差分一致；退化元件不扩张状态数；F\&D 容量概率和为一；缺 MTTR
时频率有效性为假；串并联 CTMC 流量守恒；PFD 小参数极限正确；N-2 加性算例交互为零；经验 ES 正确处理 VaR 原子。

HTTP E2E 进一步检查稳定身份、状态数、重要抽样诊断字段和 GUI 表格。精确灵敏度行必须同时返回
\fld{stable\_id}、组件稳定索引、内部位置、Birnbaum 和 FV；前端以稳定 ID 支持画布定位。测试不接受只有数组长度正确但
身份字段缺失的结果。

\subsection{与保护信息联合方法的衔接}
精确信息状态枚举与精确物理停运枚举解决不同维度。保护信息事件树通常在每个物理故障场景内部精确枚举信息元件状态，再按
场景年频率聚合；它不自动枚举全系统所有物理元件的同时停运。若需联合物理全状态与信息全状态，状态数会成为
$2^{n_p+n_c}$，应先利用条件独立、割集或重要抽样控制成本，并保持似然比覆盖所有被扭曲维度。

当前三级对照有意固定一个物理起始场景，专门测量保护和信息因素的影响；精确灵敏度入口有意固定独立物理二状态模型，专门
测量元件不可用率导数。把二者并列报告比把它们强行合成一个没有可解释基准的巨大模型更有审计价值。

\subsection{可复现证据}
对应实现位于 \srcpath{src/reliability/reliability\_assessment.cpp} 的非序贯、序贯、重要抽样、精确灵敏度、F\&D、PFD 与
串并联约化函数，以及 \srcpath{src/reliability/failure\_mode.cpp} 的二阶交互聚合。专项 C++ 测试覆盖闭式预言，
\srcpath{tests/e2e/reliability\_workflow\_e2e.mjs} 覆盖 HTTP 和 GUI。最终报告应给出命令、通过的用例/断言、依赖提交和
模型边界；若本地兄弟依赖提交与仓库 pin 不同，只能声称当前工作区验证通过，不能声称锁定环境可复现。

\subsection{概率与频率混用的诊断算例}
不可用率 $U$ 是任意时刻处于停运态的概率，年故障率 $\lambda$ 是单位时间内进入停运态的次数，二者量纲不同。取
$\lambda=2$ 次/年、$r=4$ h，则稀有事件近似 $U\approx\lambda r/8760=9.1324\times10^{-4}$。若 5 MW 负荷在每次故障
期间持续停电，频率--持续时间计算
\begin{equation}
 E=\lambda\times5\times4=40\ \text{MWh/年}.
\end{equation}
稳态概率计算同样给 $E=8760\times U\times5\approx40$ MWh/年。若误把 $\lambda=2$ 当成概率，由于概率超过一会被输入
校验拒绝；若先错误截成一，则得到 43800 MWh/年。若误把 $U$ 当年频率，则得到约 0.0183 MWh/年。数量级检查应成为结果
审计的第一步。

精确式~\eqref{eq:rel-case-unavailability} 在故障率不再很小时比 $\lambda r/8760$ 更稳健。解析器返回最终 $U$、原始字段和
约定，使审计者可以复算。测试不能只断言 $0\le U\le1$，还应检查给定 $\lambda,r$ 的闭式值。

\subsection{正常态基线污染的诊断算例}
若正常系统已经有 $S_0=1$ MW 切负荷，故障态切负荷 $S_1=3$ MW，非序贯 MC 若直接报告总损失，会把正常态问题算进故障
风险。平台同时报告
\begin{align}
 E_{raw}&=H\,\mathbb E[S(X)],\\
 E_{base}&=HS_0,\\
 E_{inc}&=H\,\mathbb E[\max(0,S(X)-S_0)].
\end{align}
假设故障概率 0.01，则 $E_{raw}=8760(0.99\times1+0.01\times3)=8935.2$ MWh/年，基线为 8760 MWh/年，故障增量只有
175.2 MWh/年。只展示 8935.2 会把拓扑或正常调度缺陷误归因于随机故障。

GUI 发现非零基线时显示正常态切负荷、原始 EENS 和增量 EENS。跨方法比较必须选择同一口径：FMEA 通常按故障增量聚合，
若拿它与 MC 原始总量比较会出现固定的 $E_{base}$ 差异。

\subsection{似然比数值下溢的诊断}
当元件很多或 $\kappa$ 很大，直接连乘似然比可能下溢。稳定计算先累加
\begin{equation}
 \log w(x)=\sum_i\left[x_i\log\frac{U_i}{Q_i}
 +(1-x_i)\log\frac{1-U_i}{1-Q_i}\right],
\end{equation}
再在可表示范围内指数化。退化概率 $U_i=0$ 或 1 不参与扭曲与对数比，避免 $0/0$。若提议分布给目标分布正概率状态零概率，
则不满足绝对连续性，入口必须拒绝，而不是产生无限权重。

诊断同时观察平均似然比和 ESS：全部权重下溢为零会使两者无定义；单个权重支配会使 ESS 接近 1。程序要求扭曲因子有限正，
并在结果中返回有限诊断值。工程上还应比较多种扭曲强度，避免只挑一个偶然接近精确值的种子。

\subsection{退化元件与状态数的诊断}
考虑 5 个元件，其中两个 $U=0$、一个 $U=1$、两个满足 $0<U<1$。完整物理组合虽然形式上有 $2^5$ 个，正概率状态只有
$2^2=4$ 个。精确枚举应报告 4 个状态，两个永不故障元件固定运行，完全不可用元件固定停运。若仍枚举 32 个状态，会出现
大量零概率后果求解，浪费计算且可能让零概率失败状态中止整个任务。

Birnbaum 对退化元件仍可通过“强制运行/停运”的反事实定义计算，但当前精确入口的随机状态空间只扩张非退化元件。报告应
区分 \fld{components\_evaluated} 与 \fld{stochastic\_components}，不能用状态数反推系统总组件数。

\subsection{秒、小时与年度的单位诊断}
保护清除窗口使用秒，恢复与修复窗口使用小时，年度频率使用次/年。一个 20 MW 故障、清除时间 0.18 s、年频率 3 次的
清除窗口贡献为
\begin{equation}
 E_c=3\times20\times0.18/3600=0.003\ \text{MWh/年}.
\end{equation}
若遗漏除以 3600，会得到 10.8 MWh/年，相差三个数量级。相反，修复时间 4 h 若又除以 3600，会把 240 MWh/年错误降为
0.0667 MWh/年。

因此三窗口实现先把每段时长规范化为小时再乘 MW；响应字段名以 \verb|_s|、\verb|_hr|、\verb|_mwh_yr| 标明单位。测试用
秒级闭式算例专门捕获这一错误，不依赖大系统结果的偶然抵消。

\subsection{比较报告的最小证据集}
一份可审查的方法对比至少应包含：系统指纹和模型名称；可靠性数据策略及默认补全数；负荷倍率和年度小时数；方法的状态空间
或事件空间；物理后果模型 scope；EENS/LOLE/LOLF 的绝对值和单位；N-0 基线；抽样数、随机种子、变异系数或 ESS；精确法
的状态数和概率残差；F\&D 的概率/频率有效性；N-2 的交互项与完整性；保护信息三级对照的两项差值；在线模式的轨迹消费
和清除时刻一致性；所有限制与失败状态。

只有百分比而没有绝对值不可接受。若基线接近零，百分比可能任意大；若两方法的结果口径不同，百分比没有可解释分母。只有
排名相关系数也不够，因为相同排名可以对应数量级完全不同的风险。GUI 的比较功能因此同时保留绝对指标、归一化离散度、
Spearman、top-k Jaccard 和共同薄弱元件。

\subsection{从偏差到定位的顺序}
当实测结果偏离理论预言时，按以下顺序定位：先核对单位和数据策略；再核对 N-0 基线与增量口径；检查状态概率或事件类概率
守恒；检查稳定 ID 与内部位置映射；检查后果求解是否对失败状态诚实返回；抽样法再检查随机种子、权重、ESS 和置信区间；
在线保护再检查量测轨迹、动作事件、失电岛和 DER 终态；最后才调整数值容差。

如果偏差远大于理论允许的 $O(N^{-1/2})$ 抽样误差或一步时间离散误差，不能通过扩大断言范围解决。应重新推导被测估计量，
构造最小闭式算例，并确认程序实际消费了发生偏差的参数。该流程保证性能或功能增强仍服从同一理论契约。


% theory_reliability_failure_analysis.tex —— 元件失效模型、故障模式 FMEA、共故障与信息物理评估
% 本文件为理论层章节，遵循 docs/modules/THEORY_WRITING_GUIDE.md。
% 数学模型依据 docs/theory/reliability_mathematical_models_and_intelligent_cyber_physical_assessment.md。

\section{元件失效模型与故障后果分析}\label{sec:rel-fa}

\begin{theorynote}
本章展开\emph{元件级}失效的详细数学模型与\emph{故障分析}（FMEA）方法，是 \S\ref{sec:rel-th}
系统级评估的物理与统计基座，依据设计文档
\srcpath{docs/theory/reliability\_mathematical\_models\_and\_intelligent\_cyber\_physical\_assessment.md}
（第 3、8、9、12 节）。覆盖：可修元件的交替更新过程与“运行年—日历年”故障强度之辨、按需失效模式
（$\lambda^{act}=\nu p^d$、按需失效概率 PFD）、确定性两阶段（隔离—修复）元件 FMEA、故障模式 FMEA
的原因—后果目录与二阶共故障（N-2）近似、以及\emph{已实现}的 Level-1 信息物理 FMEA（自动化可用/不可用
双类全期望细化、反事实归因与三维因子化筛选）。全部“已实现”结论以 \srcpath{src/reliability/} 为落点；
凡实现采用稀有事件近似或独立性假设处，均就地标注其有效边界，绝不以近似冒充精确。
\end{theorynote}

\subsection{可修元件：交替更新过程}\label{sec:rel-fa-renewal}
两状态可修元件的最一般刻画是\emph{交替更新过程}：运行时长 $T^\uparrow$ 与故障（修复）时长
$T^\downarrow$ 独立同分布、均值有限，故障后修复至“修旧如新”。设 $H$ 为年小时数。

\begin{theorem}[更新—回报不可用度]\label{thm:rel-fa-renewal}
在上述假设下，长期不可用度为
\begin{equation}
 U=\frac{\mathbb E[T^\downarrow]}{\mathbb E[T^\uparrow]+\mathbb E[T^\downarrow]}
 =\frac{\mathrm{MTTR}}{\mathrm{MTTF}+\mathrm{MTTR}},
 \label{eq:rel-fa-U}
\end{equation}
\emph{不}要求指数分布的运行 / 修复时长。仅当把模型解释为两状态连续时间马尔可夫链
（指数时长）时，才有转移强度 $\alpha=\lambda/H$、$\beta=1/r$ 与
$U=\alpha/(\alpha+\beta)=\lambda r/(H+\lambda r)$。
\end{theorem}

\begin{proof}
由更新—回报定理，长期处于“下”状态的时间比例等于每个更新周期内下状态期望时长占周期期望总时长之比，
即 $\mathbb E[T^\downarrow]/(\mathbb E[T^\uparrow]+\mathbb E[T^\downarrow])$；此结论只依赖均值存在与
遍历性，与时长分布形状无关。指数情形下 CTMC 平稳方程 $\pi_\uparrow\alpha=\pi_\downarrow\beta$、
$\pi_\uparrow+\pi_\downarrow=1$ 给出 $\pi_\downarrow=\alpha/(\alpha+\beta)$，与~\eqref{eq:rel-fa-U} 一致。
\end{proof}

\paragraph{“运行年”与“日历年”故障强度之辨。}
式~\eqref{eq:rel-fa-U} 中 $\lambda$ 必须是\emph{运行年}强度（元件在运行时的故障风险率）。若输入为按
\emph{日历年}统计的经验故障次数 $f_{cal}$，则
\begin{equation}
 U_{cal}=\frac{f_{cal}r}{H},\qquad f_{cal}=(1-U)\lambda_{up},\qquad
 \lambda_{up}=\frac{f_{cal}}{1-f_{cal}r/H}.
 \label{eq:rel-fa-cal}
\end{equation}
二者仅在小 $U$ 时近似相等：$U\approx\lambda r/H$（$\lambda r\ll H$）。平台解析器
（\srcpath{resolve\_reliability\_params}）把故障率输入\emph{按运行年强度} $\lambda$ 解释；采用观测日历
频率的研究须按~\eqref{eq:rel-fa-cal} 换算，或显式声明稀有事件近似。多工况、部分降额、非指数老化、
天气调制、相关失效与共享修复资源均超出两状态元组，需多状态 / 半马尔可夫 / 时序模型。

\subsection{按需失效模式}\label{sec:rel-fa-demand}
保护、开关等\emph{按需}功能以失败频率约简
\begin{equation}
 \lambda_m^{act}=\nu_m\,p_m^d\quad\Big[\tfrac{\text{需求}}{\text{年}}\Big]\Big[\tfrac{\text{失败操作}}{\text{需求}}\Big]
 =\Big[\tfrac{\text{失败操作}}{\text{年}}\Big],
 \label{eq:rel-fa-act}
\end{equation}
它是\emph{期望失败操作频率}，既非年概率、亦非 MTTF、亦非“需求时不可用”的概率。若需求为泊松过程、
每次独立以 $p_m^d$ 失败，则由泊松稀疏化，年内至少一次失败操作的概率为
\begin{equation}
 \Pr(N_m^{fail}\ge1)=1-e^{-\nu_m p_m^d}\ (\ne\nu_m p_m^d\ \text{除稀有近似}),\qquad
 \Pr(N_m^{fail}\ge1)=1-(1-p_m^d)^n\ (n\ \text{次独立需求}).
\end{equation}
对异构触发事件，更稳健的形式为
$\lambda_m^{act}=\sum_{k\in\mathcal K_m}\lambda_k\Pr(m\text{ 失败}\mid k)$，含信息物理类时再对类 $c$
以 $\pi_c(k)$ 加权。若共享驱动 $W$（如极端天气）同时改变需求强度与成功率，正确的边缘频率是
$\lambda_m^{act}=\sum_k\int\lambda_k(w)\Pr(m\text{ 失败}\mid k,w)f_W(w)\,\mathrm dw$，用
$\mathbb E[\lambda_k(W)]\,\mathbb E[p_m^d(W)]$ 替代将强加独立性、恰好漏掉关注的严重工况。

\paragraph{按需失效概率与时间不可用度之别。}
对休眠保护 / 开关，关键量常是按需失效概率 $\mathrm{PFD}_{avg}$ 而非时间不可用度。若危险未检出失效率
为 $\lambda_{DU}$、完备验证试验周期为 $T_I$，低率近似为 $\mathrm{PFD}_{avg}\approx\lambda_{DU}T_I/2$，
失败频率 $\lambda_m^{act}=\nu_m\mathrm{PFD}_{avg}$。平台在给出恢复时长 $r_m$ 时另算恢复占时
$U_m^{recovery}=\lambda_m^{act}r_m/(H+\lambda_m^{act}r_m)$——它\emph{只是}失败后恢复态的时间占比，
\emph{不是} $p_m^d$ 或 $\mathrm{PFD}_{avg}$；同理 $\mathrm{mttf\_hr}=H/\lambda_m^{act}$ 只是失败事件间的
平均间隔，不是器件寿命。故洁净数据契约须区分三量：$p_m^d$/$\mathrm{PFD}_{avg}$、$\lambda_m^{act}=\nu_m p_m^d$、
$U_m^{recovery}$。依据：\srcpath{resolve\_reliability\_params}（主动分支）。

\subsection{确定性两阶段元件 FMEA}\label{sec:rel-fa-fmea}
每个在运元件 $k$ 作 N-1 事件评估，事件含\emph{隔离/开关}阶段（时长 $\tau_k^{sw}$）与其后的\emph{修复}
阶段，故障在两阶段全程失电：
\begin{equation}
 \tau_k^{rep}=\max(0,\,r_k-\tau_k^{sw}).
\end{equation}
修复阶段引擎可用抢修/开关重构、DER 再调度、储能、构网 VSC 与微网孤岛；储能受功率—能量双限
\begin{equation}
 P_b^{avail}=\min\!\Big(P_b^{rated},\ \frac{(E_b-E_b^{min})\eta_b^{dis}}{\tau_s}\Big).
\end{equation}
以两阶段切负荷 $S_k^{sw},S_k^{rep}$ 与日历频率 $f_k^{cal}$，暴露一致的解析聚合为
\begin{align}
 \mathrm{EENS}&=\sum_k f_k^{cal}\big(S_k^{sw}\tau_k^{sw}+S_k^{rep}\tau_k^{rep}\big),\\
 \mathrm{LOLE}&=\sum_k f_k^{cal}\big(\tau_k^{sw}\mathbf 1[S_k^{sw}{>}\varepsilon]+\tau_k^{rep}\mathbf 1[S_k^{rep}{>}\varepsilon]\big),\\
 \mathrm{LOLF}&=\sum_k f_k^{cal}\,\mathbf 1[S_k^{sw}{>}\varepsilon\ \lor\ S_k^{rep}{>}\varepsilon],
\end{align}
其中 $f_k^{cal}=(1-U_k)\lambda_k^{up}$。\textbf{实现注}：平台直接以解析器的 $\lambda_k^{up}$ 代 $f_k^{cal}$，
即稀有事件近似 $f_k^{cal}\approx\lambda_k^{up}$，字面解释时把日历事件数高估 $1/(1-U_k)$ 倍，仅小 $U_k$ 可忽略；
长修复 / 高不可用元件须显式换算。忽略的重叠项为二阶
$\Pr(i,j\text{ 同时停})=O\big((f_i^{cal}r_i/H)(f_j^{cal}r_j/H)\big)$。LOLF 每个触发列一次损失、不计事件内
多段损失（如储能耗尽后再失）。依据：\srcpath{src/reliability/reliability\_assessment.cpp:run\_distribution\_fmea}。

\subsection{故障模式 FMEA 与二阶共故障}\label{sec:rel-fa-mode}
富组件展开为\emph{被动}与\emph{主动}失效模式，原因类含：物理、信息控制、通信、量测、保护逻辑、人因、
计划；后果含：强迫停运、降额、卡滞、拒动/误动（开/合/跳）、定值冻结、通信失联、失去构网、区域跳闸。
单模式确定性结果为
\begin{equation}
 f_m=\begin{cases}\lambda_m,&\text{被动（稀有暴露近似）},\\ \nu_m p_m^d,&\text{主动按需},\end{cases}\quad
 d_m=\tau_m^{iso}+\tau_m^{sw}+\max(r_m,r_m^{mode}),\quad
 \mathrm{EENS}_m=f_m d_m S_m.
 \label{eq:rel-fa-mode}
\end{equation}
暴露一致的被动模式应取 $f_m=f_m^{cal}=(1-U_m)\lambda_m^{up}$；主动分支已是事件/年单位，不再乘 $(1-U)$。

\begin{proposition}[二阶共故障近似]\label{prop:rel-fa-pair}
对相异模式 $i,j$，平台以
\begin{equation}
 U_i\approx\min\!\Big(1,\tfrac{f_i d_i}{H}\Big),\quad U_{ij}=U_iU_j,\quad
 d_{ij}=\frac{d_id_j}{d_i+d_j},\quad f_{ij}=\frac{f_if_j(d_i+d_j)}{H},\quad
 \mathrm{EENS}_{ij}=U_{ij}HS_{ij},
 \label{eq:rel-fa-pair}
\end{equation}
其中等效重叠时长 $d_{ij}$ 恰使 $f_{ij}d_{ij}/H=U_iU_j$。这是 N-2 暴露/排序近似，非精确不相交状态展开。
\end{proposition}

精确二阶交互展开（\emph{未实现}）应为
\begin{equation}
 \mathbb E[S]\approx S_0+\sum_iU_i(S_i-S_0)+\sum_{i<j}U_iU_j\big(S_{ij}-S_i-S_j+S_0\big),
 \label{eq:rel-fa-exact}
\end{equation}
需一致截断的状态概率。故当前实现把成对贡献\emph{加入}一阶总量，科学解读时应\emph{单列}成对项；
$O(U_iU_jU_k)$ 及以上、修复资源与时序重叠交互均省略。依据：
\srcpath{src/reliability/failure\_mode.cpp:run\_failure\_mode\_fmea}。

\subsection{Level-1 信息物理 FMEA（已实现）}\label{sec:rel-fa-cyber}
对每个 N-1 触发 $k$，以两个\emph{后果等价}的信息类细化物理 FMEA：自动化可用（$c=a$，条件概率
$A_k=\Pr(C=a\mid K=k)$）与不可用（$c=m$，$1-A_k$，人工响应，$\tau_k^m\ge\tau_k^a$，且可冻结 DER/储能/
构网/黑启动/微网等远程资源）。类条件后果 $E_{k,c}=S_{k,c}^{sw}\tau_{k,c}^{sw}+S_{k,c}^{rep}\tau_{k,c}^{rep}$，
则\emph{全期望}细化为
\begin{equation}
 \boxed{\ \mathrm{EENS}=\sum_k\lambda_k\big[A_kE_{k,a}+(1-A_k)E_{k,m}\big]
 =\sum_k\lambda_k\sum_{c\in\{a,m\}}\Pr(C{=}c\mid K{=}k)\,E_{k,c}.\ }
 \label{eq:rel-fa-cyber}
\end{equation}
它是物理 FMEA 的全期望细化，\emph{非}新指标。$A_k$ 一般 $\ne\Pr(C{=}a)$，因天气/站用电/站点损毁可同时
驱动两层——全局标量或逐元件覆盖是\emph{筛选}输入，隐含独立或用户已给条件值。

\paragraph{反事实归因与自动化效率。}
令 $E_k^{dur}$ 用人工时长但自动控制仍在。平台报告
\begin{align}
 \mathrm{EENS}^{perfect}&=\sum_k\lambda_kE_{k,a},\\
 \Delta^{cyb,dur}&=\sum_k\lambda_k(1-A_k)(E_k^{dur}-E_{k,a}),\\
 \Delta^{cyb,ctl}&=\sum_k\lambda_k(1-A_k)(E_{k,m}-E_k^{dur}),
\end{align}
且 $\mathrm{EENS}=\mathrm{EENS}^{perfect}+\Delta^{cyb,dur}+\Delta^{cyb,ctl}$，并给出无自动化上界
$\mathrm{EENS}^{noauto}=\sum_k\lambda_kE_{k,m}$、自动化效率
$\eta_A=(\mathrm{EENS}^{noauto}-\mathrm{EENS})/(\mathrm{EENS}^{noauto}-\mathrm{EENS}^{perfect})$ 与
信息物理致 SAIDI 份额 $\gamma_{SAIDI}=(\mathrm{SAIDI}-\mathrm{SAIDI}^{perfect})/\mathrm{SAIDI}$。
\begin{gapnote}
时长/控制增量是\emph{反事实归因}，非天然非负测度：其非负性依赖单调性假设
$E_{k,a}\le E_k^{dur}\le E_{k,m}$。若较慢开关反而启用更优拓扑、或冻结控制意外改善调度，任一增量可为负；
此为“完美—退化”序非单调的诊断证据，\emph{不得}静默钳零。Level-1 无显式通信拓扑、时延、信息节点电源、
共因信息失效或学习型 FDIR 策略；保护误动由故障模式 FMEA 单独处理，其 Level-1 分解字段为占位零。
\end{gapnote}

\paragraph{三维因子化筛选。}
当信息与智能维启用时，自动类有效概率因子化为
\begin{equation}
 A_k^{eff}=A_k^{info}\,P_D\,P_I\,P_R^{valid}\,P_R^{exec}\,P_P^{success},
 \label{eq:rel-fa-factor}
\end{equation}
即信息服务可用度乘检测 $P_D$、隔离 $P_I$、恢复决策有效 $P_R^{valid}$、恢复执行成功 $P_R^{exec}$、
保护成功 $P_P^{success}$；禁用维取因子 $1$，逐元件信息覆盖先替换 $A_k^{info}$。这是\emph{独立因子筛选}
近似，非联合类生成器：失败的智能功能被路由到既有退化/人工后果代理，$P_P^{success}$ 不推导继电器起动、
后备区扩展、断路器失灵或 DER 穿越轨迹，故结果声明
\texttt{joint\_class\_probability\_modelled=false}、\texttt{protection\_logic\_modelled=false}、
\texttt{protection\_frt\_reliability\_coupled=false}。依据：
\srcpath{src/reliability/reliability\_assessment.cpp:run\_distribution\_fmea}、\srcpath{failure\_mode.hpp:CyberPhysicalFMEAOptions}。

\subsection{元件类型专属失效模式目录}\label{sec:rel-fa-catalog}
上述框架的\emph{具体实例}由目录构建器 \srcpath{src/reliability/failure\_mode.cpp:build\_failure\_mode\_catalog}
生成：每个富元件按其类型展开为一至多个失效模式，每个模式携带激活类（被动 / 按需）、起因类
（物理 / 控制 / 通信 / 量测 / 保护）、映射后果（\S\ref{sec:rel-fa-mode} 的 $\Phi_m$）、经解析器
（\S\ref{sec:rel-param}）解析的参数、三段时长 $(\tau^{iso},\tau^{sw},\tau^{rep})$，以及降额模式的残余容量因子
$\rho\in(0,1]$。表~\ref{tab:rel-fa-catalog} 中的“模板默认”是\emph{可替换的筛选先验}，并非设备标准；
优先级为\textbf{算例数据 $\succ$ 选定模板 $\succ$ 筛选默认}，缺失量在严格数据策略下保持未确定。
被动模式默认以 $(\lambda\,[\text{occ/yr}],\,r\,[\text{hr}])$ 给出，按需模式以
$(p_d,\,\nu_d\,[/\text{yr}])$ 给出。
\begin{center}\footnotesize
\begin{longtable}{@{}L{2.2cm}L{3.9cm}L{4.5cm}L{3.4cm}@{}}
\caption{元件类型专属失效模式目录（模板默认为筛选先验）}\label{tab:rel-fa-catalog}\\
\toprule
\textbf{元件类型} & \textbf{失效模式} & \textbf{激活·起因·后果} & \textbf{模板默认 / 残余}\\
\midrule
\endfirsthead
\toprule
\textbf{元件类型} & \textbf{失效模式} & \textbf{激活·起因·后果} & \textbf{模板默认 / 残余}\\
\midrule
\endhead
\midrule\multicolumn{4}{r}{\footnotesize 续下页}\\
\endfoot
\bottomrule
\endlastfoot
AC/DC 母线 & \fld{busbar\_fault} & 被动·物理·强迫停运 & $\lambda{=}0.01,\ r{=}8$\\
\midrule
AC 发电机 & \fld{forced\_outage} & 被动·物理·强迫停运 & FOR，默认 $\lambda{=}4.38,\ r{=}50$\\
 & \fld{derating} & 被动·物理·降额 & $\lambda{=}0.5,\ r{=}24$；$\rho{=}0.6$\\
 & \fld{control\_unavailable} & 被动·控制·失控 & 恢复 $2$~h\\
AC 支路 & \fld{permanent\_fault} & 被动·物理·强迫停运 & $\lambda{=}0.35,\ r{=}10$\\
 & \fld{thermal\_derating} & 被动·物理·降额 & $\lambda{=}0.1,\ r{=}8$；$\rho{=}0.75$\\
变压器 2W & \fld{internal\_fault} & 被动·物理·强迫停运 & $\lambda{=}0.03,\ r{=}200$\\
 & \fld{tap\_changer\_failure} & 按需·物理·失控 & $p_d{=}0.005,\ \nu_d{=}50$\\
 & \fld{cooling\_derating} & 被动·物理·降额 & $\lambda{=}0.05,\ r{=}24$；$\rho{=}0.7$\\
变压器 3W & \fld{internal\_fault}、\fld{cooling\_derating} & 同 2W（内部故障 $\lambda{=}0.04$）& $r{=}200$；$\rho{=}0.7$\\
\midrule
静止发电机 & \fld{unit\_outage} & 被动·物理·强迫停运 & $\lambda{=}1.5,\ r{=}24$\\
 & \fld{dispatch\_control\_unavailable} & 被动·控制·失控 & 恢复 $2$~h\\
新能源机组 & \fld{unit\_outage} & 被动·物理·强迫停运 & $\lambda{=}2.0,\ r{=}48$\\
 & \fld{resource\_derating} & 被动·物理·降额 & $\rho{=}0.5$\\
AC 光伏 & \fld{plant\_outage} & 被动·物理·强迫停运 & $\lambda{=}1.5,\ r{=}24$\\
 & \fld{inverter\_failure} & 被动·物理·降额 & $\rho{=}0.5$\\
AC 储能 & \fld{unit\_outage} & 被动·物理·强迫停运 & FOR，$r{=}24$\\
 & \fld{pcs\_failure} & 被动·物理·降额 & $\rho{=}0.5$\\
 & \fld{bms\_control\_unavailable} & 被动·控制·失控 & 恢复 $2$~h\\
\midrule
AC 开关 & \fld{hardware\_outage} & 被动·物理·强迫停运 & $\lambda{=}0.05,\ r{=}4$\\
 & \fld{fail\_to\_open}、\fld{sectionalize} & 按需·物理·拒开 & $p_d{=}p_{sw},\ \nu_d{=}2$\\
 & \fld{fail\_to\_close} & 按需·物理·拒合 & 重合器取 $1{-}p_{rc}$\\
 & \fld{fail\_to\_clear}（熔断器）& 按需·保护·拒动 & $\nu_d{=}1$\\
 & \fld{stuck\_closed}、\fld{nuisance\_trip} & 被动·物理·卡合／保护·误动 & 用户先验\\
 & \fld{comm\_loss} & 被动·通信·通信丢失 & 恢复 $2$~h\\
AC 断路器 & \fld{hardware\_outage} & 被动·物理·强迫停运 & $\lambda{=}0.05,\ r{=}8$\\
 & \fld{fail\_to\_trip} & 按需·保护·拒动 & $p_d{=}0.005,\ \nu_d{=}1$\\
 & \fld{fail\_to\_close} & 按需·物理·拒合 & $p_d{=}0.01,\ \nu_d{=}2$\\
 & \fld{nuisance\_trip}、\fld{comm\_status\_failure} & 被动·保护·误动／通信·通信丢失 & 恢复 $2$~h\\
\midrule
VSC 换流器 & \fld{power\_stage\_outage} & 被动·物理·强迫停运 & FOR，$r{=}48$\\
 & \fld{derating} & 被动·物理·降额 & $\rho{=}0.7$\\
 & \fld{grid\_forming\_lost} & 按需·控制·构网失效 & $p_d{=}0.02,\ \nu_d{=}1$\\
 & \fld{setpoint\_frozen} & 被动·控制·定值冻结 & 恢复 $1$~h\\
 & \fld{comm\_loss}、\fld{measurement\_bias} & 被动·通信／量测 & 恢复 $2$/$4$~h\\
DC 支路 & \fld{pole\_fault} & 被动·物理·强迫停运 & $\lambda{=}0.20,\ r{=}24$\\
 & \fld{derating} & 被动·物理·降额 & $\rho{=}0.75$\\
DCDC 换流器 & \fld{power\_stage\_outage} & 被动·物理·强迫停运 & $\lambda{=}0.20,\ r{=}48$\\
 & \fld{derating}、\fld{setpoint\_frozen}、\fld{comm\_loss} & 被动·物理／控制／通信 & $\rho{=}0.7$\\
DC 断路器 & \fld{hardware\_outage} & 被动·物理·强迫停运 & $\lambda{=}0.10,\ r{=}8$\\
 & \fld{fail\_to\_trip}、\fld{fail\_to\_close} & 按需·保护·拒动／物理·拒合 & $p_d{=}0.01$\\
DC 储能／光伏／源 & \fld{unit\_outage}、\fld{array\_outage}、\fld{source\_outage} & 被动·物理·强迫停运 & $\lambda{=}1.5,\ r{=}24$\\
\end{longtable}
\end{center}
聚合元件（微网、虚拟电厂、移动储能、能量路由器）先由
\srcpath{src/reliability/reliability\_assessment.cpp:materialize\_aggregated\_reliability\_sources}
物化为 PCC/端口处的可调注入或 AC$\leftrightarrow$DC 传输，其 \texttt{unit\_outage} 置其出服务即切除其所供负荷。
所有模式经\emph{支持门}（\srcpath{src/reliability/failure\_mode.cpp:build\_consequence\_patch}）过滤：稳态
引擎无法表示者（如纯量测偏差、非可调设备的通信丢失、拒合）记为 \fld{unsupported} 而\emph{不}静默施加，
诚实见 \S\ref{sec:rel-src-boundary}。

\subsection{与实现的对应}\label{sec:rel-fa-impl}
\begin{center}\footnotesize
\begin{longtable}{@{}L{5.2cm}L{9.2cm}@{}}
\toprule
\textbf{理论对象} & \textbf{平台实现状态}\\
\midrule
\endfirsthead
\toprule
\textbf{理论对象} & \textbf{平台实现状态}\\
\midrule
\endhead
\bottomrule
\endfoot
交替更新不可用度~\eqref{eq:rel-fa-U}、运行年/日历年换算 &
  \implpart{解析器按运行年 $\lambda$ 解释输入；日历换算须调用方处理}\\
按需失效 $\lambda^{act}=\nu p^d$~\eqref{eq:rel-fa-act}、恢复占时 &
  \implfull{src/reliability/reliability\_assessment.cpp:resolve\_reliability\_params}\\
$\mathrm{PFD}_{avg}\approx\lambda_{DU}T_I/2$、验证试验模型 &
  \implnone（休眠失效 PFD 未建模，仅 $\nu p^d$ 频率约简）\\
两阶段元件 FMEA~\eqref{eq:rel-fa-mode} 前式、储能功率—能量限 &
  \implfull{src/reliability/reliability\_assessment.cpp:run\_distribution\_fmea}\\
故障模式目录（被动/主动、原因—后果） &
  \implfull{src/reliability/failure\_mode.cpp:build\_failure\_mode\_catalog}\\
二阶共故障近似~\eqref{eq:rel-fa-pair} &
  \implpart{加入一阶总量的 N-2 排序近似；精确展开~\eqref{eq:rel-fa-exact} \implnone}\\
Level-1 信息物理双类全期望~\eqref{eq:rel-fa-cyber}、归因、$\eta_A$ &
  \implfull{src/reliability/reliability\_assessment.cpp:run\_distribution\_fmea}（\srcpath{CyberPhysicalFMEAOptions}）\\
三维因子化筛选~\eqref{eq:rel-fa-factor} &
  \implfull{failure\_mode.hpp:CyberPhysicalFMEAOptions::IntelligentFunctionOptions}\\
联合类生成 / 保护逻辑 / FRT 耦合 &
  \implnone（声明 \texttt{joint\_class\_probability\_modelled=false} 等）\\
\end{longtable}
\end{center}

\subsection{参考文献}
\begin{thebibliography}{9}\small
\bibitem{relfa:billinton} R.~Billinton and R.~N. Allan, \emph{Reliability Evaluation of Power Systems},
  2nd ed. New York: Plenum Press, 1996（交替更新过程与元件模型）.
\bibitem{relfa:iec61508} IEC 61508-6, \emph{Functional Safety of Electrical/Electronic/Programmable
  Electronic Safety-Related Systems, Part 6}. Geneva: IEC, 2010（$\mathrm{PFD}_{avg}$ 与验证试验）.
\bibitem{relfa:designdoc} HySim-XJTU-HRPES, \emph{Reliability Mathematical Models and Intelligent
  Cyber-Physical Extension}, \srcpath{docs/theory/reliability\_mathematical\_models\_and\_intelligent\_cyber\_physical\_assessment.md}（第 3、8、9、12 节，本章模型的权威来源）.
\end{thebibliography}


% theory_reliability_restoration.tex —— 三阶段恢复可靠性模型（LinDistFlow 恢复 MILP）
% 本文件为理论层章节，遵循 docs/modules/THEORY_WRITING_GUIDE.md。
% 数学模型依据 docs/theory/reliability_mathematical_models_and_intelligent_cyber_physical_assessment.md 第 4、10 节。

\section{三阶段恢复可靠性模型}\label{sec:rel-rs}

\begin{theorynote}
本章转写\emph{已实现}的三阶段（隔离—开关恢复—修复）恢复可靠性模型，它以每个故障为一次事件，在三个
时段内分别求解\emph{线性配电潮流}（LinDistFlow）恢复 MILP，得到分时段切负荷并按 N-0 反事实加权聚合
为 EENS/LOLE。依据设计文档
\srcpath{docs/theory/reliability\_mathematical\_models\_and\_intelligent\_cyber\_physical\_assessment.md}
第 4、10 节及 \srcpath{include/hacdcpf/analysis/three\_stage\_reliability.hpp}。诚实边界：恢复 MILP 为
线性化（LinDistFlow $+$ 视在功率盒式外松弛），拒绝 LCC 换流器与多端口能量路由器，不含二次损耗与
非线性 AC 圆约束；失败或近似阶段清除 \fld{ok} 与物理有效标志，绝不静默回退。
\end{theorynote}

\subsection{阶段语义与时钟}\label{sec:rel-rs-stage}
对故障 $k$，事件分三阶段（故障元件在三段全程保持开断，仅在 $\tau_{RP}$ 零时长边界“修复”）：
\begin{center}\footnotesize
\begin{tabular}{@{}cL{3.2cm}L{5.0cm}L{4.4cm}@{}}
\toprule
阶段 & 区间 & 物理含义 & 拓扑\\
\midrule
1 & $[0,\tau_{SW}]$ & 故障清除与隔离 & 故障开断；联络额定\\
2 & $[\tau_{SW},\tau_{TP}]$ & 开关恢复 & 故障开断；合格联络优化\\
3 & $[\tau_{TP},\tau_{RP}]$ & 修复窗 & 故障仍开断；保持已接受的 Stage-2 方案\\
\bottomrule
\end{tabular}
\end{center}
时长满足 $r_k=\tau_k^{iso}+\tau_k^{sw}+\tau_k^{rep}$，阶段时长为
$\Delta t_{k,1}=\tau_{SW,k}$、$\Delta t_{k,2}=\tau_{TP,k}-\tau_{SW,k}$、$\Delta t_{k,3}=\tau_{RP,k}-\tau_{TP,k}$，
对应实现字段 $\tau_k^{iso}=\Delta t_{k,1}$、$\tau_k^{sw}=\Delta t_{k,2}$、$\tau_k^{rep}=\Delta t_{k,3}$。

\subsection{交流 LinDistFlow 恢复 MILP}\label{sec:rel-rs-milp}
每阶段联合最小化交直流有功切负荷（末项以 $10^{-8}$ 权仅打破储能/换流器调度简并、保能量）：
\begin{equation}
 \min\ \sum_{i\in\mathcal B^{ac}}w_ip_i^{sh}+\sum_{d\in\mathcal B^{dc}}w_dp_d^{sh}
 +10^{-8}\Big(\sum_bp_b^{st}+\sum_cp_c^{conv}\Big).
\end{equation}
带电与投切：源母线置为已带电根，跟网发电与储能\emph{不能}给死岛带电：
\begin{equation}
 P_i^d-p_i^{sh}\le P_i^dy_i,\ y_i\in\{0,1\};\qquad
 0\le p_i^g\le\overline P_i^gy_i,\quad 0\le p_b^{st}\le\overline P_b^{st}y_{i(b)}.
\end{equation}
有功 / 无功平衡（无功切负荷按功率因数比例 $q_i^{sh}=q_i^{sh}(P_i^d,Q_i^d)$，缺测无功以 0.9 功率因数重构）：
\begin{align}
 \sum_{j:(j,i)}P_{ji}-\sum_{j:(i,j)}P_{ij}+p_i^g+p_i^{sh}+\sum_{b:i(b)=i}p_b^{st}&=P_i^d,\\
 \sum_{j:(j,i)}Q_{ji}-\sum_{j:(i,j)}Q_{ij}+q_i^g+q_i^{sh}&=Q_i^d,\qquad q_i^{sh}=\frac{Q_i^d}{P_i^d}p_i^{sh}.
\end{align}
对平方电压 $v_i$ 与支路投切态 $z_{ij}$，LinDistFlow 压降经 Big-M 随拓扑启用（$v_i$ 为平方标幺电压，
$r_{ij},x_{ij}$ 为系统基标幺阻抗，$P_{ij},Q_{ij}$ 为 MW/Mvar，$S_{base}$ 为 MVA）：
\begin{equation}
 -M(1-z_{ij})\le v_j-v_i+\frac{2r_{ij}}{S_{base}}P_{ij}+\frac{2x_{ij}}{S_{base}}Q_{ij}\le M(1-z_{ij}),
 \qquad \underline V_i^2\le v_i\le\overline V_i^2.
\end{equation}

\begin{remark}[视在功率盒式外松弛]\label{rem:rel-rs-box}
实现对 $P,Q$ 用分离线性界 $|P_{ij}|\le z_{ij}\overline S_{ij}$、$|Q_{ij}|\le z_{ij}\overline S_{ij}$，
是 $P_{ij}^2+Q_{ij}^2\le z_{ij}\overline S_{ij}^2$ 的\emph{外、乐观}松弛：允许角点 $P{=}Q{=}\overline S$
（视在 $\sqrt2\,\overline S$）。决策级模型应改用 SOCP 或多边形内逼近
$P_{ij}\cos\theta_j+Q_{ij}\sin\theta_j\le z_{ij}\overline S_{ij}\cos(\pi/J)$（$j=1,\dots,J$）。
\end{remark}

\subsection{严格辐射能量森林}\label{sec:rel-rs-forest}
辅助商品流 $f_{ij}$ 从电压成形根建立连通性，强制辐射（生成森林）：
\begin{align}
 \sum f_{in}-\sum f_{out}&=y_i\ (i\notin\mathcal S_{root}),\qquad
 \sum f_{in}-\sum f_{out}\le0\ (i\in\mathcal S_{root}),\\
 |f_{ij}|&\le(|\mathcal B|-1)z_{ij},\quad z_{ij}\le y_i,\ z_{ij}\le y_j,\qquad
 \sum_{(i,j)}z_{ij}\le\sum_iy_i-|\mathcal S_{root}|.
\end{align}
故障元件在每阶段固定开断；常开恢复联络在 Stage 1 固定开断、Stage 2 优化、Stage 3 保持已接受方案；
开关预算 $\sum_{\ell\in\mathcal L^{NO}}z_\ell\le K^{sw}$。保护 / 联锁预检验证保护可切断故障电流、分段动作
只在所需上游跳闸后发生；非法动作序列阻断恢复准入并上报。

\subsection{事件内储能时序}\label{sec:rel-rs-storage}
各阶段时长 $\Delta t_s$ 下储能荷电按序演化（放电受额定与剩余可交付能量双限）：
\begin{equation}
 E_{b,s+1}=E_{b,s}-p_{b,s}^{st}\Delta t_s,
\end{equation}
在稳定求解序 $[\text{AC 储能}\mid\text{移动储能}\mid\text{DC 储能}]$ 上一致施加；\texttt{InTransit} 移动储能
电可用度为零、不列为可用源停运。Stage 3 仅在储能能量限改变时重解，否则复用 Stage-2 拓扑 / 解。

\subsection{耦合直流 LinDistFlow 与换流器}\label{sec:rel-rs-dc}
默认混合路径在同一阶段 MILP 内建直流网络，AC/DC 母线映射域限定（同号 ID 不混叠）。DC 支路
$v_e-v_d+\tfrac{2r_{de}}{S_{base}^{dc}}P_{de}=0$，含 Big-M 投切、带电端点约束与严格 DC 辐射森林。
VSC 以 AC$\to$DC 传输 $p_v^+$、DC$\to$AC 传输 $p_v^-$、恒效率 $\eta_v$ 表述：
\begin{equation}
 P_v^{ac}=-p_v^++\eta_vp_v^-,\qquad P_v^{dc}=\eta_vp_v^+-p_v^-,\qquad 0\le p_v^+,p_v^-\le\overline P_v.
\end{equation}
小正传输代价消去对冲简并。DC-DC 换流器用类似双向式；DC 静发 / PV / 储能带停运、功率、电压根与能量约束
入直流平衡。VPP 无稳定成员引用列表，故作单一独立边界注入，不静默展开或重复计。LCC 换流器与多端口能量
路由器因换相 / 端口平衡方程不等价于 VSC 而被拒绝。

\subsection{频率加权与 N-0 资格}\label{sec:rel-rs-n0}
设 $s_{i,k,s}^{raw}$ 为故障 $k$ 下切负荷、$s_{i,s}^{N0}$ 为同阶段时长 / 能量条件下的健康态切负荷，可靠性
增量与指标为
\begin{equation}
 s_{i,k,s}^{inc}=\max(0,s_{i,k,s}^{raw}-s_{i,s}^{N0}),\quad
 \mathrm{EENS}_i=\sum_k\lambda_k\sum_ss_{i,k,s}^{inc}\tau_{k,s},\quad
 \mathrm{LOLE}=\sum_k\lambda_k\sum_s\tau_{k,s}\mathbf 1\!\Big[\sum_is_{i,k,s}^{inc}>\varepsilon\Big].
\end{equation}
N-0 反事实自\emph{同一}初始储能能量出发、用\emph{同一}三阶段时长，但演化\emph{独立}的健康调度 / 能量轨迹
$s_{i,k,s}^{inc}=\max(0,\ s_i(G_{k,s},E_{k,s})-s_i(G_{0,s},E_{0|k,s}))$，其中 $E_{0|k,s}$ 是健康反事实在同一
外生时长序列下达到的能量态（非故障耗竭态）。这估计故障的\emph{总因果效应}（含其对储能使用的影响），
而非纯拓扑边际；仅当两路径共享初始 SOC、负荷、可再生可用性、阶段时长与所有非故障设置时有效。缺省故障集
含 AC/DC 支路，发电机 / 变压器 / 换流器 / 开关 / DER-微网族由选项控制以复现历史仅支路结果。依据：
\srcpath{src/reliability/three\_stage\_reliability.cpp:run\_three\_stage\_reliability}。

\subsection{与实现的对应}\label{sec:rel-rs-impl}
\begin{center}\footnotesize
\begin{longtable}{@{}L{5.4cm}L{9.0cm}@{}}
\toprule
\textbf{理论对象} & \textbf{平台实现状态}\\
\midrule
\endfirsthead
\toprule
\textbf{理论对象} & \textbf{平台实现状态}\\
\midrule
\endhead
\bottomrule
\endfoot
三阶段语义与时钟 & \implfull{include/hacdcpf/analysis/three\_stage\_reliability.hpp}\\
交流 LinDistFlow 恢复 MILP、能量森林、开关预算 &
  \implfull{src/reliability/three\_stage\_reliability.cpp:run\_three\_stage\_reliability}\\
事件内储能时序（AC/移动/DC 序） & \implfull{src/reliability/three\_stage\_reliability.cpp}\\
耦合 DC LinDistFlow $+$ VSC 双向传输 & \implfull{src/reliability/three\_stage\_reliability.cpp}（\fld{include\_dc\_power\_flow}）\\
N-0 反事实增量与 EENS/LOLE 聚合 & \implfull{src/reliability/three\_stage\_reliability.cpp}\\
视在功率 SOCP / 多边形内逼近 & \implnone（盒式外松弛，见注记~\ref{rem:rel-rs-box}）\\
LCC 换流器、多端口能量路由器、二次损耗 & \implnone（显式拒绝 / 线性化边界）\\
\end{longtable}
\end{center}

\subsection{参考文献}
\begin{thebibliography}{9}\small
\bibitem{relrs:baran} M.~E. Baran and F.~F. Wu, ``Optimal sizing of capacitors placed on a radial
  distribution system,'' \emph{IEEE Trans. Power Del.}, vol.~4, no.~1, pp.~735--743, 1989（LinDistFlow）.
\bibitem{relrs:designdoc} HySim-XJTU-HRPES, \emph{Reliability Mathematical Models and Intelligent
  Cyber-Physical Extension}, \srcpath{docs/theory/reliability\_mathematical\_models\_and\_intelligent\_cyber\_physical\_assessment.md}（第 4、10 节）.
\end{thebibliography}


% theory_reliability_cyber_intelligent.tex —— 智能信息物理可靠性模型（提出）
% 本文件为理论层章节，遵循 docs/modules/THEORY_WRITING_GUIDE.md。
% 数学模型依据 docs/theory/reliability_mathematical_models_and_intelligent_cyber_physical_assessment.md 第 13–19 节。

\section{智能信息物理可靠性模型（提出层）}\label{sec:rel-icp}

\begin{theorynote}
本章转写\emph{提出}（proposed）的智能信息物理可靠性框架，它把已实现的 Level-1 FMEA
（\S\ref{sec:rel-fa-cyber}）从“自动化可用/不可用”双类，细化为三耦合图上的检测—隔离—恢复—保护函数
可靠性，并以阶段概率 $\pi_c(k)$、阶段时长 $\tau_{k,c,s}$、后果补丁 $\Phi_{k,c}$ 三通道进入可靠性指标。
依据设计文档
\srcpath{docs/theory/reliability\_mathematical\_models\_and\_intelligent\_cyber\_physical\_assessment.md}
第 13--19 节。\textbf{诚实边界}：除第 17 节的静态事件抽象为\emph{部分实现}外，本章绝大部分为\emph{提出
模型，尚未在平台实现}；就地以“未实现扩展说明”标注，不得与已实现的 §12 混淆。核心方法论警示贯穿全章：
共享传感器 / 通信 / 电源 / 控制中心的路径\emph{不可}按独立相乘，须用联合分布 / 最小割集 / BDD 评估。
\end{theorynote}

\subsection{为何需要集成模型}\label{sec:rel-icp-why}
纯物理问题问“给定失效设备，存活网络能供多少负荷”；集成问题更广：故障是否被检测、多快？失效段是否被
正确隔离？保护是否无越级误动地清除目标区？哪些恢复动作可观测、可通信、可准入、可成功执行？微网 / DER
能否成岛并维持至修复？多数信息失效\emph{间接}影响可靠性——延长定位 / 开关时间、扩大隔离区、移除可行
恢复动作、冻结 DER 定值、抑制构网能力——故归入阶段概率、阶段时长与后果补丁。

\subsection{三耦合图与联合状态}\label{sec:rel-icp-graphs}
系统表为 $\mathcal G^{ICP}=(G_p,G_c,G_f,\Pi_{pc},\Pi_{cp})$：物理网 $G_p$、信息 / 通信网 $G_c$、
检测 / 隔离 / 恢复 / 保护的功能依赖图 $G_f$、信息服务$\to$受控物理资产映射 $\Pi_{pc}$、物理母线$\to$信息
资产供电映射 $\Pi_{cp}$。联合状态 $\xi_t=(x_t^p,x_t^c,e_t^c,b_t,\hat x_t,\mathcal A_t)$ 含物理 / 信息元件态、
链路质量、备用能量、故障信念 $\hat x$ 与可执行动作集 $\mathcal A_t$。

\subsection{信息功能可用度}\label{sec:rel-icp-avail}
对所需功能 $f$ 与信息态 $x^c$，布尔结构函数 $\phi_f(x^c)\in\{0,1\}$ 指示“存在有效的
传感器—控制器—执行器服务路径”，功能可用度为
\begin{equation}
 A_f=\Pr[\phi_f(X^c)=1]=\sum_{x^c}\phi_f(x^c)\Pr(X^c=x^c).
\end{equation}

\begin{proposition}[路径结构与独立性陷阱]\label{prop:rel-icp-path}
串链 $A_f=\prod_qA_q$；冗余路径 $A_f=1-\prod_{p\in\mathcal P_f}(1-A_p)$ \emph{仅当路径独立}成立。一般
（路径共享开关 / 电源 / 控制中心 / 时钟 / 管道）须按联合信息态分布评估
\begin{equation}
 A_f=\Pr\!\Big(\bigvee_{p\in\mathcal P_f}\bigwedge_{q\in p}\{X_q^c=1\}\Big),
\end{equation}
用最小割集 / 二元决策图 / 蒙特卡洛。对 FLISR，须同一信息态下三功能同时成立
$A_k^{FLISR}=\Pr[\phi_D\land\phi_I\land\phi_R]$；分别算 $A_DA_IA_R$ 相乘在共享基础设施时\emph{一般错误}。
\end{proposition}

\paragraph{通信质量（QoS），不止可用度。}
快速保护 / FLISR 的可用命令路径须满足端到端时延 / 抖动 / 丢包界
$L_p\le L_f^{max}$、$J_p\le J_f^{max}$、$P_{loss,p}\le P_f^{max}$；服务事件
$\mathcal Q_{f,p}=\{\phi_{f,p}{=}1,L_p{\le}L_f^{max},J_p{\le}J_f^{max},P_{loss,p}{\le}P_f^{max}\}$，
QoS 感知可用度 $A_f^{QoS}=\Pr(\bigvee_p\mathcal Q_{f,p})$。仅比较时延 / 丢包的\emph{均值}不是服务可用度模型；
保护流量需毫秒界、恢复控制容秒级、监视可容更长——单一全局通信可用度标量无法表征这些不同契约。

\subsection{智能故障检测}\label{sec:rel-icp-detect}
对真故障类 $k$ 与观测 $z$，检测器产生告警 $\hat k$，混淆矩阵 $C^D_{k\hat k}=\Pr(\hat k\mid k)$，
检出 / 漏报 / 误报为 $P_D(k)=1-C^D_{k0}$、$P_{miss}(k)=C^D_{k0}$、$P_{FA}=\Pr(\hat k\ne0\mid k{=}0)$。
误报须先定暴露基：以 $\nu_{eval}$ 个独立决策窗 / 年，$\lambda_{FA}=\nu_{eval}P_{FA}$；连续监测 / 时序相关
检测器须用告警 / 年实测率或点过程模型，把每采样步当独立机会会把误跳高估数量级。检测须条件于可观测性 $o$
与工况 $h$：$C^D_{k\hat k}(o,h)$。检测时延用分布 $T_k^D\sim F_k^D(t\mid o,h)$、$\tau_k^{iso}=T_k^D+T_k^{dec}+T_k^{trip}$。
因后果对时延非线性，$\mathbb E[E_k(T_k^D)]\ne E_k(\mathbb E[T_k^D])$，正确贡献为
\begin{equation}
 \mathrm{EENS}_k=\lambda_k\int E_k(t)\,\mathrm dF_k^D(t),
\end{equation}
有限互斥时延类是该积分的实用求积。检测经 $\pi_c(k),\tau_{k,c,s},\Phi_{k,c}$ 进入可靠性（及时 / 迟 / 漏 /
误报四后果类）。

\subsection{智能故障隔离}\label{sec:rel-icp-isolate}
设候选故障段集 $\mathcal Z$，隔离引擎给后验信念 $b_q=\Pr(K{=}q\mid z,G_p,G_c)$（$\sum_qb_q=1$），
以\emph{决策论}权衡欠隔离风险与过大停电：
\begin{equation}
 \hat q=\arg\min_{q\in\mathcal Z}\big[c_{miss}\Pr(K\notin q\mid z)+c_{shed}P_{isolated}(q)\big],
\end{equation}
受保护选择性与开关遮断能力约束。隔离成功概率 $p_I(k,c)=\Pr[\bigwedge_{q\in Q_k}(\text{命令路径可用})_q\land
(\text{器件动作})_q]$ 决定主区 $Z_k^{primary}$ / 后备区 $Z_k^{backup}$，共享通信路径与共直流电源须联合评估。

\subsection{智能恢复}\label{sec:rel-icp-restore}
恢复策略 $\pi_R:(\hat x_t,G_p,G_c,E_t^{st})\mapsto a_t$ 在部分可观测下选开关 / DER 调度 / 成岛 / 投荷。
单步风险感知式为
\begin{equation}
 \min_{a,y}\ \sum_iw_ip_i^{sh}+c_{sw}\sum_\ell|z_\ell-z_\ell^0|+c_{risk}\,\mathcal R(a,\hat x),\qquad
 \mathcal R(a,\hat x)=\Pr(g(a,X)>0\mid\hat x),
\end{equation}
$g$ 编码不完全隔离 / 过载 / 不安全带电 / 保护失配 / 电压越限 / 同步失败；受物理恢复模型、保护联锁与命令 /
动作可用性约束（$z_\ell^{close}\le\phi_{cmd(\ell)}$、$p_r\le\overline P_r\phi_{ctrl(r)}$、$y_{island}\le\phi_{GFM}x_{GFM}^p$）。
命令路径不可用未必使动作不可行，可转为带行程 / 操作时延的人工动作，故每动作兼具可行性与时序
$T_a\in\{T_a^{remote},\ T_a^{crew}+T_a^{manual},\ \infty\}$。

\begin{remark}[策略成功与筛选近似]\label{rem:rel-icp-policy}
须区分优化器可行性与智能策略性能：$p_R^{valid}=\Pr(a\in\mathcal A_{safe}\mid\hat x,x)$、
$p_R^{exec}=\Pr(\text{全部所需动作执行}\mid a,x^c)$。自动恢复类概率
$\pi_{auto}(k)=\Pr(D_k\cap I_k\cap R_k^{valid}\cap R_k^{exec}\mid k)$，\emph{仅在条件独立}下才可因子化为
$\Pr(D_k)\Pr(I_k)\Pr(R_k^{valid})\Pr(R_k^{exec})$。故乘积
$A_k^{FLISR}P_D p_I p_R^{valid}p_R^{exec}$（即已实现 §12 的三维因子化，式~\eqref{eq:rel-fa-factor}）只是\emph{筛选}
近似，会重复计共享依赖；无效 / 未验证动作须路由到保守人工 / 阻断类，不得计为成功恢复。多步恢复用约束
POMDP / 模型预测控制，学习型策略须过确定性安全滤波（辐射 / 电压 / 热限 / 同步 / 联锁 / DER 能力）方可执行。
\end{remark}

\subsection{智能保护（静态事件抽象为部分实现）}\label{sec:rel-icp-protect}
保护有两类失效族：\emph{可靠性}（dependability，真需求时拒动）与\emph{安全性}（security，无区内故障误动）。
对保护区 $z$，$\nu_z=\sum_{k\in\mathcal K_z}\lambda_k$、拒动频率 $\lambda_z^{FT}=\nu_zp_z^{FT}$，其后果与误动为
\begin{equation}
 \mathrm{EENS}^{FT}=\sum_z\lambda_z^{FT}\sum_sS_{z,s}^{backup}\tau_{z,s}^{backup},\qquad
 \mathrm{EENS}^{NT}=\sum_z\lambda_z^{NT}\sum_sS_{z,s}^{trip}\tau_{z,s}^{trip}.
\end{equation}
拒动概率的完整需求成功事件为
\begin{equation}
 \mathcal S_z=R_z\cap M_z\cap C_z\cap B_z\cap S_z'\cap A_z,\qquad
 p_z^{FT}=\Pr(\mathcal S_z^c\mid\text{保护需求}),
\end{equation}
其中 $R_z$ 继电逻辑健康、$M_z$ 量测有效、$C_z$ 必要时通信 / 跳闸路径及时、$B_z$ 断路器机械动作、
$S_z'$ 正确选择性定值、$A_z$ 站用电可用。
本地主保护\emph{不应}被强加广域通信依赖，否则系统性高估信息—保护耦合。自适应保护（微网模式、逆变限流、
拓扑变）以失配概率 $p_z^{mis}=\Pr(\theta_z\ne\theta_z^*(\sigma)\lor T_{update}>T_{required})$ 表述三模式
（定值更新通信失、错误定值选择、起动电流不足），其后果由短路 / 保护模型定后备区后作区变异与时长类入可靠性。

\subsection{保护—FRT 可靠性接口与微网（提出细化层）}\label{sec:rel-icp-frt}
设计文档第 18 节以约 $16$ 个子模型把第 17 节的事件概率、清除时间、拓扑变异与 DER 后果从\emph{继电器与故障
穿越（FRT）事件逻辑}细化：混合“保护—FRT”状态、DER 电压 / 频率穿越自动机（穿越 / 瞬时闭锁 / 跳闸为\emph{不同}
状态）、故障期逆变电流与 GFL/GFM 行为、定时 / 反时限动作量、主 / 后备协调与断路器失灵、重合闸 / 熔断器 /
分段器 / 反孤岛序列、DER-保护闭环与重连后可用度、保护—FRT 后果类、IEEE 1547 与保护参数接口、保真度阶梯与
所需有效标志。第 19 节给微网 / DER 物理可靠性：孤岛可行性（构网源、功率 / 能量平衡、保护可行）与储能续航
$P_b^{avail}=\min(P_b^{rated},(E_b-E_b^{min})\eta_b^{dis}/\tau_s)$。
\begin{gapnote}
第 13--19 节（除 §17 静态事件抽象部分实现外）为\emph{提出}模型，\emph{尚未}在平台可靠性评估器中实现：
无显式通信拓扑 / 时延 / 信息节点电源 / 共因信息失效 / 学习型 FDIR 策略；第 18 节的继电器起动、后备区扩展、
断路器失灵、DER-FRT 轨迹与联合类生成\emph{均未}耦合入当前评估器（已实现结果声明
\texttt{protection\_logic\_modelled=false}、\texttt{protection\_frt\_reliability\_coupled=false}）。本章用于
界定研究路线与诚实上界，不得当作平台当前算法。第 18 节细化不改变当前可靠性评估器的实现状态。
\end{gapnote}

\subsection{与实现的对应}\label{sec:rel-icp-impl}
\begin{center}\footnotesize
\begin{longtable}{@{}L{5.4cm}L{9.0cm}@{}}
\toprule
\textbf{理论对象} & \textbf{平台实现状态}\\
\midrule
\endfirsthead
\toprule
\textbf{理论对象} & \textbf{平台实现状态}\\
\midrule
\endhead
\bottomrule
\endfoot
三耦合图 $\mathcal G^{ICP}$、联合状态、QoS 可用度 & \implnone（提出）\\
信息功能结构可用度 $A_f$、FLISR 联合、最小割集 & \implnone（提出；已实现仅 §12 独立因子筛选）\\
智能检测混淆矩阵 / 时延积分 EENS$_k$ & \implnone（提出）\\
智能隔离决策论选区 / $p_I$ & \implnone（提出）\\
智能恢复风险感知 MILP / POMDP / $p_R$ & \implnone（提出；已实现三阶段 MILP 见 \S\ref{sec:rel-rs}）\\
智能保护拒动 / 误动事件抽象 $\lambda_z^{FT}$ & \implpart{静态事件抽象经故障模式 FMEA；动态协调提出}\\
保护—FRT 细化（§18）、微网孤岛 / 储能续航（§19） & \implnone（提出细化层；储能续航式已用于恢复 MILP）\\
\end{longtable}
\end{center}

\subsection{参考文献}
\begin{thebibliography}{9}\small
\bibitem{relicp:ieee1547} IEEE Std 1547-2018, \emph{IEEE Standard for Interconnection and
  Interoperability of Distributed Energy Resources}. New York: IEEE, 2018.
\bibitem{relicp:ieee2030} IEEE Std 2030.7-2017, \emph{IEEE Standard for the Specification of
  Microgrid Controllers}. New York: IEEE, 2017.
\bibitem{relicp:designdoc} HySim-XJTU-HRPES, \emph{Reliability Mathematical Models and Intelligent
  Cyber-Physical Extension}, \srcpath{docs/theory/reliability\_mathematical\_models\_and\_intelligent\_cyber\_physical\_assessment.md}（第 13--19 节）.
\end{thebibliography}

\section{保护逻辑与 DER 故障穿越的可靠性接口}\label{sec:rel-pf}

\begin{theorynote}
本章围绕平台\emph{既有}动态（IEEE~1547 块）给出已落地的 \textbf{L1 轨迹类可靠性接口}：把继电器轨迹、
时限/反时限协调、断路器机械与 DER 故障穿越（FRT）事件，条件化为互斥的\emph{后果类}并折算为
EENS/LOLE/LOLF。给定轨迹入口消费代表性电流、DER 电压/相角轨迹和类条件切负荷阶段；在线入口由 Mass-Matrix DAE
生成故障电流与 DER 轨迹并反馈支路跳闸。自动拓扑、重合序列、瞬时闭锁和同步窗分别由可执行函数承担。
\end{theorynote}

\begin{boundarynote}
\textbf{实现状态（诚实边界，依据设计文档 \S18.1）}：下表区分给定轨迹、在线动态和拓扑后果的保真度。
除非结果显式报告 \fld{protection\_frt\_reliability\_coupled=true}，动态 IEEE~1547 与短路能力是\emph{独立}
分析，\emph{不}贡献该结果的 EENS/LOLE/LOLF/SAIDI 或恢复阶段 DER 可用性。
\begin{center}\footnotesize
\begin{tabular}{@{}L{5.7cm}L{2.5cm}L{5.5cm}@{}}
\toprule
能力 & 状态 & 可靠性含义\\
\midrule
DER 电压/频率越限计时、跳闸、重连合格与功率爬坡 & 已实现（选项，动态）& L1 类生成复用 IEEE~1547 状态机；通用 FMEA/MC 不自动调用\\
GFL/GFM 动态、电流限幅与优先级、volt-var/freq-watt & 已实现（选项，动态）& 供暂态研究，未转为列联类概率或 EENS 阶段\\
短路换流器贡献（GFL 电流源、GFM 阻抗后电压源）& 短路分析已实现 & 静态故障快照，非继电/FRT 事件演化\\
拒动/误动/拒开合、保护区扩展、重合概率 & 可靠性事件抽象已实现 & 作为\emph{输入}被故障模式 FMEA 与三阶段消费；启动值与选择性\emph{未}导出\\
反时限动作积分、复位、清除延时链、主/后备/断路器失灵事件树 & \textbf{L1 已实现} & 由给定轨迹生成互斥类和协调诊断\\
在线距离/差动、重合序列、DER--保护双向反馈与 DAE & \textbf{已实现} & 给定轨迹或在线 DAE 生成互斥类并进入年度事件树\\
\bottomrule
\end{tabular}
\end{center}
\end{boundarynote}

\begin{boundarynote}
\textbf{控制器—物理元件与部分可观测边界}：换流器控制器只与其所属 VSC/IBR 的端量传感器和电流/模式执行口关联。
保护功能不按故障元件强制一对一，而由保护区、CT/PT/状态接点量测集和受控断路器/接触器集共同定义；主后备可以区域重叠并动作于不同开断设备。
系统级恢复只能消费经可用信息通道送达的遥测、位置接点、保护事件和电源资格证书，不能直接读取 DAE 的真实完整状态。不同真实状态可产生相同报告；必需观测缺失时，遥控动作或电源资格不得准入，不得以内部真值补全。平台的有限 POMDP 是独立能力；除非某结果显式调用并声明该模型，不得将其冒充为通用可靠性恢复入口的当前行为。
\end{boundarynote}

\subsection{混合保护--穿越状态}\label{sec:rel-pf-state}
\begin{definition}[事件时状态]
事件时刻的混合状态为
\begin{equation}
 \zeta(t)=\big(x^{dyn}(t),\,\sigma(t),\,\chi(t),\,\kappa(t),\,\Omega(t),\,s^{br}(t),\,s^{rec}(t),\,\vartheta(t)\big),
\end{equation}
分别为电气/控制器微分状态、拓扑与并/离网模式、DER 连接态 $\chi_r\in\{0,1\}$、复合 DER 模式
$\kappa_r=(\kappa_r^{ctrl},\kappa_r^{FRT},\kappa_r^{lim})$、保护动作量/计时器 $\Omega_z$、断路器状态
$s_z^{br}$、重合器状态 $s_z^{rec}$ 与激活定值组 $\vartheta_z$。
\end{definition}
以电压角 $\delta_i$ 记电气量、以 $\vartheta_z$ 记保护定值，避免同符号歧义。事件间遵循微分代数
\begin{equation}
 \dot x^{dyn}=F_{\sigma,\kappa}(x^{dyn},y,u,w),\qquad 0=G_{\sigma,\kappa}(x^{dyn},y,u,w),
\end{equation}
$y$ 为网络代数量（电压、支路电流、接口功率），$u$ 为控制输入、$w$ 为扰动；离散守卫在
$g_e(\zeta(t^-))=0$ 时触发 $\zeta(t^+)=\mathcal R_e(\zeta(t^-))$，产生跳闸、模式切换、断路器动作、
重合与重连。复合模式使 GFL/GFM、穿越/跳闸、限幅/未限幅\emph{不}互斥：一台 DER 可同时为 GFM 构网、
强制穿越且电流受限。

\subsection{DER 电压/频率穿越自动机}\label{sec:rel-pf-frt}
设 DER $r$ 端正序电压 $V_r=|V_{i(r)}^{+}|$、量测频率 $f_r$。对配置带 $j$，其有效谓词 $h_{r,j}\in\{0,1\}$
与\emph{逐带}连续越限计时器
\begin{equation}
 D_{r,j}(t)=\begin{cases}t-\sup\big(\{s\in[t_0,t]:h_{r,j}(s)=0\}\cup\{t_0\}\big),&h_{r,j}(t)=1,\\[2pt]0,&h_{r,j}(t)=0,\end{cases}
\end{equation}
带条件清除即复位；这比把所有异常时间累加为一个标量严格，后者会错误合并不同激励与不同带。配置跳闸守卫
\begin{equation}
 \mathcal G_r^{trip}(t)=\bigvee_{j\in\mathcal J_r^{trip}}\big[D_{r,j}(t)\ge T_{r,j}^{trip}\big],\qquad
 T_{r,j}^{trip}=\mathcal T_{r,j}(\text{类别},\text{接入规程},\text{整定},\text{固件证据}).
\end{equation}
\emph{类别本身不决定跳闸曲线}：$T_{r,j}^{trip}$ 须落在适用要求内但由实际整定给出。模式转移与连接态为
\begin{equation}
 \kappa_r(t^+)=\Psi_r\big(\kappa_r(t^-),V_r,f_r,\{D_{r,j}\}_j,\eta_r^{device}\big),\quad
 \chi_r(t^+)=\begin{cases}0,&\mathcal G_r^{trip}=1,\\1,&\mathcal G_r^{reconnect}=1,\\ \chi_r(t^-),&\text{否则}.\end{cases}
\end{equation}
重叠带谓词须在 $\Psi_r$ 中给出显式优先级，否则两个同时激活的带会规定不兼容动作。

\subsection{穿越、瞬时闭锁与跳闸是不同状态}\label{sec:rel-pf-output}
DER 出力须按动态模式条件化：
\begin{equation}
 (p_r,q_r)=\begin{cases}
 (p_r^{cmd},q_r^{cmd}),&\kappa_r^{FRT}=\text{正常},\\
 (p_r^{FRT},q_r^{FRT}),&\kappa_r^{FRT}=\text{穿越},\\
 (p_r^{MC},q_r^{MC}),&\kappa_r^{FRT}=\text{瞬时闭锁},\\
 (0,0),&\kappa_r^{FRT}\in\{\text{跳闸},\text{闭锁},\text{重连等待}\},\\
 a_r(t)(p_r^{avail},q_r^{avail}),&\kappa_r^{FRT}=\text{爬坡}.
 \end{cases}
\end{equation}
瞬时闭锁策略 $(p_r^{MC},q_r^{MC})$ 是设备相关的（可为零、可保无功支撑），\emph{不得}假定所有并网 DER
都支撑电压，也\emph{不得}把瞬时闭锁等同于物理断开。$\Psi_r$ 须给出瞬时闭锁的进入/退出守卫（含恢复驻留、
迟滞、延迟升级为跳闸），否则穿越$\leftrightarrow$瞬时闭锁会在带边界颤振而终端类不唯一。

\subsection{故障期逆变器电流与 GFL/GFM}\label{sec:rel-pf-current}
公共电流限幅圆 $i_{d,r}^2+i_{q,r}^2\le(I_r^{max})^2$。有功优先时
\begin{equation}
 i_{d,r}=\operatorname{sat}(i_{d,r}^{ref},I_r^{max}),\quad
 i_{q,r}=\operatorname{sat}\!\Big(i_{q,r}^{ref},\sqrt{(I_r^{max})^2-i_{d,r}^2}\Big),
\end{equation}
无功优先时对调 $d/q$，其中 $\operatorname{sat}(x,\bar x)=\operatorname{sgn}(x)\min(|x|,\bar x)$。平衡三相
RMS 下故障期功率受\emph{电压相关}圆约束
\begin{equation}
 p_r^2+q_r^2\le(\sqrt3\,V_{LL,r}I_r^{max})^2\ \xrightarrow{\text{标幺}}\ p_r^2+q_r^2\le V_r^2(I_r^{max})^2,
\end{equation}
\emph{不}等于稳态铭牌圆；混用相/线电压与单相/三相功率会引入 $\sqrt3$ 或 $3$ 的因子错误。GFL 为 PLL 同步的
受控\emph{电流}源；GFM 为滤波/虚拟阻抗后的受控\emph{电压}源\emph{直至限幅改变有效模式}：
\begin{equation}
 \kappa_r=\begin{cases}(\mathrm{GFM},\kappa_r^{FRT},\text{电压构网}),&|i_r^{ref}|\le I_r^{max},\\
 (\mathrm{GFM},\kappa_r^{FRT},\text{电流受限}),&|i_r^{ref}|>I_r^{max}.\end{cases}
\end{equation}
保护到达与孤岛支撑必须用\emph{限幅后}模式；标签 \texttt{GFM} 本身不认证无限电压支撑或稳定孤岛。

\subsection{保护动作量：定时限与反时限}\label{sec:rel-pf-relay}
定时限（连续启动语义）$\dot D_z=1$（$\rho_z=1$）或 $-D_z/T_z^{reset}$（$\rho_z=0$），并在 $D_z\ge T_z^{delay}$
出令。反时限过流采用通用曲线并\emph{对动作量积分}以应对故障期电流变化：
\begin{equation}
 T_z^{curve}(M)=\mathrm{TMS}_z\!\left[\frac{A_z}{M^{p_z}-1}+B_z\right]+T_z^{add},\quad
 \dot\Omega_z=\begin{cases}1/T_z^{curve}(M),&M>1,\\ -\Omega_z/T_z^{reset},&M\le1,\end{cases}
\end{equation}
$M=|\widetilde I_z|/I_z^{pickup}$，出令时刻 $T_z^{relay}=\inf\{t:\Omega_z(t)\ge1\}$。当 DER 限幅/跳闸/电压
支撑或拓扑改变在继电器动作前改变电流时，该积分是\emph{必要}的。电气清除\emph{不}在出令时刻发生，清除链为
\begin{equation}
 T_z^{clear}=T_z^{relay}+T_z^{channel}+T_z^{tripcoil}+T_z^{mechanical}+T_z^{arc},
\end{equation}
$T_z^{channel}$ 仅在方案确用通信通道（导引/传输跳闸）时非零且条件于通信服务态；纯本地硬接线取 $0$。
所有时间量须用单一时基（通常秒），曲线系数须遵循选定 IEC/IEEE/厂商族而不混用。

\subsection{主后备协调与断路器失灵}\label{sec:rel-pf-coord}
对故障 $k$，$T_k^{pri}=\min_{z\in\mathcal P_k^{pri}}T_z^{clear}$、$T_k^{bak}=\min_{z\in\mathcal P_k^{bak}}T_z^{clear}$
（永不有效动作者取 $+\infty$）。须把\emph{物理清除结果}与\emph{协调诊断}分离：主清除事件 $\mathcal E_k^{pri}$
与满意协调 $\mathcal S_k^{coord}=\mathcal E_k^{pri}\cap\{T_k^{bak}-T_k^{pri}\ge\Delta T_k^{coord}\}$——主设备可
清除故障而后备裕度不足，这是\emph{带协调告警属性的主清除}，非未分类样本。跳令后电流持续即断路器失灵
\begin{equation}
 \mathcal F_z^{BF}=\big\{u_z^{trip}(t_z^{cmd}){=}1,\,s_z^{br}\ne\text{开},\,\inf_{\tau\in[t_z^{cmd},t_z^{cmd}+T_z^{BF}]}|\widetilde I_z(\tau)|>I_z^{open}\big\},
\end{equation}
启动失灵后备并扩展清除区；机械拒开概率是\emph{另一}按需失效，若已并入标定的失灵事件概率则\emph{不}再独立
相乘。原始事件可重叠，故须由显式事件树按优先级赋类：$(1)$ 若终端仍带电则\emph{拒清};$(2)$ 否则若任一非必要
健全区被跳则\emph{误协调};$(3)$ 否则按主/后备清除赋类。

\subsection{保护--FRT 后果类与 EENS}\label{sec:rel-pf-consequence}
把接口类扩展为六元组
\begin{equation}
 c=(c_D,c_I,c_P,c_{FRT},c_{rec},c_R),\qquad
 \pi_c(k)=\Pr(C_D{=}c_D,\dots,C_R{=}c_R\mid K{=}k),
\end{equation}
联合表示检测、隔离、保护、穿越、重连与恢复结果。设故障前暴露 DER 集 $\mathcal R_k$ 与终端分类时刻
$T_k^{term}=\min(T_k^{clear},T_k^{backup,max})$。以\emph{跳闸优先}定义互斥 FRT 划分
\begin{align}
 \mathcal E_k^{FRT,trip}&=\{\exists r\in\mathcal R_k:\chi_r(T_k^{term})=0\},\\
 \mathcal E_k^{FRT,MC}&=\{\chi_r(T_k^{term}){=}1\,\forall r\}\cap\{\exists r:\kappa_r^{FRT}(T_k^{term})=MC\},\\
 \mathcal E_k^{FRT,ride}&=\{\chi_r(T_k^{term}){=}1\,\forall r\}\cap\{\kappa_r^{FRT}(T_k^{term})\ne MC\,\forall r\};
\end{align}
$\mathcal R_k=\varnothing$ 时置 $c_{FRT}=\text{不适用}$ 而非依赖空真。\textbf{关键警示}：一台小光伏跳闸与一台
关键 GFM 储能跳闸\emph{不能}仅因同属 $\mathcal E_k^{FRT,trip}$ 而赋同一后果——聚合标签仅在 $\mathcal R_k$ 内
DER 后果等价、或物理引擎同时接收器件级轨迹 $\{\chi_r,\kappa_r,a_r\}$ 时才充分。类条件物理后果与 FRT 条件
EENS 增量为
\begin{equation}
\begin{aligned}
 S_{k,c}(t)&=S\big(\Phi_{k,c_P,c_{FRT},c_I}(G_0),\mathcal A_{c_R},
   \chi(t),\overline p^{avail}(t)\big),\\
 \Delta\mathrm{EENS}^{PFRT}&=\sum_k\lambda_k\!\sum_c\pi_c(k)
   \!\int_{\mathcal T_k}\!S_{k,c}(t)\,\mathrm dt,
\end{aligned}
\end{equation}
即\emph{动态保护/FRT 条件} EENS 与\emph{静态 DER 恒可用}基线之差。代表性类：主清除$\times$穿越/瞬时闭锁/
跳闸、后备清除、误协调、拒清（失灵为致因属性）、暂态自清除、重合闭锁、孤岛保护无效。

\subsection{微网孤岛可行性、储能续航与同步}\label{sec:rel-pf-microgrid}
DER 侧后果的\emph{岛级}闭合由物理微网模型给出。微网岛 $q$ 可服务须有可用构网锚
\begin{equation}
 y_q^{island}\le\sum_{r\in\mathcal R_q^{GFM}}x_r^p\,\phi_{ctrl(r)}(x^c),
\end{equation}
二值锚可用积已隐含“至少一个锚可用”，MILP 实现须为每个积 $x_r^p\phi_{ctrl(r)}$ 引入辅助二值。有功/无功
充裕度为 $\sum_{r}p_r^{DER}+p_q^{st}+p_q^{sh}=P_q^d+P_q^{loss}$、$\sum_{r}q_r^{DER}+q_q^{sh}=Q_q^d+Q_q^{loss}$，
换流器能力用二阶锥 $(p_r^{DER})^2+(q_r^{DER})^2\le(S_r^{rated})^2$；LP/MILP 须声明其多边形是\emph{保守内}逼近
\begin{equation}
 p_r\cos\theta_j+q_r\sin\theta_j\le S_r^{rated}\cos(\pi/J),\quad j=1,\dots,J,
\end{equation}
还是\emph{乐观外}松弛。\emph{无}额定、能量、无功能力与控制通道可用性的构网标志\emph{不}足以证明孤岛可维持。
储能续航按时序 SOC
\begin{equation}
 E_{b,t+1}=E_{b,t}+\eta_b^{ch}p_{b,t}^{ch}\Delta t-\frac{p_{b,t}^{dis}}{\eta_b^{dis}}\Delta t,\quad
 \underline E_b\le E_{b,t}\le\overline E_b,
\end{equation}
并以 $u_{b,t}\in\{0,1\}$ 禁止同时充放。\textbf{诚实}：储能\emph{不得}在每个恢复阶段或每次重叠停运独立按满
功率计入——事件阶段能量记账已在三阶段实现，多事件时序需序贯仿真。可再生受天气上限
$0\le p_{r,t}\le x_{r,t}^p\bar p_{r,t}^{weather}$；通信/控制丢失应按\emph{实际回退}模式（定值保持 /
本地下垂 / 失效跳闸）取值，而非一律当强迫停运——三种回退的可靠性价值迥异。合环前须满足同步窗
$|\Delta V|\le\Delta V^{max}$、$|\Delta f|\le\Delta f^{max}$、$|\Delta\theta|\le\Delta\theta^{max}$；忽略同步的
纯拓扑恢复计划须标注为可达服务的\emph{上界}。物理$\to$信息回耦（Level~3）中信息节点后备能量按
$B_{u,t+1}=\min(\overline B_u,\max(0,B_{u,t}+P_u^{ch}\Delta t-P_u^{load}\Delta t))$ 更新，可用度
$x_{u,t}^c=x_{u,t}^{intr}\mathbf 1[y_{\pi(u),t}{=}1\lor B_{u,t}{>}0]$，剩余续航 $T_u^{remain}=B_{u,t}/P_u^{load}$；
若 $T_u^{remain}<\Delta t$ 则节点在区间内失效，须事件队列或分数时长记账，端点二值更新会误判功能可用度。

\subsection{保真度阶梯与所需有效性标志}\label{sec:rel-pf-fidelity}
建议分级集成：$L0$——静态事件抽象（\S\ref{sec:rel-fa-cyber}），保护/FRT 为\emph{输入}概率与时长；
$L1$——对代表拓扑的给定故障/保护/FRT 轨迹生成互斥类，按场景权重折算 EENS；$L2$（当前离散事件子集已实现）——
在给定量测/DER 轨迹上联立多种继电判据、主后备/断路器和显式信息状态并自动生成三窗口后果；在线 $L2$ 由
Mass-Matrix DAE 生成网络/FRT/限幅轨迹；$L3$ 由连续信息物理事件队列执行年度抽样。类生成器
$\mathcal M_{PFRT}:(k,\omega)\mapsto(C_P,C_{FRT},C_{rec},T^{clear},\Phi_{k,c},\dots,\mathcal V^{PFRT})$
必须随附有效性向量 $\mathcal V^{PFRT}$（保真度、收敛、参数溯源、启动/FRT/限幅/同步边界附近的代理误差）。
严格地，最小混杂动态状态是 $(q,x)$，网络代数量 $y$ 仅在 $g_q(x,y)=0$ 的 index-1 正则分支上作为唯一
lifting；在线轨迹是在 hybrid time domain 上由模式内 Mass-Matrix DAE flow、可测 guard、首次到达时刻、
reset 和事件后代数重初始化构成的混杂执行。代数 fiber 为空/不唯一、$\partial g/\partial y$ 奇异或重初始化
残差不合格时必须进入 unresolved execution。阶段元组只是其可靠性奖励函数。完整执行 $\rho$ 是潜在生成对象，
类生成器只能消费有限记录 $\mathfrak o=\mathsf O_k(\rho)$；该记录包含已声明的量测/状态报告、事件词、监督命令和
计量或模型认证的服务边界，缺失字段保留 unknown，不得从隐藏 DAE 坐标补齐。有限类应定义为这些有限记录函数
在执行空间中的原像，并以 unresolved 原子闭合；$\mathsf O_k$ 非单射时，不同真实执行应留在同一类内。单个状态没有集合补集；
同层 class 的补集是其余同层类的不交并，子类在父类内的相对补集是 sibling classes 的不交并。该商空间只
保证同层互斥完备和跨层唯一父类，不宣称强 Markov lumpability 或连续轨迹双模拟；任何线性奖励的无偏压缩
都必须保存类条件矩，不能以一条名义代表轨迹替代。
这里的 $(q,x)$ 仅在固定完整场景后是最小动态状态，跨场景不自动构成 Markov state；若需要 Markov kernel，
必须加入足够的外生过程状态。当前求解器审计非线性收敛、耦合 Jacobian 分解和事件后代数残差，但不输出
$\inf_t\sigma_{\min}(\partial g/\partial y)>0$ 的全轨迹证书，因此 index-1 正则性是已选分支上的局部分析假设。
在线年度耦合器 \srcpath{compare\_online\_protection\_cyber\_reliability} 对同一故障分别运行主、后备
Mass-Matrix DAE 发现过程，保留 CT/PT 后的实测继电轨迹与 DER 逐点轨迹，再构造
$\mathcal M_{PFRT}$ 的输入并调用三级事件树。设 DAE 输出主保护命令时刻 $t_p^{cmd}$，断路器链总延时为
$d_p^{br}$，则事件树使用的主清除时刻严格取
\begin{equation}
 t_p^{clear}=t_p^{cmd}+d_p^{br};
 \end{equation}
不允许由 HTTP 或 GUI 再提供一个互相矛盾的清除时间。耦合成功时
\fld{online\_network\_dae\_coupled=true}，任一 DAE 不收敛、动作未施加或轨迹维数不一致均明确失败，
不得回退到用户给定轨迹。
\begin{boundarynote}
除非导出指标\emph{确实}消费上述耦合事件类或轨迹，任何单独动态研究都\emph{不}使可靠性结果 FRT 感知。合法置
\fld{protection\_frt\_reliability\_coupled=true} 需可追溯的类生成器；相关标志（设计文档 \S18.16）还包括
\fld{der\_ride\_through\_modelled}、\fld{protection\_coordination\_modelled}、\fld{breaker\_failure\_modelled}
等。只有 \fld{trace\_classes\_validated=true} 且类概率归一时聚合器才接受场景，并始终声明
\fld{online\_dae\_coupled=false}；在线年度耦合接口成功时单独置真，通用给定轨迹聚合器不冒充在线结果。
\end{boundarynote}

\subsection{case33bw AC/DC 恢复入口状态能力审计}
\label{sec:rel-pf-case33-propagation}
恢复模型中的二值换流器口径只读取
$a_c^{bin}=\mathbf 1[c\text{ authored in service}]$；DAE 入口口径则在故障
清除后读取端口在役状态与保护状态，
$a_c^{dae}=\mathbf 1[s_c^{port}=1\land s_c^{trip}=0]$。二者的差异数为
\begin{equation}
 N_{\Delta a}=\sum_c\left|a_c^{bin}-a_c^{dae}\right|.
\end{equation}
该比较遵循经典两状态可靠性元件模型与动态保护后果的分层定义，而不把额定容量
直接当成故障后可传输容量。可复现实现在
\srcpath{tools/case33bw\_acdc\_propagation\_study.cpp}，固定使用修正
\texttt{case33bw\_acdc} 的无故障、AC 母线 18 故障、传统 DC 母线 2 故障
以及启用直流故障控制的 DC 母线 2 故障轨迹。对故障量 $z_f$ 与相同初值的
无故障量 $z_0$，传播判据为
\begin{equation}
 \Delta z=\max_{t\in[t_f,T]}|z_f(t)-z_0(t)|,
 \qquad \Delta V>10^{-4}\ {\rm pu}\quad\text{or}\quad
 \Delta P_c>10^{-3}\ {\rm MW}.
 \label{eq:case33-propagation-gate}
\end{equation}
式~\eqref{eq:case33-propagation-gate} 是本研究预先声明的数值分离门，不是保护
定值；电压与功率门在 10 MVA 基准上均为 $10^{-4}$ pu，并远高于 $10^{-9}$
代数平衡容差。2 ms 结果的初始联合潮流残差为 $5.362\times10^{-11}$，
事件后最大代数残差为 $4.62\times10^{-14}$；无故障 AC/DC 电压漂移在输出
精度内为零。

AC 故障时，静态二值模型仍给出 VSC2 可用，而 DAE 在 0.365577673 s 记录
IEEE~1547 欠压跳闸并在恢复入口给出不可用，故 $N_{\Delta a}=1$。这证明
“二值在役”与“动态可恢复入口”可以不同。传统 GFL 路径下，DC 故障只产生
$2.1082\times10^{-10}$ pu 的 AC 电压偏差和 $1.2604\times10^{-10}$ MW
的换流器交流功率偏差，证明原路径实际上是单向耦合。启用
分段直流欠压降额/闭锁模型后，相同 DC 故障的
两项偏差分别增至 $8.807411\times10^{-2}$ pu 和 $1.05$ MW；1/2/5 ms
闭锁时刻为 $0.238/0.238/0.240$ s，传播分类一致，时刻离散误差范围为
2 ms。独立 Julia oracle 不调用 C++ 控制律，重新由 CSV 计算分段包络、20 ms
定时器、传统路径恒等性和传播门，全部通过。因此准入结论是：平衡相量
可靠性恢复入口具备 AC$\leftrightarrow$DC 传播，但不具备 EMT 直流故障电流
双向传播。

\begin{center}\footnotesize
\begin{tabular}{@{}p{3.0cm}p{4.2cm}p{3.3cm}p{3.2cm}@{}}
\toprule
\textbf{近五年 IEEE 证据} & \textbf{DC 故障至 AC 侧机制/验证} &
\textbf{与本文一致处} & \textbf{本文未覆盖处}\\
\midrule
Kaler--Yazdani 2022~\cite{relpf:kaler2022} & dq 动态模型、仿真及 1.25 kW
实验；故障闭锁拓扑限制电网侧电流 & 直流故障改变交流有功/电流，闭锁型 VSC
可用平均值控制描述 & 实验硬件拓扑与开关细节\\
Bandaru 等 2024~\cite{relpf:bandaru2024} & PSCAD/EMTDC；DC 故障时有功降至
零且可保留无功电压支撑 & 本文保留有功中断及其交流电压后果 & MMC 桥臂、
子模块与专用 FRT 策略\\
Yu 等 2024~\cite{relpf:yu2024}；Lv 等 2025~\cite{relpf:lv2025} & 控制/保护
分阶段动作；饱和改变故障响应 & 采用可审计的降额后定时闭锁顺序 & 桥臂电压
重构、饱和及阀级模型\\
Pandey 等 2025~\cite{relpf:pandey2025} & RTDS/CHIL；传统 BIC 二极管可形成
不控整流馈流，DC 故障显著改变 AC 端电流 & 证明 DC$\rightarrow$AC 物理因果
必须显式存在 & 二极管续流正是本文不支持的 EMT 路径\\
\bottomrule
\end{tabular}
\end{center}
表中论文验证的是机制而非本算例数值 oracle；不同拓扑、额定值和故障阻抗下的
波形幅值不可直接作为 case33 误差真值。本文阈值 $0.95/0.75$ pu 与 20 ms
延时均为研究参数，不是由文献反演或标准指定。模型仍不覆盖反并联二极管馈流、
MMC 子模块/桥臂饱和、开关纹波、DC 断路器电弧、直流过流保护及厂家解锁逻辑。

\subsection{N-k DAE 恢复入口类与既有切负荷 MILP}
\label{sec:rel-pf-nk-entry-class}
研究驱动 \srcpath{tools/nk\_acdc\_reliability\_study.cpp} 不把恢复优化结果
混入 trajectory class。对固定同时 N-k 进入事件 $e$，Mass-Matrix DAE 轨迹
$\rho_{e,\omega}$ 向恢复模型投影为附加不可用 VSC 集
\begin{equation}
 h_e(\rho_{e,\omega})=\mathcal U^{dae}_{vsc}(e,\omega),\qquad
 \rho_{e,\omega}\sim_e\rho_{e,\omega'}
 \Longleftrightarrow h_e(\rho_{e,\omega})=h_e(\rho_{e,\omega'}).
 \label{eq:rel-nk-entry-equivalence}
\end{equation}
运行状态 $\omega$ 的 AC/DC 负荷尺度仍进入每次恢复 MILP。原模型中的
$s_i^P,s_i^Q\ge0$ 分别进入有功/无功节点平衡，并由 $\lambda_{shed}$ 惩罚；
因此 class 只复用式~\eqref{eq:rel-nk-entry-equivalence} 的 DAE 入口，不删除切负荷
变量，不删除逐状态 MILP，也不使用类平均切负荷量替代状态后果。

该轨迹不是抽象状态序列。连续状态由每台换流器配置的 GFL（PLL、功率滤波、
电流响应/限幅及可选直流电容）或 GFM（内部相角/频率、电压/无功控制、功率滤波、
电流限幅及可选直流电容）状态组成；离散状态为开关位置、FRT、闭锁与保护状态。
交流侧由设备注入与 $Y_{ac}V_{ac}$ 的复电流平衡构成，直流侧由设备注入与
$G_{dc}V_{dc}$ 的电流平衡构成，VSC 交流功率按效率和直流电容能量平衡写入直流侧。
只有显式配置直流电压反馈、欠压降额或闭锁时，直流扰动才反向改变 GFL 交流功率/
电流指令；不能仅凭直流电容方程宣称 DC$\rightarrow$AC 传播。研究算例使用
MassMatrixDae 后向欧拉，并在每步内对控制状态、交流电压实虚部与直流电压做联立
阻尼 Newton 求解；状态事件经回退/二分定位、离散复位和事件后代数重初始化。
系统级恢复 MILP 不进入该 DAE，仅在 $T_{DAE}$ 后执行一次。

若后果相关 VSC 按 $(v_1,\ldots,v_{n_v})$ 排列，并令
$b^{(\ell)}=(b_{v_1},\ldots,b_{v_\ell})$，则可从包含全部已解析运行状态的单一集合开始，
逐个加入 VSC 可用状态。第 $\ell$ 层只能保持或细分第 $\ell-1$ 层，最终层等价于
式~\eqref{eq:rel-nk-entry-equivalence} 的完整不可用 VSC 集划分。层数因而等于
$n_v$，而不是固定三层；VSC 排序只影响中间层，不影响最终类。对固定进入事件，非空类互斥且
完备，类数不超过 $2^{n_v-|\mathcal V\cap\Delta\Phi_e|}$。按各类重新排列
$f_ep_\omega E_{e,\omega}$ 的求和严格保持事件频率和 EENS；这一恒等式不构成邻近运行状态
类别预测正确性的证据，预测仍必须由留一验证和被留出状态上的恢复重求解审计。

预测验证对每个事件留出一个运行状态，在归一化
$(\alpha_{AC},\alpha_{DC})$ 平面内从其余状态选择最近邻，只继承其入口类，随后
在被留出状态上重新求解恢复：
\begin{equation}
 \widehat S(e,\omega)=R\!\left(e,
 h_e(\rho_{e,j^*(e,\omega)}),\omega\right).
 \label{eq:rel-nk-loo-restoration}
\end{equation}
由类条件均值重构总体 EENS 是全概率公式恒等式，不作为式~\eqref{eq:rel-nk-loo-restoration}
的预测证据。成本审计保留恢复项：
\begin{equation}
 T_{direct}=N(\bar t_{DAE}+\bar t_{MILP}),\qquad
 T_{class}=T_{library}+N(\bar t_{lookup}+\bar t_{MILP}).
 \label{eq:rel-nk-cost}
\end{equation}
case33 的 92 个 N-1/N-2/N-3 联合进入事件与 6 个状态形成 552 个条件，全部求解；
92 个入口类给出 $1/6$ 压缩比，留一入口类准确率 100\%，恢复 EENS 误差为 0。
静态/DAE 入口 EENS 为 0.0905562/0.260831 MWh/yr；100000 样本的分项时间投影为
5.76--5.99 倍端到端加速，净节省成立但预注册 20 倍门槛失败。旧约 150 倍口径漏计了
每个状态的恢复 MILP，不能使用。

高阶故障组合和入口类划分不能混为同一层压缩。若 $d_e$ 为同时进入停运向量，
$\vartheta_e$ 为修复阶段参数，则恢复接口的完整离散特征为
\begin{equation}
 z_{e,\omega}=(d_e,h_e(\rho_{e,\omega}),\vartheta_e).
 \label{eq:rel-nk-complete-interface}
\end{equation}
当前代码以 `(event_index, class_label)` 分组，正好保持
式~\eqref{eq:rel-nk-complete-interface}：相同不可用 VSC 集在不同初始事件下不会合并。
对固定事件，$h_e$ 的非空逆像构成唯一最粗无损划分；任何保持恢复入口的其他划分都只能进一步细分它。
若全部枚举，事件数仍按 $\sum_m\binom{n_u}{m}$ 增长，trajectory class 只减少后续样本中重复 DAE 的
入口状态数，不消除故障组合空间。由于限流、孤岛和延时保护可导致故障阶数与跳闸集合非单调，不能未经
验证把父事件类传给子事件。当前驱动对 authored N-1/N-2/N-3 目录完全枚举；更高阶只能由具有频率的
联合事件模型/采样器给出，或在未展开 EENS 上界
$B_R=\sum_{(e,\omega)\in R}f_ep_\omega\bar E_{e,\omega}$ 小于给定容差后带证书停止。
该停止上界是理论扩展，当前算例没有把它宣称为已验证数值能力。

分域树模式曾因统一 AC+DC+VSC 连通图误剪 AC tie 而产生 150 个假不可行条件；
修正后统一图预剪仅用于统一树，分域模式由各域森林等式排环。该问题不是缺少
负荷切除能力。回归测试固定 VSC--DC--VSC 旁路、两条同时 AC 故障和两条必需
AC tie。当前频率为逐阶显式联合进入测度；顺序重叠、共因校准、EMT、原生三相
DAE 与恢复后非线性 AC/DC 认证均明确不覆盖。

case123 的 175 个联合进入事件与 6 个状态形成 1050 个条件；当前仅 892 个条件
进入并完成恢复 MILP，158 个在 DAE 或事件时代数网络求解阶段失败，未解析频率为
0.0319160/yr。已解析子测度上的静态/DAE 入口 EENS 为
0.0578627/0.235781 MWh/yr；已覆盖邻域入口类准确率为 100\%，但仍有
0.000761066/yr 的留一状态无已解析邻居。故准确性不准入。两次正式复跑的 8.81--8.91 倍成本投影也只作
诊断：失败长尾计入类库，而直接侧均值只来自已解析条件。未解析质量非零时
\fld{timing\_projection\_admissible=false} 且
\fld{net\_computational\_saving\_demonstrated=false}。这些失败属于 DAE 数值/模型
闭环，而非恢复 MILP 缺少负荷切除。

\subsection{与实现的对应}\label{sec:rel-pf-impl}
\begin{center}\footnotesize
\begin{longtable}{@{}L{5.4cm}L{9.0cm}@{}}
\toprule
\textbf{理论对象} & \textbf{平台实现状态}\\
\midrule
\endfirsthead
\toprule
\textbf{理论对象} & \textbf{平台实现状态}\\
\midrule
\endhead
\bottomrule
\endfoot
静态保护事件抽象（拒动/误动/拒开合/区扩展）& \textbf{已实现}：故障模式 FMEA 与三阶段消费为输入概率与配置时长；\srcpath{src/reliability/failure\_mode.cpp:build\_consequence\_patch}\\
短路换流器贡献（GFL 电流源、GFM 阻抗后电压源）& \textbf{已实现}（短路模块，静态快照）；\srcpath{src/short\_circuit/}\\
DER IEEE~1547 越限/跳闸/重连状态机 & \textbf{已实现}（动态，选项），由 L1 轨迹分类器复用；\srcpath{src/dynamics/}\\
逐带 FRT 计时器 / 模式自动机 $\Psi_r$ / 瞬时闭锁 & \textbf{已实现}：IEEE~1547 逐带计时与 \srcpath{DERMomentaryCessationSettings}\\
定/反时限、方向闭锁、距离/差动、清除链 / 主后备协调 / 断路器失灵 & \textbf{给定轨迹 L2 已实现}；\srcpath{src/reliability/protection\_frt.cpp:evaluate\_protection\_relay / evaluate\_protection\_coordination}\\
共享信息路径、QoS、共因环境、供电与电池自治、功能绑定 & \textbf{离散状态 L2 已实现（最多 20 个二值信息元件）}；\srcpath{src/reliability/protection\_frt.cpp:enumerate\_protection\_cyber\_states}\\
静态 FMEA / 仅保护 / 信息物理联合指标对比 & \textbf{已实现}；\srcpath{src/reliability/protection\_frt.cpp:compare\_protection\_cyber\_reliability}\\
保护--FRT 后果类 $\pi_c(k)$ 与 EENS/LOLE/LOLF & \textbf{L1 已实现}（轨迹验证类 + 外部切负荷阶段）；\srcpath{src/reliability/protection\_frt.cpp:aggregate\_protection\_frt\_reliability}\\
在线 DAE、瞬时闭锁、自动后果拓扑与重合/同步反馈 & \textbf{已实现}：\srcpath{online\_protection.cpp}、\srcpath{simulate\_recloser\_fuse\_sectionalizer}、\srcpath{evaluate\_automatic\_protection\_topology}、\srcpath{evaluate\_microgrid\_synchronization}\\
\end{longtable}
\end{center}

\subsection{参考文献}
\begin{thebibliography}{99}
\bibitem{relpf:ieee1547} IEEE Std 1547-2018, \emph{IEEE Standard for Interconnection and Interoperability
 of Distributed Energy Resources with Associated Electric Power Systems Interfaces}, 2018.
\bibitem{relpf:iec60255} IEC 60255-151, \emph{Measuring relays and protection equipment -- Part 151:
 Functional requirements for over/under current protection}, 2009.
\bibitem{relpf:c37} IEEE Std C37.113, \emph{IEEE Guide for Protective Relay Applications to Transmission Lines}.
\bibitem{relpf:kundur} P. Kundur, \emph{Power System Stability and Control}. New York, NY, USA:
 McGraw--Hill, 1994.
\bibitem{relpf:yazdani} A. Yazdani and R. Iravani, \emph{Voltage-Sourced Converters in Power Systems}.
 Hoboken, NJ, USA: Wiley--IEEE Press, 2010.
\bibitem{relpf:kaler2022} S. Kaler and A. Yazdani, ``A DC-side fault-tolerant
 bidirectional AC--DC converter for applications in distribution systems,''
 \emph{IEEE Access}, vol. 10, pp. 52625--52637, 2022,
 doi: 10.1109/ACCESS.2022.3171346.
\bibitem{relpf:bandaru2024} S. Bandaru \emph{et al.}, ``AC and DC fault
 ride-through strategies for hybrid MMC based HVDC link and interaction with
 multimachine power system dynamics,'' \emph{IEEE Trans. Ind. Appl.}, 2024,
 doi: 10.1109/TIA.2024.3396790.
\bibitem{relpf:yu2024} X. Yu \emph{et al.}, ``An active DC fault current
 limiting control for half-bridge MMC based on arm voltage reconstruction,''
 \emph{IEEE Trans. Power Del.}, 2024, doi: 10.1109/TPWRD.2023.3294173.
\bibitem{relpf:lv2025} Z. Lv \emph{et al.}, ``Modeling of the DC fault response
 of MMCs with fault control capability considering converter saturation,''
 \emph{IEEE Trans. Power Del.}, 2025, doi: 10.1109/TPWRD.2024.3485910.
\bibitem{relpf:pandey2025} R. Pandey \emph{et al.}, ``Time-domain backup
 protection for converter dominated hybrid AC--DC microgrids,''
 \emph{IEEE Trans. Ind. Informat.}, 2025, doi: 10.1109/TII.2025.3574424.
\bibitem{relpf:designdoc} HySim 团队，\emph{可靠性数学模型与智能信息物理评估}，\S18（设计文档）。
\end{thebibliography}

\section{保护测量、动作逻辑、配合与拓扑后果}\label{sec:rel-protection-chain}

\subsection{从一次量到可靠性后果的因果链}
保护可靠性不能只给一个“成功概率”。平台执行的链条是
\begin{equation}
\begin{aligned}
 (V_p,I_p)&\longrightarrow \text{CT/PT}\longrightarrow \text{继电判据}
 \longrightarrow \text{动作时刻}\\
 &\longrightarrow \text{断路器清除}\longrightarrow \text{拓扑后果}
 \longrightarrow \text{EENS/LOLE/LOLF}.
\end{aligned}
\end{equation}
每一箭头都有独立输入、失败模式和时间尺度。主保护、后备保护和信息通道必须在同一轨迹时基上计算，不能把
静态故障电流、平均继电延时和另一算例的断路器时间拼成不可追溯的清除时间。

\subsection{CT/PT 一阶动态与变比}
设一次相量为 $I_p,V_p$，CT/PT 变比分别为 $n_I,n_V$，理想二次目标为
$I_p/n_I,V_p/n_V$。一阶测量模型采用精确离散更新
\begin{align}
 I_s(t_{k+1})&=I_s(t_k)+\alpha_I\left(\frac{I_p(t_{k+1})}{n_I}-I_s(t_k)\right),\\
 V_s(t_{k+1})&=V_s(t_k)+\alpha_V\left(\frac{V_p(t_{k+1})}{n_V}-V_s(t_k)\right),\\
 \alpha_I&=1-e^{-\Delta t/\tau_I},\qquad
 \alpha_V=1-e^{-\Delta t/\tau_V}.
\end{align}
$\tau=0$ 表示无滞后的代数测量。更新后再施加二次幅值饱和，保留相角而限制模值。这样 CT 饱和会降低过流元件
看到的电流并改变视在阻抗，PT 饱和则会改变距离元件的电压量。无效变比、负时间常数和非递增时间戳在继电评估前拒绝。

\subsection{定时限过流}
定时限元件在连续满足 $|I|\ge I_{pickup}$ 且方向条件成立时积累计时。若启动在 $t_a$，动作时刻为
\begin{equation}
 t_{cmd}=\inf\{t\ge t_a:t-t_a\ge T_d\}.
\end{equation}
电流跌落到启动值以下会按复位时间衰减，而不是永久保留先前启动时间。该语义能区分“两个各持续半个延时的脉冲”
与“一个连续完整脉冲”，前者不应误动作。

\subsection{反时限过流}
对电流倍数 $M=|I|/I_{pickup}>1$，IEC 形式动作时间为
\begin{equation}
 T(M)=\mathrm{TMS}\left(\frac{A}{M^p-1}+B\right)+T_{add}.
\end{equation}
轨迹电流变化时不使用某一峰值代替全程，而积累无量纲动作量
\begin{equation}
 \Omega_{k+1}=\Omega_k+\frac{\Delta t}{T(M_k)}.
\end{equation}
当 $M\le1$ 时按 $\dot\Omega=-\Omega/T_{reset}$ 复位；$\Omega\ge1$ 的首次插值时刻为命令时刻。
这一积分使 DER 电流限幅、GFM/GFL 模式切换和故障电压恢复能够真实影响保护到达时间。

\subsection{方向闭锁}
方向元件要求 \fld{directional\_current}>0。平台把方向量作为每个采样点的显式输入，不从电流幅值猜测方向。
在线 DAE 入口目前把故障支路电流按约定正方向写入方向量；离线轨迹调用者必须自行保证方向参考与继电安装方向一致。
若方向信号缺失，应禁用方向元件或提供经过极化算法计算的方向量，不能默认“幅值大即正向”。

\subsection{复数距离保护}
有效相量点计算
\begin{equation}
 Z_{app}=\frac{V_s}{I_s}.
\end{equation}
当 $|I_s|$ 太小时视为无效点。圆形幅值区判据为 $|Z_{app}|\le Z_{reach}$。Mho 区以线路角
$\theta_l$ 定向，圆心 $C=\tfrac12 Z_{reach}e^{j\theta_l}$、半径 $R=Z_{reach}/2$：
\begin{equation}
 |Z_{app}-C|\le R.
\end{equation}
该圆经过原点，天然提供方向性。测试用正向 $5\ \Omega$ 和反向 $-1\ \Omega$ 验证边界，防止把 mho 退化成
只比较模值的无方向圆。

\subsection{四边形距离区}
四边形区在旋转到线路坐标后，以前向/反向电阻和电抗界定义：
\begin{equation}
 -R_{rev}\le \Re(Ze^{-j\theta_l})\le R_{fwd},\qquad
 -X_{rev}\le \Im(Ze^{-j\theta_l})\le X_{fwd}.
\end{equation}
四个界必须为非负有限量。它可独立放宽故障电阻覆盖而不扩大反向范围，适用于高阻接地故障筛选。平台逐区按延时评估，
返回首个实际动作区的序号和名称；区序列是稳定配置顺序，不用到达时间排序后改写身份。

\subsection{差动保护}
差动判据使用差流 $I_d$ 和制动流 $I_r$：
\begin{equation}
 I_d\ge I_{pickup}+kI_r\quad\lor\quad I_d\ge I_{high}.
\end{equation}
高定值段绕过比例制动。测试分别构造高制动不动作和高定值动作点。差动量必须来自同一时间基准的多端测量；
若通信对时不满足要求，应通过信息功能门控使该保护通道不可用，而不是继续使用失同步相量。

\subsection{自适应整定}
自适应上下文给出启动值、距离范围和延时的正比例缩放：
\begin{align}
 I'_{pickup}&=k_I I_{pickup},\\
 Z'_{reach}&=k_Z Z_{reach},\\
 T'_d&=k_TT_d.
\end{align}
只有 \fld{communication\_available=true} 时应用新整定；通信不可用保留基础定值。这对应“失联保持最后安全本地定值”策略，
不是普遍适用的自动回退。若实际设备策略是闭锁或跳闸，应通过后果类和功能绑定显式建模。

\subsection{断路器清除链}
继电器出令并不等于电弧熄灭。清除时刻为
\begin{equation}
 t_{clear}=t_{cmd}+t_{channel}+t_{coil}+t_{mechanical}+t_{arc}.
\end{equation}
四段延时必须为有限非负数。通道延时只用于确实依赖通信的方案，本地硬接线应取零。机械拒动通过按需成功概率形成
主断路器失败分支，不能同时又把同一统计数据并入“失灵保护概率”，否则会重复计算。

\subsection{主后备配合}
主、后备分别计算清除时刻 $t_p,t_b$。要求配合裕度 $\Delta T_{req}$ 时，选择性条件为
\begin{equation}
 t_b-t_p\ge\Delta T_{req}.
\end{equation}
平台把“主保护已清除但裕度不足”保留为已清除事件并附选择性诊断，不把它误分类成拒动。主保护未动作时，后备仍可独立
清除；后备未动作时清除时刻无定义，不能用零参与裕度相减。

\subsection{保护事件树概率}
设主继电、主断路器、后备链成功概率为 $p_R,p_B,p_{bak}$。互斥分支为
\begin{align}
 p_{primary}&=p_Rp_B,\\
 p_{backup,R}&=(1-p_R)p_{bak},\\
 p_{backup,B}&=p_R(1-p_B)p_{bak},\\
 p_{unclear,R}&=(1-p_R)(1-p_{bak}),\\
 p_{unclear,B}&=p_R(1-p_B)(1-p_{bak}).
\end{align}
五项严格和为一。检测或跳闸信息功能不可用时，相应继电成功率被置零，并以明确失效原因进入后备或未清除类。
这不是把通信不可用率简单加到设备故障率，而是在同一互斥事件树中条件化。

\subsection{重合器、熔断器与分段器自动机}
重合序列输入每次带故障通流时间、死区时间和最大次数。第 $j$ 次跳闸后，临时故障可清除并在死区后重合；永久故障
继续下一次。熔断器热积累以总熔断时间归一：
\begin{equation}
 h_f\leftarrow h_f+\frac{\Delta t_{energized}}{T_{fuse}},qquad h_f\ge1\Rightarrow\text{熔断}.
\end{equation}
分段器在上游无电流死区达到计数门限时开断。若设瞬时闭锁，则第一次跳闸后直接闭锁，不再重合。输出事件保留时刻、
动作类型和次数，并给出熔断、分段开断、重合器闭锁和最终故障清除状态。

\subsection{序列与供电可靠性指标}
瞬时跳闸后成功重合可能只构成瞬时停电，不一定进入 IEEE~1366 持续停电指标；是否计入 MAIFI 取决于研究阈值。
平台序列模拟器只给物理时序，不擅自把所有跳闸计为 SAIFI。年度聚合必须由后果窗口明确给出持续时长和切负荷阈值，
从而避免把毫秒级动作按一次长期修复事件计入 EENS。

\subsection{自动拓扑后果}
\srcpath{evaluate\_automatic\_protection\_topology} 接收稳定开关、交流支路和直流支路 ID，复制系统后执行开断，
再从交流平衡源、可成岛构网源、直流源和可用换流通道作域限定可达性搜索。不可达母线上的在役负荷形成切负荷：
\begin{equation}
 S=\sum_{i\in\mathcal B_{deenergized}^{ac}}P_i^d+
 \sum_{d\in\mathcal B_{deenergized}^{dc}}P_d^d.
\end{equation}
输出分别给出交流和直流母线稳定 ID，不能因两域母线号相同而合并。未知稳定 ID 直接拒绝，防止把“未找到设备”解释为
“设备已经打开”。

\subsection{拓扑后果的边界}
该后果是源可达性筛选，不证明电压、无功、热稳定或暂态稳定可行。结果明确声明
\fld{model\_scope=automatic-protection-topology-source-reachability} 与中文限制。若用于精细 EENS，需把开断后的系统送入
混合 AC/DC 后果 LP 或潮流/OPF 校核；可达性结果一般是可服务负荷的乐观上界。

\subsection{无故障误动}
在无故障判别窗中，误报、通道成功和断路器成功共同导致真实停电。频率与指标为
\begin{align}
 \lambda_{NT}&=\nu_{win}p_{FA}p_{ch}p_{br},\\
 \mathrm{EENS}_{NT}&=\lambda_{NT}S_{trip}\tau_{rest},\\
 \mathrm{LOLE}_{NT}&=\lambda_{NT}\tau_{rest}\mathbf1[S_{trip}>\varepsilon],\\
 \mathrm{LOLF}_{NT}&=\lambda_{NT}\mathbf1[S_{trip}>\varepsilon].
\end{align}
该链由两个入口消费：独立函数
\srcpath{evaluate\_protection\_misoperation}，以及通用 FMEA 的
\srcpath{ProtectionMisoperationOptions}。后者把增量加入系统总量并返回分解残差，是可执行计算字段。

\subsection{保护配合对结果的方向性影响}
延长主保护延时通常增加故障清除窗口 EENS，并可能缩小后备裕度；提高主保护成功率通常把概率质量从后备/未清除类移向
主清除类。但若主保护误切更大健康区，简单提高成功率未必降低总风险。因此平台把动作概率与类条件后果分开，只有在
$E_{primary}\le E_{backup}\le E_{unclear}$ 的单调后果假设下，成功率改善才保证 EENS 不增。

\subsection{保护验证矩阵}
注册测试覆盖：CT/PT 一阶闭式响应与饱和、mho 正反向、四边形边界、差动制动与高定值、通信失联保持基础整定、
主后备裕度合格/不合格、临时故障重合、永久故障分段、熔断竞争、瞬时闭锁、AC/DC 域限定拓扑、未知稳定 ID、无故障误动
指标。任何新增保护特性必须同时给出动作点、拒动作点和边界点，单个“能动作”算例不足以证明判据实现正确。

\section{保护拓扑、误动与配合的可复核推导}\label{sec:rel-protection-topology-cases}

本章把保护装置的动作结果从“成功率乘一个后果”展开为可执行的因果链。分析对象不是单个继电器的动作精度，而是从
量测、判据、出口、断路器、拓扑分离到恢复窗口的系统级可靠性贡献。所有概率均为无量纲，时间先以秒完成保护配合，进入
年度指标前再除以 $3600$ 转为小时；支路、母线和 DER 均使用对外稳定 ID，不以向量位置代替。

\subsection{故障事件与无故障误动事件必须分开}
设物理故障场景集合为 $\mathcal E_F$，无故障判别窗集合为 $\mathcal E_N$。物理故障场景 $e$ 的年频率为
$\lambda_e$，联合类为 $c$，类概率为 $\pi(c\mid e)$，三窗口后果为
\begin{equation}
 W(e,c)=\sum_{w=1}^{3}S_{e,c,w}\tau_{e,c,w}.
\end{equation}
故障触发的期望未供电量为
\begin{equation}
 \mathrm{EENS}_{F}=\sum_{e\in\mathcal E_F}\lambda_e
 \sum_c\pi(c\mid e)W(e,c).
\end{equation}
无故障时不存在上述起始故障频率。若一年有 $\nu_{win}$ 个独立判别窗，单窗误跳概率为 $p_{FA}$，跳闸通道与断路器
成功率分别为 $p_{ch}$、$p_{br}$，则真正造成开断的误动频率为
\begin{equation}\label{eq:rel-misop-frequency}
 \lambda_{NT}=\nu_{win}p_{FA}p_{ch}p_{br}.
\end{equation}
设误动切除负荷为 $S_{trip}$，恢复时间为 $\tau_{rest}$，则
\begin{align}
 \Delta\mathrm{EENS}_{NT}&=\lambda_{NT}S_{trip}\tau_{rest},\\
 \Delta\mathrm{LOLE}_{NT}&=\lambda_{NT}\tau_{rest}
   \mathbf 1[S_{trip}>\varepsilon],\\
 \Delta\mathrm{LOLF}_{NT}&=\lambda_{NT}
   \mathbf 1[S_{trip}>\varepsilon].
\end{align}
这里 $p_{ch}$ 和 $p_{br}$ 出现在成功链而不是失败链中，因为误判只有在跳闸命令被成功传输且断路器成功开断时才形成
停电。若把它们写成 $1-p$，会得到“通道越可靠，误动停电越少”的错误方向。

\subsection{误动算例与守恒检查}
取 $\nu_{win}=100$ 次/年、$p_{FA}=0.01$、$p_{ch}=0.5$、$p_{br}=0.8$，则
\begin{equation}
 \lambda_{NT}=100\times0.01\times0.5\times0.8=0.4\ \text{次/年}.
\end{equation}
若 $S_{trip}=2$ MW、$\tau_{rest}=0.25$ h，则
\begin{equation}
 \Delta\mathrm{EENS}_{NT}=0.4\times2\times0.25=0.2\ \text{MWh/年},
\end{equation}
同时 $\Delta\mathrm{LOLE}_{NT}=0.1$ h/年、$\Delta\mathrm{LOLF}_{NT}=0.4$ 次/年。程序不仅返回这四个量，
还校核
\begin{equation}\label{eq:rel-eens-decomposition}
 r_E=\mathrm{EENS}_{total}-\mathrm{EENS}_{physical}
 -\Delta E_{information}-\Delta E_{intelligent}
 -\Delta\mathrm{EENS}_{NT}.
\end{equation}
专项算例要求 $|r_E|\le10^{-12}$ MWh/年。这个残差可以发现两个常见错误：只显示误动字段却没有把它加入总量，或者既在
故障模式目录中计一次、又在无故障判别窗中重复计一次。

\subsection{误动参数的边际影响}
在其他参数固定且 $S_{trip}>\varepsilon$ 时，式~\eqref{eq:rel-misop-frequency} 给出
\begin{align}
 \frac{\partial\Delta\mathrm{EENS}_{NT}}{\partial p_{FA}}
 &=\nu_{win}p_{ch}p_{br}S_{trip}\tau_{rest},\\
 \frac{\partial\Delta\mathrm{EENS}_{NT}}{\partial p_{ch}}
 &=\nu_{win}p_{FA}p_{br}S_{trip}\tau_{rest},\\
 \frac{\partial\Delta\mathrm{EENS}_{NT}}{\partial\tau_{rest}}
 &=\lambda_{NT}S_{trip}.
\end{align}
因此误报率降低一半会使误动 EENS 精确降低一半；但通道可用率降低虽然也减少误动停电，却会增加真实故障时的拒动风险。
工程比较必须把两类事件同时纳入，不能凭误动一项得出“降低通信可靠性有益”的结论。

\subsection{量测链的离散更新}
保护判据消费 CT/PT 二次量而不是网络一次量。对时间常数 $T_m$、采样步长 $\Delta t$ 的一阶量测环节，平台采用精确
零阶保持离散式
\begin{equation}\label{eq:rel-instrument-exact-step}
 y_k=\alpha y_{k-1}+(1-\alpha)\frac{x_k}{n},\qquad
 \alpha=\exp(-\Delta t/T_m),
\end{equation}
其中 $n$ 为一次/二次变比。$T_m=0$ 时直接取 $x_k/n$。若二次侧饱和上限为 $y_{sat}$，则复相量按幅值投影
\begin{equation}
 \widehat y_k=
 \begin{cases}
 y_k,&|y_k|\le y_{sat},\\
 y_{sat}y_k/|y_k|,&|y_k|>y_{sat}.
 \end{cases}
\end{equation}
这一投影保留相角而限制幅值。CT 饱和会降低过流元件看到的电流，可能推迟或阻止动作；PT 动态会改变距离元件的
$Z_{app}=V/I$ 轨迹。程序把变比、时间常数、饱和上限同时送入继电器量测轨迹，测试以式
~\eqref{eq:rel-instrument-exact-step} 为闭式预言逐点比较。

\subsection{定时限与反时限的统一动作积分}
定时限过流在 $I(t)>I_p$ 时累计拾取时间，低于拾取值时按复归时间衰减。连续超过定值 $T_d$ 后，命令时刻为
\begin{equation}
 t_{cmd}=\inf\left\{t:\int_0^t\mathbf1[I(\xi)>I_p]\,d\xi\ge T_d\right\}.
\end{equation}
IEC 60255 型反时限曲线的局部动作时间写成
\begin{equation}
 T(I)=\mathrm{TMS}\left[\frac{A}{(I/I_p)^P-1}+B\right]+T_{add},
\end{equation}
继电器状态量按
\begin{equation}
 z_{k+1}=z_k+\frac{\Delta t_k}{T(I_k)}
\end{equation}
积分，首次达到 $z\ge1$ 时动作。该写法允许电流随在线 DAE 变化，而不是把故障初值代入曲线一次。若限流型 DER 使
$I/I_p$ 从 4 降到 1.2，$T(I)$ 会显著增加；这正是静态短路电流法与在线耦合结果产生差异的来源。

\subsection{方向、距离与差动判据}
方向元件先检查方向量为正，再允许过流或距离判据累计。距离保护的 mho 圆以线路角 $\theta_l$ 和整定阻抗 $Z_r$ 定义，
判据为
\begin{equation}
 \left|Z_{app}-\frac{Z_r}{2}e^{j\theta_l}\right|
 \le\frac{Z_r}{2}.
\end{equation}
四边形区则分别限制旋转坐标中的电阻和电抗上下界，适合独立设置耐受故障电阻的范围。差动元件使用
\begin{equation}
 I_{diff}\ge\max(I_{pickup},k_s I_{rest})
 \quad\text{或}\quad I_{diff}\ge I_{high}.
\end{equation}
这些判据的输出不是最终可靠性指标，而是主/后备动作时刻与动作布尔值。随后断路器链增加通道、线圈、机械和燃弧延时：
\begin{equation}\label{eq:rel-breaker-chain}
 t_{clear}=t_{cmd}+t_{ch}+t_{coil}+t_{mech}+t_{arc}.
\end{equation}
任何一项为负、非有限数或使清除时刻超过动态时域，入口都会拒绝，而不是截成零。

\subsection{主后备配合的含义}
对主保护 $i$ 和后备保护 $j$，配合裕度为
\begin{equation}
 M_{ij}=t_{clear,j}-t_{clear,i}-M_{req}.
\end{equation}
$M_{ij}\ge0$ 表示满足选择性。配合报告同时保留两个清除时刻、要求裕度和实际裕度。可靠性事件树仍按真实动作结果生成
主清除、后备清除和未清除类；单纯的裕度告警不会凭空扩大切负荷。若运维策略规定裕度不足时闭锁重合或扩大隔离区，调用方
应在后果模型中显式表达该策略。

设主继电与主断路器按需成功率分别为 $p_{r1},p_{b1}$，后备对应 $p_{r2},p_{b2}$，则忽略信息系统时
\begin{align}
 p_P&=p_{r1}p_{b1},\\
 p_B&=(1-p_P)p_{r2}p_{b2},\\
 p_U&=(1-p_P)(1-p_{r2}p_{b2}),
\end{align}
且 $p_P+p_B+p_U=1$。程序对每个信息状态重新计算这些互斥类，并在最终聚合前检查概率和。

\subsection{开断后的交流拓扑判定}
设在役交流图为 $G=(V,E)$，保护跳开的稳定支路为 $e_t$，电压源母线集合为 $V_s$。开断后可带电集合为
\begin{equation}
 V_E=\operatorname{Reach}_{G\setminus\{e_t\}}(V_s),
\end{equation}
失电母线集合为 $V_D=V\setminus V_E$。平台从外部电网、在役发电机、构网储能和构网 VSC 的交流端建立 $V_s$，沿除
$e_t$ 外的在役支路做确定性图遍历。每个 $v\in V_D$ 使用稳定母线 ID 进入结果。

在线闭环在 $t_{clear}$ 同时施加三类事件：清除故障并联支路、跳开被保护支路、把 $V_D$ 上恒功率负荷缩放为零。第三项
不是把停电后果删除；它只使无源孤岛的 DAE 代数方程符合“负荷已失电”的物理状态。年度 EENS 仍由三窗口后果计算，因而
不会因动态负荷退出而漏计能量。

\subsection{径向馈线手算}
考虑 $1-2-3$ 径向馈线，母线 1 有外部电网，母线 2、3 各有 1 MW 负荷。若保护跳开支路 $1-2$，则
$V_E=\{1\}$、$V_D=\{2,3\}$，动态闭环应在清除时刻退出 2 MW 恒功率负荷。若跳开 $2-3$，则
$V_E=\{1,2\}$、$V_D=\{3\}$，只退出 1 MW。专项测试构造单支路径向系统，要求失电母线稳定 ID 精确等于
\verb|{2}|、负荷退出事件确实被 DAE 应用且闭环收敛。

若在 $1-2$ 之间有两条并联支路，只跳开其中一条后 $V_D=\varnothing$。此时不生成负荷退出，另一条支路承担供电。这个
对照证明算法依据连通性而不是“任何跳闸都切除下游负荷”的固定模板。

\subsection{自动拓扑后果与区段切负荷}
给定故障支路、保护区和开关状态，自动后果入口分别生成主保护、后备保护与未清除状态。主保护只切除最小故障区；后备保护
可以按配置扩展到上游保护区；未清除状态保持故障影响至物理修复。每个输出含受影响负荷稳定 ID、总切负荷 MW 和连通性
诊断。AC 与 DC 母线映射保持域限定，数值相同的 AC bus 5 与 DC bus 5 不会混为同一节点。

当用户给出的保护区引用不存在、支路已退出或稳定 ID 重复时，入口明确拒绝。没有映射到负荷的合法空区返回零后果并保留
成功状态；“空区”和“引用错误”不能用同一个零值表示。

\subsection{重合器、熔断器与分段器状态序列}
自动序列按时间顺序处理重合器跳闸、死区、重合、熔断竞争和分段器无电流计数。设重合器第 $m$ 次曲线动作时刻为
$t_{R,m}$，熔断器熔断时刻为 $t_F$。在 fuse-saving 策略中，若 $t_{R,1}<t_F$，第一次快速跳闸先清除电流，瞬时故障
在死区后消失，熔断器保持；永久故障则在后续慢曲线中可能满足 $t_F<t_{R,m}$，由熔断器隔离支线。

分段器不直接开断故障电流。它在上游重合器造成的无电流死区中累计计数，达到门槛后开断，再允许上游重合器恢复健康区段。
程序输出每个事件的时间、动作类型、设备稳定 ID、计数和最终开合状态。若启用瞬时闭锁，第一次跳闸后序列终止，不会伪造
后续重合。

临时故障算例要求序列终态为重合器闭合、熔断器完好、分段器闭合；永久主干故障要求重合器最终闭锁；永久支线故障在设定
熔断竞争条件下要求熔断器开断而健康主干恢复。测试同时检查事件时间单调，防止“后发生的熔断被排在先发生的重合之前”。

\subsection{DER 穿越、瞬时闭锁与跳闸}
DER 状态机至少区分穿越、瞬时闭锁和物理跳闸。设进入闭锁的低压门槛为 $V_{in}$，退出门槛为 $V_{out}>V_{in}$，退出
驻留时间为 $T_{dwell}$，最大闭锁时间为 $T_{max}$。轨迹第一次满足 $V<V_{in}$ 时进入闭锁；之后连续满足
$V>V_{out}$ 达 $T_{dwell}$ 才恢复。若闭锁持续超过 $T_{max}$，终态升级为跳闸。

迟滞 $V_{out}-V_{in}$ 防止电压在边界附近抖动导致每个采样点切换状态。输出保留首次闭锁时刻、首次恢复时刻、终态和是否
超时。保护事件树分别消费主清除闭环轨迹、后备清除闭环轨迹和未清除发现轨迹；三类不能共享一条 DER 终态，因为不同清除
时刻可能跨越 $T_{max}$。

\subsection{微网同步检查}
重合闸或并列前需同时满足电压、频率和相角窗口。设两侧测量为 $(V_1,f_1,\theta_1)$ 与
$(V_2,f_2,\theta_2)$，平台计算
\begin{align}
 \Delta V&=|V_1-V_2|,\\
 \Delta f&=|f_1-f_2|,\\
 \Delta\theta&=\left|\operatorname{wrap}_{[-\pi,\pi)}(\theta_1-\theta_2)\right|.
\end{align}
只有三项均不超过门槛、同步测量可用且合闸命令通道可用时才允许合闸。相角按周期折回十分关键：$179^\circ$ 与
$-179^\circ$ 的差应为 $2^\circ$，不是 $358^\circ$。测量缺失或命令通道失效均返回禁止合闸及明确原因，不能把缺数当成
零偏差。

\subsection{保护因素对三种方法的影响方向}
静态 FMEA 不执行保护，故继电延时、断路器机械时间和配合裕度不会改变其结果。仅保护方法消费这些量：清除时间增加会扩大
第一窗口未供电量；主保护拒动会把概率质量从主清除移向后备或未清除。信息物理联合方法在仅保护基础上再受检测、跳闸、
隔离和恢复功能状态影响。

若后果满足 $E_P\le E_B\le E_U$，提高主保护成功率使仅保护与联合 EENS 不增；若主保护误切更大区域导致 $E_P>E_B$，
方向可能反转，程序不强制单调。保护误动只进入配置了无故障判别窗的 FMEA 总量，不应改变以单一物理故障构造的三级对照
静态行。报告必须把这种口径差异写清楚。

\subsection{保护链验收条件}
保护相关实现通过以下不变量验收：量测变比与饱和逐点满足闭式离散式；定/反时限动作积分不跨越采样时间；距离区边界和差动
斜率边界分别有内外点；主后备清除时刻与断路器链满足式~\eqref{eq:rel-breaker-chain}；重合序列事件时间单调；拓扑结果使用
稳定 ID；径向开断后失电岛负荷退出且年度后果仍计入；误动分解满足式~\eqref{eq:rel-eens-decomposition}；同步检查同时消费
三种差值与两项信息可用性。对应实现位于 \srcpath{src/reliability/protection\_frt.cpp}、
\srcpath{src/reliability/online\_protection.cpp} 与
\srcpath{src/reliability/reliability\_assessment.cpp}，注册测试以闭式数值和终态共同断言，而不是只检查函数不抛异常。

\section{信息系统拓扑、服务质量与序贯事件队列}\label{sec:rel-cyber-chronology}

\subsection{信息元件、路径和功能}
信息元件 $q$ 包含固有可用率 $A_q$、单包交付概率 $d_q$、时延 $l_q$、抖动 $j_q$、供电母线、备用能量、功耗和
共因组。功能 $f$ 由一组替代成功路径 $\mathcal P_f$ 描述。路径 $p$ 的确定性服务量为
\begin{equation}
 L_p=\sum_{q\in p}l_q,\qquad J_p=\sum_{q\in p}j_q,\qquad D_p=\prod_{q\in p}d_q.
\end{equation}
只有满足 $L_p\le L_f^{max}$、$J_p\le J_f^{max}$、$D_p\ge D_f^{min}$ 的路径参与功能可用性。零上限表示关闭对应
确定性门限，不表示所有路径时延为零。

\subsection{共享依赖不能按路径独立相乘}
若两条路径共享交换机、电源或控制中心，路径成功事件相关。平台在元件位空间上精确枚举，同一元件只生成一个二值位：
\begin{equation}
 A_f=\sum_{x^c}\mathbf1\!\left[\bigvee_{p\in\mathcal P_f}
 \bigwedge_{q\in p}x_q^c=1\right]p(x^c).
\end{equation}
因此共享元件失效会同时切断所有经过它的路径。离散事件树最多处理二十个二值信息元件，超过限制明确拒绝；结构可靠性
独立接口采用成功路径容斥，最多二十四条约简路径。两个限制针对不同指数维度，不能互相替代。

\subsection{最小割集}
成功路径族的最小击中集就是功能最小割集。平台逐路径扩展候选击中集，并删除包含已有小集合的超集。割集用于解释功能
脆弱点，不直接假定割集互斥。若需要由割集近似不可用度，必须处理割集交叠；当前结构接口直接计算成功路径并集概率，
避免一阶割集和在高不可用率下越过一。

\subsection{联合功能可用性}
FLISR 同时需要检测、隔离和恢复功能，联合成功事件为
\begin{equation}
 \Phi_{DIR}(x)=\phi_D(x)\land\phi_I(x)\land\phi_R(x).
\end{equation}
平台先从每个功能选一条成功路径，再合并所需元件，删除冗余联合路径并做精确容斥。该结果一般不等于
$A_DA_IA_R$。只有三功能依赖元件集合互不相交且条件独立时，标量乘积才成立。

\subsection{共因环境条件化}
天气、站点损毁或控制中心故障以互斥外层环境 $e$ 表示，概率 $p_e$ 和为一。每个环境指定强制失效的共因组和失电母线。
条件于环境后，剩余元件才按独立固有状态枚举：
\begin{equation}
 \Pr(x^c)=\sum_ep_e\Pr(x^c\mid e).
\end{equation}
这比把共因概率分别乘到每条路径更严格，因为同一环境同时改变多个元件。环境概率不归一、未知共因组或非法概率均在
枚举前拒绝。

\subsection{物理供电与备用电池}
信息元件绑定物理母线 $b(q)$。若环境使该母线失电，元件仅在备用自治时间覆盖信息停电时仍可用：
\begin{equation}
 T_q^{aut}=\frac{E_q^{backup}}{P_q^{draw}},\qquad
 x_q^{eff}=x_q^{intr}\mathbf1[T_q^{aut}\ge T^{outage}].
\end{equation}
$P_q^{draw}=0$ 且有备用能量可解释为本模型内不耗能；负功耗或负能量非法。静态离散环境比较持续时间与自治时间，序贯
事件队列则连续更新剩余能量，能表示信息节点在一次长停电中途失效。

\subsection{保护功能绑定}
事件输入把检测、主跳闸、后备跳闸、隔离和恢复分别绑定到信息功能索引。索引 $-1$ 表示本地独立功能，不需要通信；
非负索引必须在功能数组范围内。检测不可用会同时阻断自动隔离与依赖检测的主后备逻辑，恢复只有在隔离成功后才可能成功：
\begin{equation}
 I=D\land\phi_I,qquad R=I\land\phi_R.
\end{equation}
该因果关系避免“未检测故障但恢复成功”的非法类。人工恢复不是把 $R$ 强行置真，而是在三窗口后果中选用人工时长。

\subsection{静态信息状态事件树}
对每个信息状态，保护事件树再展开继电与断路器按需成功分支。联合类概率为
\begin{equation}
 \pi_c=p_e\,p(x^c\mid e)\,p(c_P\mid x^c,e).
\end{equation}
所有正概率类之和必须在 $10^{-12}$ 容差内为一。零概率类不输出，从而减少 JSON 体积；这不改变概率总量。每一类保留
环境名、信息元件状态、功能结果、保护清除类型、失效原因和 DER-FRT 终态。

\subsection{三窗口后果}
联合类的事件能量由清除、切换和修复三个连续窗口组成：
\begin{align}
 E_c={}&S_c^{clear}\frac{t_c^{term}}{3600}
 +S_c^{switch}\tau_c^{switch}
 +S_c^{residual}\max(0,\tau^{repair}-\tau_c^{switch}).
\end{align}
主清除、后备清除和未清除分别选择不同清除切负荷；隔离失败选择扩大区切负荷；恢复失败保留残余切负荷。DER 终态造成的
发电损失等值切负荷逐设备相加，并限制恢复比例到 $[0,1]$。

\subsection{序贯 CTMC 元件过程}
序贯信息物理模拟为每个信息元件建立故障/修复事件。若年故障率为 $\lambda_q$、平均修复时间为 $r_q$，等待时间分别为
\begin{equation}
 T_q^{up}\sim\operatorname{Exp}(\lambda_q/H),\qquad
 T_q^{down}\sim\operatorname{Exp}(1/r_q).
\end{equation}
事件队列使用连续小时，保留亚小时动作，不向上取整到整点。相同时间事件用单调序号稳定排序，保证固定随机种子下执行顺序
确定。

\subsection{共因事件过程}
共因组以发生率和平均持续时间生成开始/结束事件。开始时组内全部元件强制失效，结束后恢复到各自固有状态；若个别元件
本身仍处于故障修复期，则共因结束不会错误地把它置为可用。结果返回共因事件数和
\fld{common\_cause\_modelled}，便于区分“输入为空”和“模型未消费”。

\subsection{物理故障和供电恢复事件}
物理故障可以按泊松过程生成，也可提供确定发生时刻用于解析回归。故障使绑定母线失电并安排供电恢复事件；其间信息节点
从备用电池取能。供电恢复只解除物理失电约束，不会修复固有信息故障或共因故障。

\subsection{电池连续耗尽}
在事件间隔 $[t_k,t_{k+1})$ 内，失电信息节点能量按
\begin{equation}
 E_q(t_{k+1})=\max\{0,E_q(t_k)-P_q^{draw}(t_{k+1}-t_k)\}
\end{equation}
更新。若耗尽时刻落在区间内部，模拟器插入耗尽边界而不是等到下一外部事件。这保证保护功能在电池耗尽前后可以选择不同
分支，并返回全程最小电池能量。

\subsection{包交付的序贯抽样}
静态枚举把 $A_qd_q$ 合成条件成功概率；序贯路径在每次物理故障需求时对包交付作伯努利抽样，同时保留元件当前健康、
QoS 门限和共享路径。结果标志 \fld{packet\_delivery\_sampled} 只有实际执行过需求抽样时才为真，不能由配置字段预置。

\subsection{保护类抽样}
给定当前确定信息状态，模拟器生成条件保护类，再按类概率抽取主清除、后备清除或未清除。机械成功概率在每次需求上重新
抽样；不会把断路器某次拒动永久保留为元件故障，除非调用方另建一个具有修复过程的断路器状态元件。

\subsection{停电区间并集}
多个故障可能产生重叠停电。简单累加每个事件的 $S\tau$ 会重复计算相同负荷和时段。当前序贯接口对同一代表性切负荷强度
形成区间集合，排序后合并：
\begin{equation}
 \mathcal I=\bigcup_k[t_k^{start},t_k^{end}).
\end{equation}
LOLE 取并集长度，LOLF 取不连续合并区间数，EENS 取区间长度乘切负荷。结果返回合并区间数和
\fld{outage\_interval\_union\_modelled}。若同时需要多个不同负荷点，应扩展为逐负荷点区间并集，不能用系统最大切负荷代替。

\subsection{静态与序贯结果的差异来源}
静态事件树用环境概率和代表时长直接求期望，适合参数筛选；序贯队列显式表示故障先后、电池跨事件状态、重叠停电和共因
持续过程。两者在低频、无重叠、每次事件前电池充满且持续时间固定时应趋近；电池未充满、事件重叠或共因持续时，序贯结果
不应被强求等于静态结果。

\subsection{边界和复杂度}
静态元件位枚举复杂度 $O(2^{n_c}n_fn_p)$；序贯队列主要复杂度为 $O(N_e\log N_e+N_kn_p)$，其中 $N_e$ 是事件数、
$N_k$ 是物理故障需求数。\fld{maximum\_events} 防止异常高率输入导致无界运行；达到上限明确报错，不返回不完整年度均值。

\subsection{解析回归算例}
测试安排两个确定物理故障，其中第二次发生在信息节点备用电池耗尽之后。第一次自动保护成功，第二次因信息功能丢失进入
保守后果。预期 EENS 为 $10(1+1/3600)$ MWh，LOLE 为 $1+1/3600$ 小时；该非整数小时结果证明事件时标未被整点化。
另一个测试构造重叠停电，验证 LOLE 取区间并集而非各事件持续时间之和。

\subsection{输出真实性标志}
序贯结果分别报告亚小时计时、物理供电耦合、共因、保护概率、包交付、电池轨迹和区间并集标志。标志由实际执行路径设置：
空共因数组不能得到“共因已建模”，没有电池节点不能得到“电池轨迹已建模”。该约束防止 GUI 仅因用户勾选选项就显示
并未消费的数据能力。

\section{信息系统与智能决策的参数影响及方法对照}\label{sec:rel-cyber-impact-cases}

信息系统不是一个可以直接乘在故障率上的总可用率。本章从共享元件、备选路径、服务质量、共因环境、物理供电、电池自治和
智能决策依次推导联合事件概率，并说明这些因素如何改变静态 FMEA、仅保护和信息物理联合三种结果。重点是保持同一故障
频率与同一后果基准，使差值能够归因。

\subsection{元件状态、路径状态与功能状态}
设信息元件集合为 $\mathcal C=\{1,\ldots,n\}$，元件 $i$ 的工作状态为 $x_i\in\{0,1\}$，固有可用率为 $a_i$。
在独立环境中，完整状态 $x$ 的概率为
\begin{equation}\label{eq:rel-cyber-state-probability}
 p(x)=\prod_{i=1}^{n}a_i^{x_i}(1-a_i)^{1-x_i}.
\end{equation}
一条路径 $r$ 是元件索引集合 $C_r$。路径结构可用当且仅当所有成员工作：
\begin{equation}
 X_r(x)=\prod_{i\in C_r}x_i.
\end{equation}
功能 $f$ 可以有若干备选路径 $\mathcal R_f$，其结构状态为
\begin{equation}
 X_f(x)=1-\prod_{r\in\mathcal R_f}(1-X_r(x)).
\end{equation}
平台不先把每条路径压缩成独立概率再做并联，因为两条路径可能共享交换机、电源或通信中心。它枚举元件位状态并对满足
$X_f=1$ 的状态概率求和，因此共享依赖只计算一次。

\subsection{串联、独立并联与共享并联的数值差异}
两元件串联且 $a_1=a_2=0.99$ 时，功能可用率为 $0.99^2=0.9801$。两条完全独立的单元件并联路径可用率为
\begin{equation}
 1-(1-0.99)^2=0.9999.
\end{equation}
若两条名义路径都经过可用率 $0.98$ 的共享中心，路径末端各为 $0.99$，则联合可用率不是把
$0.98\times0.99$ 当成两条独立路径并联，而是
\begin{equation}
 A_f=0.98\left[1-(1-0.99)^2\right]=0.979902.
\end{equation}
可见共享中心使“新增冗余链路”的收益从接近两个九降到受 $0.98$ 上限约束。程序的最小割集会识别共享中心单点割集和
两个末端同时失效的二元割集；这比只报告路径条数更能说明薄弱环节。

\subsection{QoS 是确定性准入，不是事后折扣}
每个信息元件给出时延 $d_i$、抖动 $j_i$ 和包交付概率 $q_i$。路径聚合为
\begin{align}
 d_r&=\sum_{i\in C_r}d_i,\\
 j_r&=\sum_{i\in C_r}j_i,\\
 q_r&=\prod_{i\in C_r}q_i.
\end{align}
功能契约给出 $d_f^{max}$、$j_f^{max}$、$q_f^{min}$。路径必须同时满足
\begin{equation}
 d_r\le d_f^{max},\qquad j_r\le j_f^{max},\qquad q_r\ge q_f^{min},
\end{equation}
才能参与结构可用性。一个固有可用率为 $0.999999$ 但保护消息时延 200 ms、契约上限 20 ms 的路径，对该保护功能的贡献
仍为零。把它乘成 $0.999999\times20/200$ 会制造没有物理依据的部分成功状态。

QoS 改善的影响具有门槛性。路径时延从 25 ms 降到 21 ms 时，20 ms 契约下功能可用率不变；降到 19 ms 时，整条路径
进入可用集合，功能可用率发生离散跃迁。因此参数扫描应同时报告准入路径集合，而不能只画一条平滑灵敏度曲线。

\subsection{共因环境的条件枚举}
设环境状态 $g\in\mathcal G$ 的概率为 $p_g$，环境指定失效共因组和失电物理母线。条件于 $g$ 后，属于失效共因组的元件
状态强制为零；其余元件再按式~\eqref{eq:rel-cyber-state-probability} 枚举。联合概率为
\begin{equation}
 p(g,x)=p_g p(x\mid g),\qquad
 \sum_g\sum_xp(g,x)=1.
\end{equation}
若正常环境概率为 $0.98$、通信中心共因失效环境为 $0.02$，即使所有独立元件可用率趋近于一，依赖该中心的功能可用率也
不会超过 $0.98$。这解释了独立元件升级后的收益饱和。

环境概率必须构成完备分布。若用户只给出若干异常环境且概率和小于一，入口不会暗自归一化并改变事件频率；调用方应显式
加入正常环境。概率为负、和超出容差或共因组引用不存在均拒绝执行。

\subsection{信息设备供电与电池自治}
信息元件可以绑定物理母线稳定 ID。若环境使该母线失电，元件仅在备用电池能够覆盖信息停电持续时间时继续工作。设电池
能量为 $E_b$ Wh、功耗为 $P_b$ W，则自治时间
\begin{equation}
 T_b=\begin{cases}E_b/P_b,&P_b>0,\\+\infty,&P_b=0.\end{cases}
\end{equation}
静态事件树中，若 $T_b\ge T_{out}$，该次环境下供电依赖不使元件失效；否则元件在整个代表事件中视为不可用。序贯队列
则在 $t_{deplete}=t_{out,start}+T_b$ 产生电池耗尽事件，因此同一次物理停电的前半段和后半段可对应不同的信息功能状态。

例如 $E_b=20$ Wh、$P_b=10$ W，自治时间为 2 h。面对 1.5 h 停电，静态与序贯模型都保持功能；面对 3 h 停电，静态
模型把功能判为失效，序贯模型保留前 2 h 自动操作能力，最后 1 h 进入人工或失败类。若关键控制动作发生在停电后 10 min，
序贯结果会优于静态二值近似。

\subsection{联合功能而非边际功能的必要性}
检测、主跳闸、后备跳闸、隔离和恢复可能依赖同一信息元件。分别计算五个边际可用率再相乘隐含了功能独立，通常会低估共享
中心失效的相关性。平台对同一个状态 $x$ 同时计算五个功能布尔值，得到功能向量
\begin{equation}
 y(x)=(y_D,y_{TP},y_{TB},y_I,y_R),
\end{equation}
再按相同向量归并概率。若一个共享中心控制所有功能，只有“全部可用”和“全部不可用”两个主要类；把五个相同边际概率
相乘会错误地产生大量部分成功概率。

对故障场景 $e$，保护信息联合类概率可写为
\begin{equation}
 \pi(c\mid e)=\sum_{g,x}p(g,x)\,
 \mathbf1\{\Phi(e,g,x)=c\},
\end{equation}
其中 $\Phi$ 顺序执行检测、主后备保护、隔离、恢复与 DER-FRT 分类。类概率和必须为一；概率为零的类不参与后果聚合。

\subsection{检测混淆矩阵}
真实故障类别为 $z$，检测输出为 $\widehat z$，混淆矩阵元素
$C_{z\widehat z}=P(\widehat z\mid z)$。每个真实类别行必须非负且和为一。检测结果还带有时延类 $\ell$，其概率为
$P(\ell\mid z,\widehat z)$。检测类的完整概率为
\begin{equation}
 p(z,\widehat z,\ell)=p(z)C_{z\widehat z}P(\ell\mid z,\widehat z).
\end{equation}
漏检进入检测失败类；错误分类可以进入错误隔离区，并不等同于漏检。提高总检出率若同时把概率从正确分类移向错误分类，可能
增加切负荷，因此“召回率更高”不能单独证明可靠性改善。

检测时延会扩展故障未隔离窗口。若该窗口切负荷为 $S_d$，时延从 $t_a$ 增至 $t_b$ 的 EENS 增量为
\begin{equation}
 \Delta E_d=\lambda p(z,\widehat z,\ell)S_d
 \frac{t_b-t_a}{3600}.
\end{equation}
秒到小时的单位换算必须在此处完成。对 10 MW、2 次/年、概率 0.5 的类别，时延增加 0.2 s 只增加
$5.56\times10^{-4}$ MWh/年；若程序显示 2 MWh/年，通常是遗漏了 $3600$。

\subsection{最小风险隔离}
候选隔离动作 $a$ 具有切除负荷 $S_a$、误隔离代价 $C_a^{false}$、漏隔离风险 $C_a^{miss}$ 和安全约束。给定检测后验
$b_z=P(z\mid\widehat z)$，动作风险为
\begin{equation}
 R(a\mid b)=C_a^{op}+\sum_zb_z C(a,z),
\end{equation}
其中 $C_a^{op}$ 可包含切负荷能量和开关动作成本。算法先删除引用不存在、无法覆盖故障或违反安全门槛的动作，再在可行集合
中选最小风险动作。若没有可行动作，结果显式为失败，不返回默认第一项。

考虑后验 $b=(0.8,0.2)$，窄区动作对类别 1 的代价为 1、对类别 2 的漏隔离代价为 20，宽区动作对两类代价均为 5，则
\begin{align}
 R_{narrow}&=0.8\times1+0.2\times20=4.8,\\
 R_{wide}&=0.8\times5+0.2\times5=5.0.
\end{align}
窄区略优；若类别 2 后验升到 0.25，窄区风险变为 5.75，算法改选宽区。这个切换点来自后验和代价，而不是写死的置信度
阈值。

\subsection{风险约束恢复}
恢复动作 $a$ 给出恢复负荷 $L_a$、不安全概率 $p_a^{unsafe}$、执行成功率 $p_a^{exec}$ 和期望剩余损失。平台先施加
\begin{equation}
 p_a^{unsafe}\le\bar p_{unsafe}
\end{equation}
硬约束，再在可行动作中最小化风险调整损失或最大化风险调整恢复量。一个恢复 10 MW 但不安全概率 0.2 的动作，在门槛
0.05 下不会因收益高而被选中。收紧门槛可能提高短期 EENS，却减少二次事故暴露；可靠性和安全性不能压成一个没有单位说明
的加权分数。

若动作恢复 $L_a$，执行失败时仍损失 $S_f$，成功后损失 $S_s$，则条件期望损失
\begin{equation}
 \overline S_a=p_a^{exec}S_s+(1-p_a^{exec})S_f.
\end{equation}
信息系统影响 $p_a^{exec}$ 时，应改变该条件期望或类概率；不能只在结果上附加“智能恢复已启用”布尔值。

\subsection{有限 POMDP 精确信念树}
部分可观测恢复模型用状态 $s\in\mathcal S$、动作 $a\in\mathcal A$、观测 $o\in\mathcal O$、转移
$T(s'\mid s,a)$、观测模型 $O(o\mid s',a)$ 和阶段代价 $c(s,a)$。给定信念 $b(s)$，观测后的贝叶斯更新为
\begin{equation}
 b'(s')=\frac{O(o\mid s',a)\sum_sT(s'\mid s,a)b(s)}
 {\sum_{\bar s}O(o\mid\bar s,a)\sum_sT(\bar s\mid s,a)b(s)}.
\end{equation}
有限时域价值递推为
\begin{equation}
 V_h(b)=\min_a\left[\sum_sb(s)c(s,a)+
 \sum_oP(o\mid b,a)V_{h-1}(\tau(b,a,o))\right].
\end{equation}
实现遍历完整有限信念树，不用固定“先测量后动作”规则。专项算例构造一个立即恢复风险高、先获取信息后可按状态选择动作的
系统，精确解必须选择信息收集动作。非法转移概率、观测概率或维数不一致直接拒绝。

\subsection{连续时间序贯事件队列}
信息元件 $i$ 的失效率、修复率为 $\lambda_i^c,\mu_i^c$，其两状态 CTMC 的下一事件时间分别服从指数分布。物理故障、信息
失效、修复、共因开始/结束、包投递、母线失电/复电和电池耗尽都进入按时间排序的同一队列。处理同一时刻事件时采用确定性
优先级，避免容器迭代顺序改变结果。

每次物理故障到达时，算法读取该时刻的信息状态和剩余电池能量，生成保护、隔离和恢复后果区间。年度 LOLE 不是简单相加
所有区间长度，而是对区间并集求测度：
\begin{equation}
 \mathrm{LOLE}_{y}=\left|\bigcup_{k\in\mathcal O_y}
 [t_k^{start},t_k^{end})\right|.
\end{equation}
EENS 对同时停电区间的切负荷可按系统约定聚合，但绝不能让同一 1 h 停电因两个重叠故障变成 2 h LOLE。专项测试构造
重叠事件并检查并集长度上界。

\subsection{包交付抽样与结构状态的区别}
元件结构工作不代表本次报文必达。对已通过 QoS 门槛的路径，包交付仍按 $q_r$ 抽样；若有多条可用路径，功能成功是至少一条
报文成功。结构失效是持续状态，包丢失是按消息发生的随机事件。把包交付概率并入元件长期可用率会错误地产生修复过程，
把结构失效只按单包抽样又会丢失持续相关性。

序贯结果返回包交付是否被抽样、信息供电耦合是否生效、电池能量轨迹是否生效和最小剩余能量。只有这些执行证据为真时，
才能声明相应因素进入年度指标。

\subsection{三级方法的共同基准}
三级对照的静态行、仅保护行和联合行必须共享：场景 ID、起始频率 $\lambda_e$、主/后备/未清除切负荷、隔离后切负荷、恢复后
切负荷、自动/人工恢复时间和修复时间。静态行把未清除切负荷保持到修复：
\begin{equation}
 E_e^S=S_{e,U}\tau_{e,repair}.
\end{equation}
仅保护行删除全部信息元件和功能绑定，只保留保护类：
\begin{equation}
 E_e^P=\sum_{c_P}\pi_P(c_P\mid e)W(e,c_P).
\end{equation}
联合行保留完整信息状态：
\begin{equation}
 E_e^C=\sum_{g,x,c_P}p(g,x)\pi(c_P\mid e,g,x)W(e,g,x,c_P).
\end{equation}
最终三行分别乘同一个 $\lambda_e$ 并求和。

\subsection{保护收益和信息代价}
平台定义
\begin{align}
 \Delta E_{prot}&=\mathrm{EENS}^{S}-\mathrm{EENS}^{P},\\
 \Delta E_{cyber}&=\mathrm{EENS}^{C}-\mathrm{EENS}^{P}.
\end{align}
$\Delta E_{prot}$ 表示相对静态未清除基线的保护收益；$\Delta E_{cyber}$ 表示相对理想信息条件的联合增量。二者不是 Shapley
分解，也不含无故障误动，不能相加后宣称得到全部信息物理贡献。

在后果单调且信息仅会使功能从成功变失败时，$\Delta E_{cyber}\ge0$。若远方自动控制错误会造成误切，而信息失效反而阻止
该错误动作，则差值可以为负。程序保留实际符号，不把负值截零；报告必须解释对应类的后果。

\subsection{闭式三级对照算例}
基准场景取年频率 2 次，未清除切负荷 10 MW、修复 10 h，因此静态行
\begin{equation}
 \mathrm{EENS}^{S}=2\times10\times10=200\ \text{MWh/年},
 \quad \mathrm{LOLE}^{S}=20\ \text{h/年}.
\end{equation}
主保护成功率 0.9、后备成功率 1，主/后备清除延时分别为 0.1 s 和 0.5 s；隔离、恢复后果使用同一三窗口。闭式事件树得到
\begin{align}
 \mathrm{EENS}^{P}&=0.4007777778\ \text{MWh/年},\\
 \mathrm{LOLE}^{P}&=0.2000777778\ \text{h/年}.
\end{align}
共享信息中心可用率为 0.8，且检测、跳闸、隔离、恢复都绑定该中心时，联合结果为
\begin{align}
 \mathrm{EENS}^{C}&=21.5217333333\ \text{MWh/年},\\
 \mathrm{LOLE}^{C}&=4.1600622222\ \text{h/年}.
\end{align}
因此保护收益为 $199.5992222222$ MWh/年，信息失效相对仅保护的增量为 $21.1209555555$ MWh/年。这个结果同时证明保护
逻辑和信息状态都实际改变了后果；若三行完全相同，说明运行链没有消费相应模型。

\subsection{信息可用率扫描}
保持静态与保护输入不变，把共享中心可用率从 0.8 降到 0.4。静态行必须严格保持 200 MWh/年，保护行也必须保持
$0.4007777778$ MWh/年；只有联合行上升。该测试比“联合结果非零”更有辨识力，因为它同时检查比较基准未被信息参数污染。

若有两条独立路径，单一路径可用率从 0.8 降到 0.4 对功能可用率的影响取决于另一条路径；若两条路径共享中心，中心参数
的边际影响最大。参数扫描必须与最小割集一起解释，不能只比较一个总百分比。

\subsection{有限差分敏感性检查}
对任一连续概率参数 $\theta$，可用中心差分
\begin{equation}
 D_h(\theta)=\frac{E(\theta+h)-E(\theta-h)}{2h}
\end{equation}
检查解析方向。$h$ 应避免跨越 $[0,1]$ 边界和 QoS 离散门槛。对纯串联独立路径，解析导数可直接由乘积得到；对共享依赖、
共因和功能绑定，精确状态枚举提供无抽样噪声的差分基准。序贯模型有随机误差时，应使用公共随机数或置信区间比较，不能以
单次样本差的符号判定错误。

\subsection{信息物理验收条件}
信息系统相关实现至少满足：元件状态概率和为一；共享路径与直接位枚举一致；QoS 超限路径不进入可用集合；最小割集确实使
功能失败且删除任一成员后不再是割集；共因环境条件概率守恒；母线失电和电池耗尽改变后续事件结果；包投递抽样标志为真；
检测混淆矩阵逐行守恒；隔离与恢复只在可行动作中选择；POMDP 与小规模手算价值一致；停电区间使用并集；三级对照的静态行
不随信息参数变化而联合行按类后果变化。实现分别位于
\srcpath{src/reliability/reliability\_assessment.cpp}、
\srcpath{src/reliability/intelligent\_cyber\_physical.cpp} 和
\srcpath{src/reliability/protection\_frt.cpp}，HTTP 与 GUI 的三级结果使用同一 C++ 聚合入口。

\section{智能检测、隔离、恢复与在线动态闭环}\label{sec:rel-intelligent-online}

\subsection{混淆矩阵的概率契约}
检测模型把无故障类记为 $0$，真实故障类和报告类为 $1,\ldots,K$。矩阵
$C_{ij}=\Pr(\widehat K=j\mid K=i)$ 的每一行必须非负且和为一。对真实故障 $k$：
\begin{align}
 P_D(k)&=1-C_{k0},\\
 P_C(k)&=C_{kk},\\
 P_M(k)&=C_{k0}.
\end{align}
检测与正确分类不是同一指标：报告了错误故障段属于“已检测但误分类”，会改变隔离候选和后果区，而不是自动计为漏报。

\subsection{误报暴露基}
无故障行的 $1-C_{00}$ 是每判别窗误报概率。年度误报频率必须乘实际判别窗数：
\begin{equation}
 \lambda_{FA}=\nu_{win}(1-C_{00}).
\end{equation}
若连续采样高度相关，$\nu_{win}$ 应取独立决策次数或直接使用现场告警率；把每毫秒采样点都当独立窗会把误动频率夸大。

\subsection{检测时延类积分}
每个真实故障的检测时延用互斥离散类 $(p_h,t_h)$ 近似分布。与报告类别联合后，类概率为 $C_{kj}p_h$，及时检测概率为
\begin{equation}
 P_{timely}(k)=\sum_{j>0}\sum_{h:t_h\le T_{lim}}C_{kj}p_h.
\end{equation}
条件平均检测时延为
\begin{equation}
 \mathbb E[T_D\mid D]=
 \frac{\sum_{j>0,h}C_{kj}p_ht_h}{P_D(k)}.
\end{equation}
平台输出每个联合类，因而调用方可把时延映射为非线性故障能量，而不是只消费均值。

\subsection{后验风险隔离}
给定故障段后验 $b_k$ 和候选隔离方案 $q$，覆盖概率为
\begin{equation}
 B_q=\sum_{k\in\mathcal K_q}b_k.
\end{equation}
候选期望成本为
\begin{equation}
 J_q=c_{miss}(1-a_qB_q)+c_{shed}P_q^{isolated}+c_tT_q,
\end{equation}
其中 $a_q$ 是开关动作成功概率。若所需信息功能不可用，候选标记为不可执行，不参与最小化；不把不可执行方案的成本设成
某个任意大常数，以免数值尺度变化后被错误选中。

\subsection{隔离结果的含义}
最小风险方案不是“预测最可能故障段”的同义词。高后验但隔离范围巨大时，另一个覆盖多个相邻段、切负荷较小的候选可能有
更低风险。输出保留每个候选的可执行性、覆盖概率和成本，便于审计权重。后验必须和为一，覆盖索引越界或负成本直接拒绝。

\subsection{风险约束恢复}
候选恢复动作 $a$ 给出恢复负荷、开关成本、不安全概率、执行成功概率、远程时间和人工时间。硬风险约束为
\begin{equation}
 p_a^{unsafe}\le\overline p^{unsafe}.
\end{equation}
可远程执行时，期望恢复量为 $p_a^{exec}P_a^{restore}$；信息功能不可用但允许人工时，使用人工时长和相应执行概率。目标为
\begin{equation}
 J_a=c_E\left(P^{d}-p_a^{exec}P_a^{restore}\right)T_h+c_{sw,a}.
\end{equation}
违反风险约束的动作进入拒绝列表，不能通过较高恢复量补偿安全违规。

\subsection{人工回退}
人工回退不是信息功能“成功”。它改变动作时长和可能的执行概率，但仍允许系统最终恢复。输出
\fld{used\_manual\_fallback} 和 \fld{selected\_duration\_s}，年度后果据此选择自动或人工切换窗口。
若人工时间也不可用或动作无本地操作能力，应把持续时间视为无穷并从候选集中移除。

\subsection{有限 POMDP 模型}
平台提供有限状态、动作、观测和有限时域的精确信念树。模型为
\begin{equation}
 \mathcal M=(\mathcal S,\mathcal A,\mathcal O,T,O,c,c_T,b_0,H).
\end{equation}
转移张量 $T^a_{ss'}$ 和观测张量 $O^a_{s'o}$ 的每一概率行必须和为一；初始信念 $b_0$ 也必须归一。

\subsection{信念更新}
执行动作 $a$ 后先预测
\begin{equation}
 \bar b(s')=\sum_sT^a_{ss'}b(s),
\end{equation}
观察到 $o$ 后更新
\begin{equation}
 b'(s')=\frac{O^a_{s'o}\bar b(s')}
 {\sum_jO^a_{jo}\bar b(j)}.
\end{equation}
零概率观测分支不递归。实现不做信念量化或启发式剪枝，因此对给定初始信念和有限时域是精确的。

\subsection{Bellman 递推}
终端值为 $V_0(b)=\sum_sb_sc_T(s)$。剩余 $h$ 步时
\begin{equation}
 V_h(b)=\min_a\left[\sum_sb_sc(s,a)+
 \sum_o\Pr(o\mid b,a)V_{h-1}(\tau(b,a,o))\right].
\end{equation}
输出最优首动作、期望总成本和访问信念节点数。由于复杂度随时域、动作和观测指数增长，接口对状态数、动作数、观测数和时域
设明确上限；超限拒绝，不声称近似求解了大规模 POMDP。

\subsection{信息收集动作算例}
测试构造“立即隔离”和“先观察”两类动作。立即隔离成本确定，观察动作有小成本但产生能区分真实故障的观测，下一步可选
针对性隔离。精确信念树应选择先观察，并得到低于立即隔离的期望总成本。该测试同时验证观测价值确实通过后验改变后续动作，
不是只在结果中返回一个固定策略名。

\subsection{Mass-Matrix DAE 在线入口}
在线保护输入包括故障交流母线、相别、发生时刻、故障阻抗、被保护支路稳定 ID、继电模型、断路器链、CT/PT 参数、DER
监视器和动态求解选项。入口强制求解器为 Mass-Matrix DAE，打开逐步快照、设备输出和 DER 保护状态机。

\subsection{第一次发现运行}
第一次运行只施加故障，不预设保护动作。对故障母线相电压 $V_f(t)$ 和故障导纳 $Y_f$，继电一次电流轨迹取
\begin{equation}
 I_f(t)=\begin{cases}0,&t<t_f,\\Y_fV_f(t),&t\ge t_f.\end{cases}
\end{equation}
该轨迹经过 CT/PT 后送入继电判据，得到 $t_{cmd}$。因此换流器限流、GFL/GFM 动态和网络电压变化会影响继电到达，
而不是由输入字段预置“已经消费动态反馈”。

\subsection{第二次闭环运行}
若继电动作，计算 $t_{clear}$ 并重新构建同一动态系统，同时加入故障、交流支路跳闸和故障清除事件。两种清除事件必须都在
\fld{applied\_event\_records} 中出现；缺任一项即视为闭环失败。若继电不动作，结果仍成功返回发现轨迹和未清除 DER 轨迹，
但 \fld{protection\_action\_applied=false}。

\subsection{DER-FRT 轨迹消费}
每个 DER 监视器用稳定字符串 ID 和交流母线 ID 绑定。发现运行保留未清除故障轨迹；闭环运行保留对应保护清除轨迹。
IEEE~1547 状态机逐点消费电压幅值和相角，输出是否曾跳闸、是否重连、首跳/首重连时间和终端恢复比例。

\subsection{GFL/GFM 动态反馈证据}
\fld{converter\_dynamic\_feedback\_consumed} 只有在动态快照中实际找到 GFL、GFM 或 VSC 设备输出时才为真。测试分别建立
真实 \srcpath{VSCGridFollowing} 和 \srcpath{VSCGridForming} 动态设备，检查输出类型、稳定组件 ID 和非空遥测；仅把
\fld{ac\_grid\_forming} 设为真而没有设备输出不能通过。

\subsection{主后备在线年度耦合}
\srcpath{compare\_online\_protection\_cyber\_reliability} 对同一故障运行主、后备两个 DAE 配置。它把两次发现运行的
CT/PT 实测轨迹送入保护配合，把主闭环 DER 轨迹绑定主清除类、后备闭环轨迹绑定后备类、发现运行未清除轨迹绑定拒清类。
这三套轨迹不能互换，否则后备清除时间改变的 FRT 后果会丢失。

\subsection{时间一致性守卫}
事件树重新评估继电轨迹后得到的主、后备清除时刻必须与 DAE 结果在最大时间步容差内一致：
\begin{equation}
 |t_{clear}^{tree}-t_{clear}^{DAE}|\le\max(\Delta t_p,\Delta t_b)+10^{-9}.
\end{equation}
不一致表示轨迹、整定或断路器链在转换中被改变，接口明确失败。成功后才设置
\fld{online\_network\_dae\_coupled=true} 和 \fld{dae\_trajectories\_consumed\_by\_event\_tree=true}。

\subsection{数值成本模型}
每个在线场景需两次主保护 DAE 和两次后备保护 DAE（发现与闭环；不动作时只有发现），再做一次离散信息状态枚举。
粗略成本为
\begin{equation}
 T\approx N_s(2T_{DAE,p}+2T_{DAE,b}+2^{n_c}T_{tree}).
\end{equation}
因此在线耦合适合代表性故障类或重要场景，不应未经筛选地对百万次年度事件逐次求 DAE。大样本研究可先用在线接口生成
条件类库，再由序贯队列抽取类；类库的拓扑、工况和参数适用域必须随结果保存。

\subsection{动态耦合的有效边界}
当前在线入口使用故障母线电压与故障导纳形成继电电流，适合验证因果闭环和设备动态影响；它不是完整三相线路端保护模型，
未替代专业 EMT、行波或互感器磁滞仿真。结果的 \fld{model\_scope} 明确为
\srcpath{mass-matrix-dae+trajectory-derived-relay+branch-trip+der-frt}，工程定值校核仍需匹配所需保真度。

\section{在线动态耦合、年度消费与验证证据}\label{sec:rel-online-dae-validation}

本章说明在线保护模式如何从网络状态生成量测轨迹、如何把保护动作反馈给第二遍动态仿真，以及这些轨迹如何进入年度可靠性
事件树。验收重点不是“调用了动态求解器”，而是主清除、后备清除和未清除三种后果确实消费不同轨迹，且 DAE 清除时刻与
事件树清除时刻一致。

\subsection{动态网络与设备状态}
在线模式使用质量矩阵形式的微分代数方程
\begin{align}
 M\dot x&=f(x,v,u,t),\\
 0&=g(x,v,u,t),
\end{align}
其中 $x$ 为发电机、换流器、控制器等动态状态，$v$ 为交流节点复电压与可选直流代数量，$u$ 包含故障和开关事件。
Backward Euler 步在 $t_{k+1}$ 求解
\begin{equation}
 R(z_{k+1})=
 \begin{bmatrix}
 M(x_{k+1}-x_k)/\Delta t-f(x_{k+1},v_{k+1},u_{k+1})\\
 g(x_{k+1},v_{k+1},u_{k+1})
 \end{bmatrix}=0,
\end{equation}
其中 $z=(x,v)$。每个时间步用稀疏 Newton 迭代；求解失败时在线可靠性入口返回具体阶段和动态消息，不生成后续事件类。

动态时域必须满足 $t_{end}>t_{fault}$，时间步必须为有限正数。故障阻抗
$Z_f=R_f+jX_f$ 不能为零向量；当前输入指定交流母线稳定 ID、相别和故障阻抗。投影后找不到母线、被保护支路已退出或稳定
ID 不存在时，入口在第一遍 DAE 前拒绝。

\subsection{第一遍故障发现运行}
第一遍在原系统上于 $t_f$ 施加故障并联导纳
\begin{equation}
 Y_f=1/(R_f+jX_f).
\end{equation}
对故障母线目标相，快照电压为 $V_k$，一次故障电流为
\begin{equation}
 I_k=\begin{cases}0,&t_k<t_f,\\Y_fV_k,&t_k\ge t_f.\end{cases}
\end{equation}
程序不是把 $1/Z_f$ 当成恒定电流，而是每个快照用动态电压更新，因此电源控制、换流器限流和电压恢复都会改变继电器输入。

一次 $V_k,I_k$ 先通过 CT/PT 精确离散环节，再形成
\verb|measured_relay_trajectory|。继电器判据只消费这条二次量轨迹。若轨迹不达到拾取条件，结果仍可成功返回发现运行与 DER
分类，但不会伪造支路跳闸和闭环运行。

\subsection{动作时刻与断路器因果链}
继电器得到命令时刻 $t_{cmd}$ 后，按式~\eqref{eq:rel-breaker-chain} 计算清除时刻。清除时刻必须不晚于动态时域终点；否则
延长时域或调整整定，不能把动作夹在最后一个采样点。保护成功概率不参与 DAE 动作时刻，它在年度事件树中决定“设备按需
成功”分支；这一区分使物理动作时间与统计拒动概率保持独立。

例如故障发生于 0.05 s，定时限延时 0.02 s，断路器机械延时 0.01 s，其余延时为零，则理论清除时刻为 0.08 s。离散轨迹
首次拾取受时间网格影响，实际值允许与理论连续值相差一个 $\Delta t$。程序在事件树转换后要求
\begin{equation}\label{eq:rel-dae-event-time-consistency}
 |t_{clear}^{DAE}-t_{clear}^{tree}|\le
 \max(\Delta t_P,\Delta t_B)+10^{-9}\ \mathrm{s}.
\end{equation}

\subsection{第二遍动作反馈运行}
第二遍从同一原始系统重建动态模型，仍在 $t_f$ 施加相同故障，并在 $t_{clear}$ 同时施加：
\begin{enumerate}
 \item 被保护交流支路稳定 ID 的跳闸事件；
 \item 故障母线的清除故障事件；
 \item 开断后所有无源交流岛的负荷退出事件。
\end{enumerate}
如果只在第一遍轨迹上截断电流而不重新求解，支路开断后的节点电压、频率和换流器状态都不会改变，不能称为闭环。结果中的
\fld{protection\_action\_applied} 只有在动态事件记录同时包含目标支路跳闸和故障清除时才为真。

失电岛负荷退出由图连通性自动生成，结果返回
\fld{deenergized\_ac\_bus\_ids} 和
\fld{dead\_island\_load\_shedding\_applied}。这些动态负荷退出只保证 DAE 状态符合拓扑；年度失负荷仍由三窗口后果统计，
不会重复从动态功率积分一次。

\subsection{为什么需要主、后备各两遍}
主保护和后备保护可能有不同的支路、整定、CT/PT、断路器延时和监视 DER。若先做一遍故障轨迹，再只把后备延时加在主保护
结果上，会漏掉后备支路跳开后的不同拓扑。在线年度场景因此包含两个
\srcpath{OnlineProtectionDAEInput}：主输入和后备输入。每个输入各自执行发现与反馈两遍，总计四遍 DAE。

主清除类消费主输入的闭环 DER 轨迹；后备清除类消费后备输入的闭环 DER 轨迹；未清除类消费主输入第一遍持续故障轨迹。
记三组轨迹为 $\Gamma_P,\Gamma_B,\Gamma_U$，则 DER 终态分别为
\begin{equation}
 q_P=\Psi(\Gamma_P),\qquad q_B=\Psi(\Gamma_B),\qquad
 q_U=\Psi(\Gamma_U).
\end{equation}
如果主保护在 DER 闭锁超时前清除而后备保护在超时后清除，可能有 $q_P=$穿越、$q_B=$跳闸。使用一条共享轨迹会直接改变
后果类概率与切负荷。

\subsection{稳定身份一致性}
主、后备监视器必须以同一个 DER 稳定 ID 表示同一设备。适配器不能分别生成
\verb|device:primary| 与 \verb|device:backup|，否则事件树无法判断两条轨迹是否属于同一设备。HTTP 在线模式用
\verb|ac_bus:<id>:online-der-monitor| 作为默认监视身份，主后备完全一致；轨迹容器的位置只用于内部并行数组。

转换前检查主后备监视器维数与稳定 ID；不一致时拒绝。这个检查曾在 E2E 中捕获适配器错误：两遍 DAE 都成功，但事件类
生成器拒绝消费错位轨迹。修正稳定身份后，路由返回的
\fld{dae\_trajectories\_consumed\_by\_event\_tree} 才能为真。

\subsection{换流器动态反馈证据}
若系统含 GFL 或 GFM VSC，动态快照的设备输出应出现对应设备类型与稳定组件索引。在线结果扫描发现和闭环两遍快照，只有
读取到真实换流器动态输出才置
\fld{converter\_dynamic\_feedback\_consumed=true}。专项测试分别放入 GFL 与 GFM VSC，要求输出类型、组件索引和状态值
非空。这个标志防止仅凭系统含有 VSC 就宣称换流器动态参与。

GFL 限流可降低故障电流并延迟过流保护；GFM 内电势与虚拟阻抗会改变故障期间电压和电流相位。两者对保护动作的影响由
第一遍 DAE 轨迹自然进入，不在可靠性层另乘经验修正系数。

\subsection{在线事件树的构造}
对在线场景 $e$，主后备 DAE 产生测量轨迹和 DER 轨迹。适配器构造与给定轨迹入口相同的
\srcpath{ProtectionCyberEventInput}：
\begin{align}
 \text{primary.trajectory}&\leftarrow\Gamma_{relay,P}^{DAE},\\
 \text{backup.trajectory}&\leftarrow\Gamma_{relay,B}^{DAE},\\
 \text{ders}&\leftarrow\Gamma_{P},\\
 \text{backup\_ders}&\leftarrow\Gamma_{B},\\
 \text{uncleared\_ders}&\leftarrow\Gamma_{U}.
\end{align}
信息元件、功能、环境、绑定和三窗口后果保持用户输入。这样在线与给定轨迹模式的差别仅在轨迹来源，后续类生成与年度聚合
完全共享，便于做同口径对照。

\subsection{年度三窗口积分}
每个联合类的后果分为电气清除、切换恢复和残余修复三段。设清除切负荷 $S_c$、清除时刻 $t_c$，切换后切负荷 $S_s$、
恢复时间 $\tau_s$，残余切负荷 $S_r$、修复时间 $\tau_r$，则
\begin{equation}
 W_c=S_c\frac{t_c}{3600}+S_s\tau_s+S_r\max(0,\tau_r-\tau_s).
\end{equation}
具体类选择主清除、后备清除或未清除的 $S_c,t_c$，再按检测、隔离、恢复功能状态选择自动或人工窗口。所有持续时间先检查
单调关系；自动恢复不能晚于修复，负持续时间拒绝。

在线动态只覆盖秒级清除与 DER 状态，小时级切换和修复仍由可靠性后果模型给出。把 DAE 一直运行 10 h 既无必要，也会把
运维恢复决策误当成电磁或机电动态。

\subsection{在线与给定轨迹的差异解释}
给定轨迹适合保护定值回放、标准波形和闭式测试。它的优点是确定、快速、容易复现；缺点是整定、支路跳闸或换流器控制变化
不会反过来改变轨迹。在线模式的轨迹由当前系统生成，动作反馈会改变第二遍网络，适合研究限流、低电压穿越和拓扑分离。

两种模式使用同一条给定轨迹且关闭所有动态反馈时，结果应在时间离散容差内接近。若差异很大，应逐项核查：故障阻抗与给定
电流是否一致、CT/PT 是否一致、在线电压是否塌落、换流器是否限流、支路跳开是否形成失电岛、DER 是否跨越闭锁或跳闸
门槛。不能把差值统称为“动态误差”。

\subsection{计算成本模型}
设场景数为 $N_e$，每遍时间步数为 $N_t$，单步稀疏 DAE 成本为 $C_{DAE}$。在线三级对照主导成本约为
\begin{equation}
 C_{online}\approx4N_eN_tC_{DAE}+C_{tree},
\end{equation}
其中系数 4 来自主发现、主反馈、后备发现、后备反馈。给定轨迹模式成本约为
$C_{trace}\approx N_eN_tC_{relay}+C_{tree}$，通常远小于在线模式。

这一模型给出三个可检验预言：时域加倍且步长不变，在线耗时近似加倍；步长减半且时域不变，耗时近似加倍；场景数加倍，
动态部分近似线性增加。稀疏因子复用和 Newton 迭代数会造成偏离，但不会改变四遍 DAE 的结构常数。HTTP 限制单次最多
100 个场景，避免无界请求占用会话。

\subsection{时间步收敛检查}
设用步长 $h$ 和 $h/2$ 得到清除时刻 $t_h,t_{h/2}$ 与 EENS $E_h,E_{h/2}$。定时限动作由采样跨越触发，清除时刻差通常
满足
\begin{equation}
 |t_h-t_{h/2}|\le h+epsilon.
\end{equation}
EENS 的清除窗口差满足近似上界
\begin{equation}
 |E_h-E_{h/2}|\le\sum_e\lambda_e S_{e,max}\frac{h}{3600}
\end{equation}
，不包括因 DER 终态跨门槛导致的离散跃迁。若步长变化使 DER 从“超时跳闸”变成“恢复”，应提高时间分辨率并报告状态
门槛，而不是只看 EENS 相对误差。

\subsection{HTTP 请求契约}
可靠性主路由使用 \fld{method=protection\_cyber\_compare}。顶层
\fld{protection\_cyber.online\_dae} 选择轨迹来源。在线场景中的
\fld{online\_dae} 对象可给出故障母线、主后备支路稳定 ID、故障时刻、阻抗、时间步、终止时刻、CT/PT 和动态初始化选项。
支路或母线填零/省略时，GUI 请求允许后端按在役交流支路选择默认值，并在诊断中返回实际后果；工程批处理应显式给稳定 ID，
避免模型拓扑变化后默认对象改变。

主、后备对象仍给出继电器、断路器和按需成功概率。在线模式忽略其中人为给出的量测幅值，但保留轨迹字段以维持同一请求
schema；真正参与计算的轨迹点数由响应
\verb|online_diagnostics.*_relay_trajectory_points| 给出。GUI 通过“在线 DAE”开关、故障母线、主后备支路和 DAE 步长控件
构造该请求。

\subsection{HTTP 响应证据}
在线成功响应至少包含：
\begin{itemize}
 \item \fld{validity.online\_network\_dae\_coupled=true}；
 \item \fld{dae\_trajectories\_consumed\_by\_event\_tree=true}；
 \item 主、后备轨迹点数均大于零；
 \item 主、后备动作反馈均为真；
 \item 主、后备 DER-FRT 消费均为真；
 \item DAE 与事件树清除时刻满足式~\eqref{eq:rel-dae-event-time-consistency}；
 \item 径向开断时返回失电母线稳定 ID 和负荷退出应用标志；
 \item 静态、仅保护、联合三行指标与保护收益、信息增量。
\end{itemize}
其中任一动态证据缺失，前端把工作流标记为受限或失败，而不是仅凭 HTTP 200 显示“在线耦合完成”。

\subsection{端到端构造算例}
E2E 使用三母线系统：母线 1 为电源，母线 2、3 为负荷；主后备支路使用稳定 ID 1 和 2。故障设在母线 2，
$t_f=0.05$ s，$R_f=0.2$ pu，$\Delta t=0.01$ s。主、后备定时限分别为 0.02 s 和 0.05 s，断路器机械延时均为
0.01 s。测试先通过 JSON 装载系统，再调用 HTTP 在线模式。

验收不是硬编码某一个浮点清除时刻，而是要求：两条 DAE 轨迹非空；两条支路动作进入各自闭环；DER 监视身份相同；每个
DAE 清除时刻与事件树时刻差不超过 0.010000001 s；联合比较有效性为真。随后浏览器实际勾选“在线 DAE”，运行同一路由，
结果区必须出现“在线 DAE 与年度事件树闭环诊断”“主后备均已反馈”“主后备均已消费”。

\subsection{径向与并联对照}
专项 C++ 测试先构造两条并联支路。跳开一条后仍连通，失电母线集合为空，闭环 DAE 通过剩余支路恢复。随后删除第二条支路
形成径向系统，跳开唯一支路后失电母线集合精确为目标负荷母线，自动负荷退出被应用，闭环仍收敛。这一成对测试同时证明：
拓扑算法不会把所有跳闸都当成停电，也不会在真实失电岛中保留不可解的恒功率负荷。

\subsection{异常路径与诚实失败}
下列情况明确失败并带原因：故障母线不存在；被保护支路不存在或不在役；故障阻抗无效；时间步非正；时域早于故障或清除
时刻；断路器延时为负；第一遍或第二遍 DAE 不收敛；主后备监视器维数或稳定 ID 不一致；事件树类概率不守恒；清除时刻
超过一步容差。HTTP 将 C++ 标准异常转为 400 和中文 \verb|error|，不会返回一组全零指标。

动态库个别设备可能抛出非标准异常。在线适配器在主 DAE、后备 DAE、事件类转换和年度聚合四个边界分别捕获并转成带阶段
名称的标准错误，避免 HTTP 只显示“未知服务器错误”。这项防护不掩盖失败，只改善可诊断性。

\subsection{结果适用边界}
当前在线可靠性入口的网络量测采用正序交流网络并接受单相故障相别输入；年度后果使用用户给定三窗口，不把动态瞬时功率
直接外推为全年能量。LCC 换流器和显式三相不平衡保护动态不属于该入口的模型口径。它们在系统中出现时，结果必须通过
\fld{model\_scope} 与限制文字说明，不能解释成相应设备已被详细保护模型覆盖。

这类边界与“代码没有执行理论功能”不同：本章提出的正序在线 DAE、量测、继电、主后备、断路器、失电岛、DER-FRT 和年度
事件树均有执行入口和测试；模型没有声称覆盖的 LCC/显式三相物理不能用当前结果替代。

\subsection{验证矩阵与复现顺序}
推荐按以下顺序复现，便于定位偏差：先运行保护判据与 CT/PT 闭式单元测试；再运行在线 DAE 的并联/径向、GFL/GFM 和主后备
年度消费专项；随后构建 HTTP 服务器，运行可靠性工作流 E2E；最后在桌面和移动视口检查在线控件、结果表、控制台错误与页面
横向溢出。若数值与理论预言偏差超过一步时间误差，先检查轨迹、稳定 ID 和事件记录，再检查年度类概率，不能直接放宽断言。

实现锚点为 \srcpath{src/reliability/online\_protection.cpp:run\_online\_protection\_dae}、
\srcpath{src/reliability/online\_protection.cpp:compare\_online\_protection\_cyber\_reliability}、
\srcpath{tests/run\_gui\_server.cpp:online\_protection\_cyber\_scenario\_from\_json} 与
\srcpath{tests/e2e/reliability\_workflow\_e2e.mjs}。这些锚点分别承担动态、年度聚合、HTTP 适配和浏览器验收，缺少任一层都不能
称为产品闭环。

\section{静态、仅保护与信息物理联合方法对比}\label{sec:rel-method-comparison}

\subsection{对比的公平性条件}
三种方法必须共享同一系统指纹、故障频率、切负荷定义、修复时间和年度小时数。只有保护与信息处理层逐级增加：
\begin{longtable}{@{}L{3.0cm}L{5.0cm}L{5.8cm}@{}}
\toprule
方法 & 事件逻辑 & 后果\\
\midrule
静态 FMEA & 故障发生即按未清除/全修复处理 & $S_u\tau_r$\\
仅保护 & 主/后备继电、断路器、配合与 FRT & 类条件三窗口\\
信息物理联合 & 仅保护事件树再条件于路径、QoS、共因、供电和电池 & 联合类条件三窗口\\
\bottomrule
\end{longtable}
若三行使用不同负荷倍率或不同故障率，差值不能解释为保护或信息系统贡献。

\subsection{静态 FMEA 基线}
对场景 $k$，静态基线把未清除切负荷保持到修复：
\begin{align}
 E_k^{static}&=S_{k,u}\tau_{k,r},\\
 \mathrm{EENS}^{static}&=\sum_k\lambda_kE_k^{static},\\
 \mathrm{LOLE}^{static}&=\sum_k\lambda_k\tau_{k,r}\mathbf1[S_{k,u}>\varepsilon],\\
 \mathrm{LOLF}^{static}&=\sum_k\lambda_k\mathbf1[S_{k,u}>\varepsilon].
\end{align}
它是保守基线，不执行保护逻辑，因此不应把其中的“未清除”解释为所有保护都拒动；它表示传统单持续时间 FMEA 的后果口径。

\subsection{仅保护模型}
仅保护模型删除全部信息元件、路径、环境和功能绑定，保留主后备轨迹及按需成功概率。其指标为
\begin{equation}
 \mathrm{EENS}^{P}=\sum_k\lambda_k\sum_{c_P}\pi(c_P\mid k)E(k,c_P).
\end{equation}
本地保护不依赖通信时，这一行反映真实保护快速清除的收益。若保护方案本来依赖通道，则为公平起见，应把通道作为保护链
固有成功概率或在联合行中绑定信息功能，并明确“仅保护”是理想通道反事实。

\subsection{信息物理联合模型}
联合行保留完整信息状态：
\begin{equation}
 \mathrm{EENS}^{ICP}=\sum_k\lambda_k\sum_{x^c,c_P}
 p(x^c\mid k)p(c_P\mid x^c,k)E(k,x^c,c_P).
\end{equation}
信息系统可以改变检测、主/后备跳闸通道、隔离和恢复，因而既改变保护类概率，也改变切换时长与残余切负荷。

\subsection{两个可解释差值}
平台报告
\begin{align}
 \Delta E_{cyber}&=\mathrm{EENS}^{ICP}-\mathrm{EENS}^{P},\\
 \Delta E_{protection}&=\mathrm{EENS}^{static}-\mathrm{EENS}^{P}.
\end{align}
$\Delta E_{cyber}$ 是相对理想信息条件的退化增量；$\Delta E_{protection}$ 是相对静态未清除基线的保护收益。二者不是
互斥 Shapley 分解，在非单调后果或误动存在时可以为负。

\subsection{单调性条件}
若对每个类满足
\begin{equation}
 E_{primary}\le E_{backup}\le E_{unclear},\qquad
 E_{auto}\le E_{manual}\le E_{failed},
\end{equation}
则降低信息可用率会把概率质量移向更差类，联合 EENS 不减；静态 FMEA 与信息可用率无关。该性质进入 E2E 回归：把信息
可用率从 $0.8$ 降至 $0.4$，静态值必须不变、联合值必须上升。若后果不满足单调性，测试应改为概率守恒和逐类解释，不能
强制错误的方向性。

\subsection{解析对照算例}
当前 HTTP 基准使用一项故障频率和三窗口后果，得到：
\begin{center}
\begin{tabular}{@{}lrrr@{}}
\toprule
方法 & EENS/(MWh/年) & LOLE/(h/年) & 解释\\
\midrule
静态 FMEA & 200.000000 & 20.000000 & 未清除后果保持到修复\\
仅保护 & 0.4007778 & 0.2000778 & 主/后备快速清除\\
信息物理联合 & 21.5217333 & 4.1600622 & 信息失效使部分事件进入人工/拒清类\\
\bottomrule
\end{tabular}
\end{center}
三行使用同一故障频率与切负荷数据。仅保护相对静态显著下降证明保护逻辑真正改变后果；联合相对仅保护上升证明信息路径
状态实际进入事件树，而非只显示一个可用率字段。

\subsection{为何不能直接与普通 FMEA 数值混比}
通用 \srcpath{run\_distribution\_fmea} 会从完整网络求解阶段切负荷，而三级对照入口消费用户提供的三窗口切负荷。
只有当三窗口后果由同一网络、同一故障和同一负荷倍率自动生成时，二者绝对值才可互比。当前 GUI 对照的目标是比较三种事件
逻辑，不宣称复现通用 FMEA 的每一网络后果。

\subsection{保护成功率敏感性}
主清除概率 $p_p=p_Rp_B$。在其余条件固定且后果单调时，主继电成功率导数为
\begin{equation}
 \frac{\partial\mathbb E[E]}{\partial p_R}
 =p_B(E_p-E_{fallback,B})+(1-p_B)(E_{fallback,R}-E_{fallback,B}),
\end{equation}
符号由后果差决定。提高成功率只有在主清除后果更优时才有正价值；这说明概率参数和区变异后果必须一起校准。

\subsection{配合裕度敏感性}
配合裕度不直接改变事件概率，而改变 \fld{coordination\_satisfied} 诊断。若应用策略规定裕度不足时闭锁自动重合或扩大隔离，
调用方应把诊断映射到不同后果类；当前对照保留实际清除结果，不把单纯裕度告警自动升级为停电。

\subsection{通信可用率敏感性}
单路径串联系统中，功能可用率近似为元件成功概率乘积；改善最薄弱串联元件的边际收益最大。冗余路径中，新增一条完全独立
路径的收益为现有全部路径失败概率乘新增路径成功率；若共享电源，收益显著降低。平台精确元件位枚举能够反映该差异。

\subsection{QoS 敏感性}
提高固有可用率不会挽救超过时延上限的路径。路径先经过确定性 QoS 门限，只有合格路径再参与概率枚举。因此优化通信网络时
应分别考察“路径是否满足实时契约”和“合格路径的可用率”，不能用单一九个九指标掩盖保护时延超限。

\subsection{共因敏感性}
两条名义冗余路径若共享共因组，其联合失效概率至少包含共因环境概率。即使每个独立元件可用率趋近一，功能不可用率也不会
低于该共因概率。联合对照中降低独立失效率但保持共因不变会出现收益饱和，这是合理现象而非枚举误差。

\subsection{供电和电池敏感性}
备用能量对持续时间呈阈值效应：$E_b/P_b\ge T_{outage}$ 时静态环境中节点保持可用，略低于阈值时整段视为失效；序贯模型
则在耗尽时刻发生状态切换。阈值附近应使用序贯事件队列或对停电持续时间积分，避免静态二值近似造成突跳。

\subsection{检测混淆敏感性}
提高总检出率但把概率从正确类别移到错误类别，可能扩大隔离区并增加 EENS。评估应同时报告 $P_D$、正确分类率和每类隔离
后果。只优化召回率会忽略错误分类的系统代价。

\subsection{恢复风险约束敏感性}
收紧最大不安全概率会删除高恢复量动作，短期 EENS 可能上升，但避免不安全带电和二次事故。可靠性和安全性不是同一目标；
平台把风险约束作为硬门槛，不能用较低 EENS 证明越过安全门槛的策略可接受。

\subsection{在线 DAE 与给定轨迹对比}
给定轨迹模式便于复现实验和批量统计，但轨迹不会因整定、换流器控制或支路跳闸反馈自动改变。在线 DAE 模式从系统状态生成
轨迹，并把清除动作反馈到第二次动态运行，能揭示 DER 限流导致的保护延时或 FRT 跳闸。两者的结果差异应解释为轨迹生成
保真度差异，而非随机误差。

\subsection{年度序贯与解析事件树对比}
解析事件树直接对互斥类求期望，无抽样误差，适合独立代表事件。序贯队列能处理事件重叠、电池跨事件状态和共因持续过程，
但有抽样误差。低事件率、每次事件前状态重置时，两者均值应一致；有记忆和重叠时，序贯结果是更合适的口径。

\subsection{验收矩阵}
\begingroup\footnotesize
\begin{longtable}{@{}L{4.1cm}L{5.1cm}L{5.1cm}@{}}
\toprule
验证对象 & 必须成立的关系 & 失败含义\\
\midrule
\endfirsthead
\toprule
验证对象 & 必须成立的关系 & 失败含义\\
\midrule
\endhead
\bottomrule
\endfoot
事件树概率 & 正概率类和为一 & 分支遗漏或重复\\
三级静态基线 & 不随信息可用率变化 & 比较基准被污染\\
联合后果单调算例 & 信息退化时 EENS 不减 & 功能绑定或后果选择错误\\
在线清除时间 & DAE 与事件树差不超过时间步 & 轨迹转换不一致\\
DER 轨迹身份 & 主/后备/未清除稳定 ID 一致 & 错位消费设备轨迹\\
重要抽样恒等 & $\kappa=1,w=1,N_{eff}=N$ & 似然比错误\\
精确枚举 & 状态概率和为一 & 退化元件或位映射错误\\
Birnbaum & 与有限差分一致 & 边际状态配对错误\\
保护误动分解 & EENS 分解残差为零 & 增量未进入总量或重复计数\\
区间并集 & LOLE 不超过各事件时长和 & 重叠停电重复计算\\
\end{longtable}
\endgroup

\subsection{结论使用规则}
若三级结果同时给出，应优先报告三行绝对值、两个差值和全部有效性标志。不能只展示“联合方法更差”而不说明这是对理想信息
的增量，也不能只展示“保护收益”而隐藏保护误动。所有百分比变化都应注明分母；静态基线接近零时不报告无意义的巨大百分比。

% ── 实现转写层（源码等价）──
\section{源码等价数学实现}\label{sec:rel-source-model}

本章按入口分别转写当前代码。不同入口对基线切负荷、时间尺度和事件频率的处理不同，公式不得
跨入口复用。所有路径均位于 \srcpath{src/reliability/}。

\subsection{统一参数解析器的实际优先级}

\srcpath{resolve\_reliability\_params} 先将非有限或非正数视为“未提供”；
\fld{forced\_outage\_rate} 仅在 $0<f<1$ 有效，按需失败概率仅在 $0<p_d\le1$ 有效。
报告小时数 $H$ 仅在策略值有限且为正时采用，否则固定回退 8760。已知年故障强度
$\lambda$ 与修复时长 $r$ 时，公共收尾函数执行
\begin{equation}
 \mu=H/r,\qquad U=\begin{cases}\lambda/(\lambda+\mu),&r>0,\\0,&r\le0,
 \end{cases}\qquad \texttt{mttf\_hr}=\begin{cases}H/\lambda,&\lambda>0,\\0,&\text{否则}.
 \end{cases}
\end{equation}
实际解析顺序为：
\begin{enumerate}
 \item 有效的主动按需数据：$\lambda=\nu_dp_d$，修复时长优先
 \fld{cyber\_recovery\_hr}，其次 \fld{mttr\_hr}，再次 \fld{mttr\_hours}；
 \item \fld{failure\_rate\_per\_year}+\fld{mttr\_hr}；
 \item 显式 \fld{mttf\_hours}，有修复时额外覆盖 $U=r/(\texttt{mttf\_hours}+r)$；
 \item \fld{mtbf\_hours}，循环时长约定下先令
 $\texttt{mttf}=\max(\texttt{mtbf\_hours}-\texttt{mttr\_hours},0)$，否则直接使用 MTBF；
 \item FOR 与修复时长：$U=f$、$\lambda=fH/((1-f)r)$、MTTF$=r(1-f)/f$；
 \item 仅 FOR，或仅故障率；缺失量保持未确定并记录警告；
 \item 策略允许且模板的 $\lambda,r$ 均为正时，只填补缺失项，再执行公共收尾函数。
\end{enumerate}
依据为 \srcpath{src/reliability/reliability\_assessment.cpp:resolve\_reliability\_params}。这是一组控制流，不可改写为无优先级的
对称参数换算表。

\subsection{非序贯 MC}

每个有效元件独立比较均匀随机数与解析不可用度：
$X_c=[U(0,1)<U_c]$；基态已停运元件的 $U_c$ 被置零，严格数据策略下缺少算例可靠性数据的
$U_c$ 也被置零。依据为 \srcpath{src/reliability/reliability\_assessment.cpp:sample\_state} 和
\srcpath{src/reliability/reliability\_assessment.cpp:run\_nonsequential\_mc}。

对第 $q$ 个状态，令 $D_q$ 为后果引擎切负荷，$D_0$ 为预计算健康态切负荷，源码累加
\begin{align}
 D_q^{inc}&=\max(0,D_q-D_0),\\
 \texttt{edns}&=\frac1n\sum_qD_q,&
 \texttt{eens}&=8760\,\texttt{edns},\\
 \texttt{incremental\_edns}&=\frac1n\sum_qD_q^{inc},&
 \texttt{incremental\_eens}&=8760\,\texttt{incremental\_edns},\\
 \texttt{plc}&=\frac1n\sum_q[D_q^{inc}>\varepsilon],&
 \texttt{lole}&=8760\,\texttt{plc}.
\end{align}
注意原始 \fld{eens\_mwh\_yr} 使用 $D_q$，而失负荷判断、节点 EENS、收敛历史与尾部样本使用
$D_q^{inc}$。方差按总体二阶矩
$v=n^{-1}\sum(D_q^{inc})^2-(\overline D^{inc})^2$，仅在
$n>1$、均值正且 $v>0$ 时令
$\texttt{cov}=\sqrt v/(\overline D^{inc}\sqrt n)$。只有
$n>100$、CoV$>0$ 且小于阈值才提前收敛。
依据为 \srcpath{src/reliability/reliability\_assessment.cpp:run\_nonsequential\_mc} 和
\srcpath{src/reliability/reliability\_assessment.cpp:eval\_missing（lambda）}。

尾部风险不是直接对单小时乘 8760。源码以有放回方式，每个合成年抽取 8760 个
$D_q^{inc}$，合成年数为
$\min(\max(200,n),2000)$，再把年 EENS/LOLE 传给尾部函数；见
\srcpath{src/reliability/reliability\_assessment.cpp:eval\_missing（lambda）}。

\subsection{序贯 MC}

严格数据策略下缺数据元件的 MTTF 被置为 $10^{15}$ 小时。每个在运元件从上态开始，源码用
\begin{equation}
 T_{up}=-\texttt{mttf}[c]\ln u,qquad
 T_{down}=-\texttt{mttr}[c]\ln u,qquad u\sim U(0,1)
\end{equation}
交替填充小时状态矩阵，见 \srcpath{src/reliability/reliability\_assessment.cpp:run\_sequential\_mc} 和
\srcpath{src/reliability/reliability\_assessment.cpp:n0\_for\_hour（lambda）}。每一失负荷小时直接累加该小时后果引擎返回的
\fld{curtailment\_mw}；代码没有在序贯年汇总中扣除对应小时的 N-0 切负荷。对每年得到
$E_y$、失负荷小时 $L_y$，并由布尔失负荷序列的 false$\to$true 回合数得到 $F_y$。最终
\begin{equation}
 \texttt{eens}=N_y^{-1}\sum_yE_y,\quad
 \texttt{lole}=N_y^{-1}\sum_yL_y,\quad
 \texttt{lolf}=N_y^{-1}\sum_yF_y,
\end{equation}
$\texttt{edns}=\texttt{eens}/\fld{hours\_per\_year}$，
$\texttt{plc}=\texttt{lole}/\fld{hours\_per\_year}$。CoV 使用年 EENS 的总体二阶矩，收敛条件与
非序贯相同。依据为 \srcpath{src/reliability/reliability\_assessment.cpp:eval\_down\_hours（lambda）}。

\subsection{元件 FMEA 与信息物理混合}

每个元件的自动化和人工路径先分别生成阶段 ENS 与失负荷持续时间。未启用信息物理时
\begin{equation}
 EENS_k=\lambda_k\,\texttt{automatic.ens\_mwh}.
\end{equation}
启用后，令 $p_{down}=1-p_{auto}$，源码分解为
\begin{align}
 E_{perfect}&=\lambda E_{auto},\\
 \Delta E_{duration}&=\lambda p_{down}(E_{durationOnly}-E_{auto}),\\
 \Delta E_{control}&=\lambda p_{down}(E_{manual}-E_{durationOnly}),\\
 EENS_k&=E_{perfect}+\Delta E_{duration}+\Delta E_{control},\\
 LOLE_k&=\lambda\left(p_{auto}T_{auto}+p_{down}T_{manual}\right).
\end{align}
失负荷频率为
$\lambda(p_{auto}[auto\ causes\ loss]+p_{down}[manual\ causes\ loss])$。
系统值直接累加逐元件值。依据为
\srcpath{src/reliability/reliability\_assessment.cpp:evaluate\_class（lambda）} 和
\srcpath{src/reliability/reliability\_assessment.cpp:nodal\_value（lambda）}。因此旧式
$\lambda(ENS_{sw}+ENS_{rep})$ 只有在 \texttt{automatic.ens\_mwh} 恰按该方式构成时才成立，
不能替代当前信息物理分支。

\subsection{故障模式 FMEA 与二阶共故障}

单模式持续时间与频率直接赋为
\begin{align}
 d_m&=\texttt{isolation\_hr}+\texttt{switching\_hr}
 +\max(\texttt{params.repair\_hr},\texttt{mode.repair\_hr}),\\
 f_m&=\texttt{params.lambda\_per\_year},\\
 EENS_m&=f_md_mS_m,\qquad LOLE_m=[S_m>\varepsilon]f_md_m.
\end{align}
主动模式的 $f_m=\nu_dp_d$ 已在解析器中完成；当前评估器没有再乘
\fld{probability\_given\_initiated}。不支持的后果补丁被记录但贡献保持零。
依据为 \srcpath{src/reliability/failure\_mode.cpp:evaluate\_single\_mode（lambda）}。

二阶入口仅缓存不同组件、$f_m>0$ 且 $d_m>0$ 的支持模式，并令
\begin{align}
 U_m&=\min(1,f_md_m/8760),& U_{ij}&=U_iU_j,\\
 EENS_{ij}&=U_{ij}\,8760\,S_{ij},& LOLE_{ij}&=U_{ij}\,8760,\\
 f_{ij}&=f_if_j(d_i+d_j)/8760,&
 d_{ij}&=\begin{cases}d_id_j/(d_i+d_j),&d_i+d_j>0,\\0,&\text{否则}.
 \end{cases}
\end{align}
同组件模式、$U_{ij}$ 小于阈值、硬冲突补丁或零切负荷对均不贡献；枚举受
\fld{max\_pairs\_evaluated} 限制。依据为 \srcpath{src/reliability/failure\_mode.cpp:evaluate\_single\_mode（lambda）} 和
\srcpath{src/reliability/failure\_mode.cpp:evaluate\_single\_mode（lambda）}。

\subsection{三阶段结果聚合}

每个故障先对三个阶段的节点切负荷逐点扣除健康态并截为非负
$p^{raw}_{ki}=\max(0,\texttt{shed\_by\_load}^{fault}_{ki}-\texttt{shed\_by\_load}^{healthy}_{ki})$，
再按阶段总量缺口重标定：令
$\Delta=\max(0,\texttt{shed\_kw}^{fault}-\texttt{shed\_kw}^{healthy})$，
\begin{equation}
 p^{inc}_{ki}=p^{raw}_{ki}\cdot\frac{\Delta}{\sum_j p^{raw}_{kj}}
 \quad\left(\sum_j p^{raw}_{kj}\le10^{-12}\ \text{时整段置零}\right),
\end{equation}
汇总为 $P_1,P_2,P_3$（kW）。依据为
\srcpath{src/reliability/three\_stage\_reliability.cpp:incremental\_shed（lambda）}。源码逐故障赋值
\begin{align}
 \texttt{ens\_kwh}&=P_1\tau_{iso}+P_2\tau_{sw}+P_3\tau_{rep},\\
 d&=[P_1>10^{-6}]\tau_{iso}+[P_2>10^{-6}]\tau_{sw}
 +[P_3>10^{-6}]\tau_{rep},\\
 \texttt{eens\_contribution\_mwh\_yr}&=\lambda\,
 \texttt{ens\_kwh}/1000,\\
 \texttt{lole\_contribution\_hr\_yr}&=\lambda d,\qquad
 \texttt{lolf\_contribution\_occ\_yr}=[d>0]\lambda.
\end{align}
系统 \fld{eens\_kwh\_yr} 不是简单累加逐故障字段后再乘一次：代码按负荷点直接累加
\begin{equation}
 \sum_{k,i}\lambda_k
 (p_{ki1}^{inc}\tau_{k,iso}+p_{ki2}^{inc}\tau_{k,sw}+p_{ki3}^{inc}\tau_{k,rep}).
\end{equation}
SAIFI/SAIDI 只要某负荷点任一阶段增量超过 $10^{-6}$ 就将该故障频率计入中断次数，并按各阶段
阈值分别累计持续时间。依据为
\srcpath{src/reliability/three\_stage\_reliability.cpp:shed\_sum（lambda）} 和
\srcpath{src/reliability/three\_stage\_reliability.cpp:append\_isolation\_action（lambda）}。

\subsection{频率--持续时间递推}

该入口只读取在运 AC 发电机，容量为 \fld{pmax\_mw}。MTTR 缺失时固定回退 50 小时；有效 FOR
$q\in(0,1)$ 时 $p=1-q$、$\mu=1/r$、$\lambda=q/((1-q)r)$，否则按完美元件
$p=1,q=0,\lambda=0,\mu=1$。容量离散步长固定为 10 MW，容量步数用整数截断
$c=\lfloor C/10\rfloor$；$c\le0$ 的机组不进入递推。

数组初值为 $P_0[0]=1$、其余 $P_0,F_0$ 为零。代码存储的是累积数组，并逐级执行
\begin{align}
 P_{new}[x]&=pP[x]+qP[x-c],\\
 F_{new}[x]&=pF[x]+qF[x-c]+\lambda p(P[x-c]-P[x]),
\end{align}
其中 $x<c$ 时硬编码 $P[x-c]=1,F[x-c]=0$。备用索引先取
$r=\operatorname{clamp}(\lfloor(C_{tot}-L)/10\rfloor,0,n-1)$，再用 $r+1$ 读取
LOLP 与频率：
\begin{equation}
 LOLP=P[r+1],\quad LOLF=8760F[r+1],\quad LOLE=8760LOLP,\quad
 LOLD=\begin{cases}LOLE/LOLF,&LOLF>0,\\0,&\text{否则}.
 \end{cases}
\end{equation}
越过数组上界时 LOLP/LOLF 均为零。依据为
\srcpath{src/reliability/reliability\_assessment.cpp:run\_frequency\_duration\_analysis}。

\subsection{VaR、CVaR 与配电指标}

尾部函数以 EENS 样本数 $n$ 为基准；LOLE 长度不等时 resize 到 $n$ 并以零补齐。排序后
$i_{VaR}=\min(\lceil\alpha n\rceil-1,n-1)$，CVaR 是从该索引到末尾的算术平均，包含 VaR
样本本身。五个报告分位数使用
$\min(\lfloor pn\rfloor,n-1)$，而不是同一个 ceil 规则。依据为
\srcpath{src/reliability/reliability\_assessment.cpp:compute\_tail\_risk}。

配电指标的客户数优先取 AC 负荷 \fld{n\_customers}，否则取
$\max(1,10p_{MW})$；没有显式负荷但有 \fld{pd\_mw} 的 AC 母线同样回退。DC 仅在 CIF/CID
向量长度超过 AC 母线数时加入，依次优先母线客户数、负荷客户数、
$\max(1,10(\sum|p_{load}|+\max(0,pd)))$。客户总数小于 1 时返回零结构。
具体权重与索引布局见 \srcpath{src/reliability/reliability\_assessment.cpp:compute\_distribution\_indices}。

\subsection{覆盖边界}

本章只保证所列聚合式与当前代码一致。切负荷后果引擎内部的 AC-only DC-OPF、混合 AC/DC LP
和三阶段恢复 MILP 是不同模型；其覆盖范围由结果的 \fld{model\_scope}、
\fld{model\_limitations} 与有效性标志限定，不能仅凭“EENS”字段名互相替代。

% 第 8 章：选项与结果数据结构
\section{选项与结果数据结构}\label{sec:rel-io}

\subsection{蒙特卡洛选项 \texttt{ReliabilityOptions}}
\compmeta{ReliabilityOptions}{include/hacdcpf/reliability/reliability\_assessment.hpp:ReliabilityOptions}
\begin{paramtable}
\fld{max\_iterations} & $n_{max}$ & — & $10^3\sim10^5$ & 10000 & 非序贯采样数上限 / 序贯年数上限\\
\fld{cov\_threshold} & — & — & $0.01\sim0.1$ & 0.05 & 变异系数收敛阈值\\
\fld{curtail\_threshold\_mw} & $\varepsilon$ & MW & 小正数 & 0.01 & 判定失负荷事件的切负荷阈值\\
\fld{seed} & — & — & 整数 & 0 & 随机种子（0=非确定）\\
\fld{load\_scale\_factor} & $\rho^\star$ & — & $\ge1$ & 1.0 & 负荷加压系数（$>1$ 降裕度）\\
\fld{hours\_per\_year} & $H$ & 小时/年 & 8736/8760 & 8736 & 序贯年报告小时数\\
\fld{data\_policy} & — & — & 枚举 & 见 \S\ref{sec:rel-param} & 缺失数据处理策略\\
\fld{opf\_options} & — & — & 结构 & 见 OPF 手册 & 状态评估 DC-OPF 选项\\
\fld{compute\_tail\_risk} & — & — & 布尔 & true & 是否计算 VaR/CVaR\\
\fld{var\_confidence} & — & — & $0\sim1$ & 0.95 & VaR/CVaR 置信水平\\
\fld{compute\_distribution\_indices} & — & — & 布尔 & false & 是否算 SAIFI/SAIDI/ASAI\\
\fld{use\_importance\_sampling} & — & — & 布尔 & false & 重要抽样方差缩减（实验）\\
\fld{importance\_lambda} & — & — & $>0$ & 2.0 & 重要抽样扭曲因子\\
\fld{enable\_parallel} & — & — & 布尔 & true & 独立状态/年并行\\
\fld{parallel\_threads} & — & — & 整数 & 0 & 0=硬件并发数\\
\end{paramtable}

\subsection{元件 FMEA 选项 \texttt{FMEAOptions}}
\compmeta{FMEAOptions}{include/hacdcpf/reliability/reliability\_assessment.hpp:FMEAOptions}
\begin{paramtable}
\fld{load\_scale\_factor} & $\rho^\star$ & — & $\ge1$ & 1.0 & 负荷加压系数\\
\fld{switching\_time\_hr} & $\tau_{sw}$ & 小时 & $0.1\sim2$ & 0.5 & 默认开关/隔离时间\\
\fld{voll} & — & \$/MWh & — & 10000 & 失负荷价值（仅破平）\\
\fld{enable\_microgrid\_islanding} & — & — & 布尔 & true & 隔离后按孤岛评估\\
\fld{enable\_repair\_reconfiguration} & — & — & 布尔 & true & 修复阶段提升本地 DER\\
\fld{enable\_switch\_reconfiguration} & — & — & 布尔 & true & 修复阶段显式开关重构\\
\fld{max\_repair\_switch\_actions} & — & — & 整数 & 2 & 同时考虑的开关动作数上限\\
\fld{max\_repair\_opf\_calls} & — & — & 整数 & 200 & 修复搜索 OPF 预算（0=不限）\\
\fld{enable\_storage\_dispatch} & — & — & 布尔 & true & 储能作可调急救源\\
\fld{enable\_grid\_forming\_vsc\_support} & — & — & 布尔 & true & 构网 VSC 锚定孤岛\\
\fld{data\_policy} & — & — & 枚举 & 补缺 & 缺失数据策略\\
\fld{cyber\_physical} & — & — & 结构 & disabled & 信息物理 L1 条件\\
\end{paramtable}

\subsection{故障模式 FMEA 选项 \texttt{FailureModeFMEAOptions}}
\compmeta{FailureModeFMEAOptions}{include/hacdcpf/reliability/failure\_mode.hpp:FailureModeFMEAOptions}
\begin{paramtable}
\fld{max\_order} & $k$ & — & 1/2 & 1 & 最高同时故障阶数（2=N-2）\\
\fld{min\_pair\_unavailability} & — & — & 小正数 & $10^{-10}$ & 舍弃可忽略共故障对的阈值\\
\fld{max\_pairs\_evaluated} & — & — & 整数 & 50000 & 对枚举硬上限\\
\fld{curtail\_threshold\_mw} & $\varepsilon$ & MW & 小正数 & 0.01 & 失负荷阈值\\
\fld{catalog} & — & — & 结构 & — & 目录过滤（激活/成因/元件类）\\
\fld{capabilities} & — & — & 结构 & — & 后果引擎能力声明\\
\fld{enable\_parallel} & — & — & 布尔 & true & 独立模式/对并行\\
\end{paramtable}

\subsection{三阶段选项 \texttt{ThreeStageReliabilityOptions}}
\compmeta{ThreeStageReliabilityOptions}{include/hacdcpf/analysis/three\_stage\_reliability.hpp:ThreeStageReliabilityOptions}
\begin{paramtable}
\fld{max\_switch\_operations} & — & — & $1\sim5$/不限 & INT\_MAX & 阶段 2 允许闭合联络开关数\\
\fld{unavailable\_tie\_switch\_ids} & — & — & id 列表 & 空 & 拒合联络（评估拒合影响）\\
\fld{include\_dc\_power\_flow} & — & — & 布尔 & true & 同阶段纳入 DC LinDistFlow/DER\\
\fld{include\_generator\_faults} & — & — & 布尔 & false & 纳入发电机停运\\
\fld{include\_transformer\_faults} & — & — & 布尔 & false & 纳入双绕组变压器停运\\
\fld{include\_converter\_faults} & — & — & 布尔 & false & 纳入 VSC/DC-DC 停运\\
\fld{include\_switch\_faults} & — & — & 布尔 & false & 纳入开关/断路器停运\\
\fld{revalidate\_stage3\_plan} & — & — & 布尔 & false & 固定阶段 2 二元量重解阶段 3\\
\fld{reliability\_configuration} & — & — & 结构 & 空 & 会话保护定义\\
\end{paramtable}

\subsection{结果结构摘要}
\begin{itemize}
  \item \fld{ReliabilityResult}（:ThreeStageReliabilityOptions）：\fld{eens\_mwh\_yr}、\fld{lole\_hr\_yr}、
        \fld{lolf\_occ\_yr}（序贯）、\fld{plc}、\fld{nodal\_eens\_mwh\_yr}、
        \fld{critical\_components}（共现归因秩）、\fld{tail\_risk}、\fld{distribution\_idx}、
        \fld{model\_scope}/\fld{validity}/\fld{data\_quality}；
  \item \fld{FMEAResult}（:run\_three\_stage\_reliability\_from\_string）：系统级指标 + 逐条 \fld{contingencies}（按 EENS 贡献降序）
        + \fld{cyber\_physical} + \fld{ValidityFlags}（15 个能力标志，:FMEAResult::ValidityFlags）；
  \item \fld{FailureModeFMEAResult}（\srcpath{include/hacdcpf/reliability/failure\_mode.hpp:FailureModeFMEAResult}）：\fld{contingencies} 与
        \fld{co\_contingencies}（N-2）、\fld{coverage}、\fld{n\_pairs\_evaluated}；
  \item \fld{ThreeStageReliabilityResult}（:FailureConsequenceKind）：\fld{saifi}/\fld{saidi\_min}/\fld{eens\_kwh\_yr}、
        逐故障 \fld{faults}（含 \fld{switching\_sequence}）、网络计数、\fld{validity}。
\end{itemize}

\section{运行时接口、GUI 契约与验证证据}\label{sec:rel-runtime-contract}

\subsection{统一可靠性入口}
生产 GUI 使用 \texttt{POST /api/session/run\_reliability}。请求至少给出 \fld{method}，不同方法只消费各自命名对象：
\begin{center}\footnotesize
\begin{tabular}{@{}L{4.0cm}L{5.0cm}L{5.2cm}@{}}
\toprule
方法值 & 主要请求对象 & 核心输出\\
\midrule
\texttt{nsq} & \fld{monte\_carlo}、\fld{load} & EENS、LOLE、PLC、重要抽样诊断\\
\texttt{seq} & \fld{monte\_carlo}、\fld{load} & 年度 EENS/LOLE/LOLF、尾部风险\\
\texttt{fmea} & \fld{dimensions}、\fld{restoration} & 阶段后果、信息物理分解\\
\texttt{exact\_sensitivity} & \fld{exact\_sensitivity} & 精确状态、Birnbaum、FV\\
\fld{protection\_cyber\_compare} & \fld{protection\_cyber.scenarios} & 静态/仅保护/联合三行对比\\
\texttt{three\_stage} & 三阶段恢复对象 & 交直流恢复 MILP 指标\\
\bottomrule
\end{tabular}
\end{center}
未知方法、类型错误、越界概率或缺必需字段返回 HTTP 400，不自动改用 FMEA 或普通 MC。

\subsection{重要抽样请求}
非序贯蒙特卡洛对象新增：
\begin{verbatim}
"monte_carlo": {
  "max_samples": 400,
  "min_samples": 400,
  "use_importance_sampling": true,
  "importance_lambda": 4.0
}
\end{verbatim}
扭曲因子必须有限且大于零。序贯方法收到重要抽样选项会拒绝，因为静态似然比不适用于 CTMC 路径。响应的
\fld{importance\_sampling} 给出是否使用、扭曲因子、有效样本量和平均似然比。

\subsection{精确灵敏度请求}
\begin{verbatim}
"method": "exact_sensitivity",
"exact_sensitivity": {
  "maximum_components": 20,
  "curtailment_threshold_mw": 0.01
}
\end{verbatim}
上限取 $1$ 到 $20$。这里的元件数是 $0<U<1$ 的随机元件数；确定运行/停运元件仍参与后果状态。响应逐行给出
\fld{component\_type}、\fld{component\_position}、\fld{component\_index}、\fld{stable\_id}、
\fld{birnbaum\_eens}、\fld{eens\_derivative} 和 \fld{fussell\_vesely}。

\subsection{三级对照请求结构}
每个场景必须给出场景 ID、年故障频率、主/后备保护链、配合裕度、未清除终端时间、信息元件、信息功能、环境、功能绑定和
三窗口后果。服务器限制最多一百个场景和二十个信息元件，防止指数状态枚举失控。

主/后备保护链包含继电特性、逐点测量轨迹、断路器四段延时和按需成功概率。当前 GUI 提供定时限过流的常用参数；HTTP
解析器还支持反时限、距离和差动模型。测量轨迹时间必须递增，电流、视在阻抗、差流和制动流必须有限。

\subsection{信息元件与路径 JSON}
信息元件使用数组位置构造路径，但响应和诊断保留字符串 ID。每个元件可给固有可用率、包交付概率、时延、抖动、供电母线、
备用能量、功耗和共因组。路径中的索引越界立即失败，不删除非法路径后继续计算。

信息功能包含替代路径和三项 QoS 门限。空路径表示不需要任何信息元件、功能恒可用；没有任何替代路径表示功能不可用。两者
语义相反，解析器不得把空数组与含一个空路径混为一谈。

\subsection{环境 JSON}
环境数组是互斥状态，概率和必须为一。每行给环境名、失败共因组、失电母线和信息停电持续时间。交流和直流供电母线在当前
保护信息接口中使用字符串稳定标识，调用方应包含域前缀，避免 AC/DC 同号混淆。

\subsection{三窗口后果 JSON}
后果对象分别给主清除、后备清除、未清除、隔离成功、隔离失败、恢复成功和恢复失败的切负荷，以及自动恢复、人工恢复和修复
时间。所有功率和时间必须非负。服务器不根据名称猜单位：清除时间用秒，恢复/修复时间用小时，响应统一输出 MWh/年和小时/年。

\subsection{三级对照响应}
\fld{method\_comparison} 下有 \fld{static\_fmea}、\fld{protection\_only} 和
\fld{cyber\_conditioned} 三个对象，每个给 EENS、LOLE、LOLF；另给保护收益和信息增量。顶层 \fld{metrics} 使用联合行，
避免 GUI 图表和下载表使用不同主结果。

\subsection{有效性标志}
三级响应逐项报告继电逻辑、保护配合、断路器失败、信息拓扑、QoS、共因、信息节点物理供电、DER 穿越和在线 DAE。给定轨迹
HTTP 入口的 \fld{online\_network\_dae\_coupled=false} 是事实边界；C++ 在线耦合接口成功消费 DAE 轨迹后才置真。

\subsection{保护误动配置}
通用 FMEA 的 \fld{dimensions.information.protection\_misoperation} 接收：启用开关、无故障判别窗数、每窗误跳概率、
通道成功率、断路器成功率、断开负荷和恢复时长。响应返回误动频率、EENS/LOLE/LOLF 增量和 EENS 分解残差。该对象未启用时
增量自然为零；启用后该计算项实际进入系统总指标。

\subsection{GUI 方法选择}
可靠性面板显式列出“精确灵敏度（状态枚举）”和“保护-信息物理三级对照”。重要抽样开关和扭曲因子只在 NSQ 显示；精确
方法只显示随机元件上限；三级对照显示保护、信息和三窗口后果参数。切换方法时隐藏控件同时禁用，避免隐藏字段被意外提交。

\subsection{GUI 三级结果}
结果区用同一表展示三方法 EENS、LOLE、LOLF，并另列保护收益和信息失效增量。精确灵敏度用稳定 ID、类型和位置展示各元件
Birnbaum、EENS 导数与 FV。重要抽样显示扭曲因子、ESS 和平均似然比，使用户能判断方差缩减是否退化。

\subsection{旧 Level-1 与新 Level-2 的区分}
GUI 仍保留通用 FMEA 的 Level-1 智能独立因子筛查，它把检测、隔离、恢复和保护成功概率相乘或消费显式联合类，后果压缩为
自动/人工两类。\texttt{protection\_cyber\_compare} 是 Level-2 离散事件树，显式计算继电轨迹、主后备、断路器和信息元件
状态。两者不能都标成“完整信息物理仿真”，结果页必须显示模型范围。

\subsection{HTTP 错误契约}
下列情况必须返回 400：概率不在 $[0,1]$、概率行不归一、负时长/功率、非递增轨迹、未知继电特性、功能绑定越界、路径索引
越界、状态元件超限、精确随机元件超限、重要抽样因子非正。运行时求解失败也返回明确错误文本，不输出全零结果冒充成功。

\subsection{稳定 ID 契约}
对外资产身份使用组件 \fld{.index} 或带类型前缀的稳定 ID；\fld{component\_position} 仅作容器诊断。保护自动拓扑入口按稳定
ID 查找开关/支路，精确灵敏度同时返回稳定 ID 和位置。任何 JSON 客户端不得把表格行号回写为组件 ID。

\subsection{注册单元测试}
\srcpath{test\_reliability\_resolver} 覆盖参数解析、MC、重要抽样、精确灵敏度、FMEA、保护逻辑、事件树、拓扑后果和误动分解。
\srcpath{test\_intelligent\_cyber\_physical\_reliability} 覆盖混淆矩阵、风险隔离、风险恢复、POMDP、序贯事件队列、在线 DAE 和
DAE 到年度事件树的消费。\srcpath{test\_three\_stage\_reliability} 覆盖恢复 MILP。

\subsection{重要抽样回归证据}
$\kappa=1$ 的四百样本测试要求平均似然比精确为一、ESS 精确为四百；$\kappa=4$ 要求返回有限正 ESS 且不超过样本数；非法
$\kappa=0$ 必须抛出参数错误；序贯方法启用重要抽样必须拒绝。这四项同时防止“参数被解析但抽样分布未改变”的假实现。

\subsection{精确灵敏度回归证据}
混合 AC/DC 小系统含两个随机元件、一个确定元件，预期只枚举四个状态。VSC 稳定 ID 设为 $77$，验证不会返回位置 $0$。
五兆瓦直流负荷解析值为 EENS $4380$、Birnbaum $43800$、FV $1$，误差容差由双精度求和决定。

\subsection{保护信息事件树回归证据}
测试构造共享信息路径、QoS 不合格路径、共因环境和供电电池，检查类概率归一及功能可用率；再把信息可用率从 $0.8$ 降至
$0.4$，要求静态 FMEA 不变且联合 EENS 增加。主/后备断路器概率分别改变时，对应失效原因类概率按解析乘积改变。

\subsection{在线 DAE 回归证据}
两母线双回线系统在 $0.05$ 秒施加故障，定时限继电器从 DAE 电流轨迹动作并跳开稳定 ID 为 $101$ 的支路。测试检查发现和
闭环两次求解成功、故障清除与支路跳闸事件均实际施加、DER 状态机消费轨迹。主后备年度耦合再检查两条清除时刻一致性和
联合指标有效性。

\subsection{GUI E2E}
\srcpath{reliability\_workflow\_e2e.mjs} 通过真实 HTTP 服务器装载算例，执行三级对照、精确灵敏度和重要抽样，并检查方法选项、
控件显隐和返回字段。E2E 不用前端模拟响应，因此能发现服务器分支未接线、字段命名漂移和 JSON 数值类型错误。

\subsection{浏览器视觉验收}
桌面与 $390\times844$ 移动视口需检查：三级控制组不横向溢出，标签和输入不重叠，表格可在局部容器滚动，方法切换不引起
错误控件残留，控制台无异常。视觉检查不能替代接口测试，但能发现功能虽然可调用却无法在目标视口操作的问题。

\subsection{构建与依赖口径}
可靠性代码属于静态库 \srcpath{hacdcpf}，GUI 服务器是测试目录下的独立生产路由实现。完整构建依赖兄弟仓库 MIPSolvers。
若本地 MIPSolvers HEAD 与仓库记录版本不同，测试只能标为本地非可复现验证，不能把通过结果写成锁定依赖基线。

\subsection{发布前零残留审计}
发布前扫描可靠性手册和源码中的悬空状态词。每个命中必须属于明确的模型适用边界，而不能是手册承诺但代码没有执行的功能；
源码新增行不得保留任务标记、空壳函数或固定假值。

\subsection{页数和语言验收}
主手册必须由中文理论、公式推导、接口契约、算例和验证证据构成，有效正文不少于一百页。英文参考文献题名和源码标识不算
英文正文，也不能用英文附录、重复表格、空白页或超大字号补页。最终 PDF 使用 XeLaTeX 编译，并经 Poppler 全页渲染抽检。

\subsection{完成判据}
只有同时满足代码可链接、专项测试通过、HTTP/GUI E2E 通过、中文手册超过一百页、PDF 无布局缺陷、文档承诺零悬空和依赖
边界已记录，才可称可靠性理论功能闭环。任一验收项缺失都必须在状态文档中列为阻断项，不得以局部单元测试通过替代整体完成。

% 第 9 章：验证与边界局限
\section{验证与边界局限}\label{sec:rel-verify}

\subsection{注册测试}
\begin{center}\footnotesize
\begin{tabular}{@{}L{5.2cm}L{4.0cm}L{5.2cm}@{}}
\toprule
测试目标 & 注册位置 & 覆盖\\
\midrule
\srcpath{test\_reliability\_resolver} & \srcpath{tests/CMakeLists.txt} & MC、FMEA、故障模式 FMEA、解析器、数据策略、信息物理\\
\srcpath{test\_three\_stage\_reliability} & \srcpath{tests/CMakeLists.txt} & 三阶段隔离/重构/修复与指标\\
\srcpath{test\_intelligent\_cyber\_physical\_reliability} & \srcpath{tests/CMakeLists.txt} & CT/PT、保护配合、设备序列、信息共享依赖、在线 DAE、DER-FRT、三级对照\\
\srcpath{test\_resilience\_assessment} & \srcpath{tests/CMakeLists.txt} & 交叉调用 \srcpath{run\_sequential\_mc}/\srcpath{run\_distribution\_fmea}\\
\srcpath{reliability\_workflow\_e2e} & GUI E2E & FLISR 三向敏感性、日历年频率、物理割集、三阶段、在线 DAE 与三级对照\\
\bottomrule
\end{tabular}
\end{center}
\srcpath{reliability\_workflow\_e2e} 的三向检查为方法内敏感性（负荷加压$\Rightarrow$EENS 增、
失自动化$\Rightarrow$EENS 增且保留时延与冻结控制分量、正被动故障率增$\Rightarrow$模式频率与
总 EENS 增），\textbf{不}主张不同后果引擎的 EENS 数值可直接互比。

\subsection{边界与局限（如实标注）}
\begin{itemize}
  \item \textbf{网络后果口径}：MC/FMEA 在纯交流上为
        \srcpath{"ac-only-dcopf"}；混合系统为
        \srcpath{"hybrid-acdc-network-lp"}，此时
        \fld{validity.dc\_load\_curtailment\_included}=true、\fld{vsc\_dc\_power\_flow\_modelled}=true。
        两种口径的 \fld{ac\_voltage\_reactive\_feasibility\_certified}=false；
  \item \textbf{重要度语义}：\fld{critical\_components.importance} 是共现/归因秩（多故障态整体切负荷
        归给每个停运元件，可相加超 100\%），\textbf{非} Birnbaum 边际重要度 $\partial\mathrm{EENS}/\partial U_c$；
  \item \textbf{收敛判据}：CoV 仅对 $\bar S>0$ 有定义，零损失运行跑满迭代且记录零 CoV 不代表收敛；
  \item \textbf{尾部风险}：非序贯 VaR/CVaR 为 iid 小时近似，低估天气/修复/储能持续性引致的离散度；
  \item \textbf{信息物理分层}：L1 标量接口矩阵保留为快速筛选；L2 对照入口执行共享通信拓扑、QoS、
        共因环境、物理供电与电池自治、保护/FRT 轨迹和互斥联合事件树，保护误动进入总 EENS/LOLE/LOLF；
  \item \textbf{三阶段范围}：耦合 LinDistFlow（线性化损耗），LCC 与多端口能量路由器被拒
        （清 \fld{ok}）；VSC/DC-DC 恒效率（无固定/二次损耗）；发电机/变压器/换流器/开关故障为选入项；
  \item \textbf{在线保护边界}：在线入口执行 CT/PT、定/反时限、方向、mho/四边形距离、差动、主后备、
        断路器失灵、重合器/熔断器/分段器和 DER-FRT 闭环；其网络为正序、单相故障输入，不认证三相 EMT 或行波保护。
\end{itemize}

\subsection{独立方程 oracle：参数解析器交叉验证}\label{sec:rel-resolver-xref}
注册测试 \srcpath{reliability\_resolver\_cross\_validation} 把生产可靠性参数解析器
\srcpath{src/reliability/reliability\_assessment.cpp:resolve\_reliability\_params} 的规范化输出与一个
不链接 hacdcpf 的独立 Python oracle 对照。证据生成器
\srcpath{tools/reliability\_validation/validate\_reliability\_xref.cpp:emit\_case} 在覆盖全部转换分支的
八个确定性故障模式输入上运行解析器（$\lambda$+MTTR 的运行/日历口径、传统 MTBF 的两种约定、显式 MTTF、
含与不含修复时间的 FOR、以及主动按需），oracle \srcpath{tools/reliability\_validation/run\_cross\_validation.py}
依据 Billinton \& Allan 交替更新闭式独立重算 $U=\lambda/(\lambda+\mu)$（$\mu=H/r$）、MTBF/MTTF$\to\lambda$、
$\lambda=f/((1-f)r)\,H$、日历基改正 $\lambda=f_{cal}/(1-f_{cal}r/H)$ 与主动等效 $\lambda_{act}=\nu_d p_d$，
并逐字段核对解析器输出。

复现命令为
\begin{verbatim}
ctest --test-dir build/macos-release \
  -R '^reliability_resolver_cross_validation$' --output-on-failure
\end{verbatim}

准入阈值为 $10^{-9}$；本轮八个算例六类参数（$\lambda$、修复时间、不可用度、MTTF、日历频率、主动等效频率）
的最大误差为 $0$（独立重算与生产解析器在同一 IEEE~754 双精度下逐位一致），负控制（人为扰动 FOR 分支
$\lambda$ 的 $1\%$）如期触发失败。该 oracle 校核确定性参数解析闭式；系统级 EENS/SAIDI 的 MC 与 FMEA
引擎仍由各自套件验证，本节不以参数解析器一致性冒充系统级外部准确度。

\subsection{证据维护}
本手册为对照 2026-08-17 检出的实现修订；行号引用如与后续代码变更不符，以
\srcpath{include/hacdcpf/reliability/}、\srcpath{include/hacdcpf/analysis/three\_stage\_reliability.hpp}
与注册测试为准。数学模型对照见
\srcpath{docs/theory/reliability\_mathematical\_models\_and\_intelligent\_cyber\_physical\_assessment.md}、
\srcpath{docs/guides/reliability\_calculation\_workflow\_and\_case\_guide.md}、
\srcpath{docs/theory/reliability\_assessment\_models.md}。


\section{跨模块工业级技术评价基线}

本章规定本手册所述模块进入工程算例、批量研究或生产服务前必须共同满足的评价基线。
具体模型公式、变量、约束和专用算法以正文相应章节为准；本章不扩大源码已经实现的范围。

\subsection{输入、输出、边界条件与数据结构审计}
输入审计应逐项检查有限性、单位、上下界、枚举值和引用完整性。系统对象必须先按所需层级完成
校验；AC 与 DC 母线采用域限定映射，组件稳定 \texttt{index}、向量位置和临时图索引不得混用。
输出审计必须记录单位、索引空间、模型范围、收敛或可行状态、后端、容差、迭代/节点计数和告警。
空系统、空时域、孤岛、重复 ID、零容量、退化约束与不可用可选后端均属于边界测试集。

\subsection{工程假设与数值稳定性分析}
模型使用的平衡、线性化、稳态、独立性、概率分布、时间离散或网络等值假设必须与应用场景匹配。
数值评价至少包含变量缩放、绝对/相对残差、可行性违反、条件数或枢轴质量、步长/阻尼、迭代停滞
和随机种子复现性。若采用正则化、平滑、回退、松弛或近似，应同时报告触发条件、参数值、对主指标
的敏感性以及是否改变模型口径；不得用“已返回结果”替代收敛或有效性证明。

\subsection{复杂度与资源分析}
令 $n_x$、$n_c$、$n_t$、$n_s$、$|V|$、$|E|$ 分别表示变量、约束、时段、场景、节点和边数。
报告应分别给出建模、求解、后处理的时间与内存。图遍历通常为 $O(|V|+|E|)$；逐时段/逐场景
流程至少线性增长为 $O(n_t n_s)$ 次子问题调用；稠密线性代数最坏可达 $O(n_x^3)$，稀疏方法则应
报告结构非零数、填充率和因子分解峰值内存；MILP/NLP 不给出虚假的多项式最坏界，而报告节点数、
间隙、时限和实测扩展曲线。复杂度结论必须基于固定硬件、固定后端和可复现算例族。

\subsection{异常工况处理}
非法输入应在进入求解器前返回结构化错误；不可行、无界、不收敛、时限、内存不足和后端缺失必须
相互区分。部分结果只有在对应 \fld{ValidityFlags}、\fld{model\_scope}、
\fld{model\_limitations} 或等价字段允许时才可使用。异常路径不得静默丢弃 AC/DC 资产、制造平衡
节点、把 NaN 改写为零或把近似解标为全局最优；系统原始对象应保持可恢复和可重试。

\subsection{与其他模块的接口关系}
接口链路遵循“工程语义模型 $\to$ 校验 $\to$ 投影/装配 $\to$ 求解或分析 $\to$ 结果回投影”。
跨模块传递时保留稳定 ID、单位、符号、时间基准和场景身份；消费潮流、OPF、动态或可靠性结果前，
先检查其收敛、覆盖范围和有效性标志。HTTP/JSON/Python 层只做语义保持适配，不重新解释求解结果。

\subsection{测试与验证建议}
最低验证矩阵包括：手算小例、标准算例、边界/异常输入、随机性质测试、算法或后端交叉校核、
序列化往返、稳定 ID 回投影、AC/DC 同号母线、重复运行确定性和规模扩展测试。数值算法还应加入
残差独立重算、雅可比有限差分或自动微分校核；近似模型应与高保真模型比较并给出误差分布，而非
只报告单个平均值。回归报告必须列明构建配置、算例版本、容差、通过/失败/跳过数及未验证范围。

\subsection{工业级评估指标}
\begin{itemize}
  \item \textbf{正确性}：守恒残差、约束最大违反、解析/基准误差、结果回投影一致性；
  \item \textbf{鲁棒性}：收敛率、不可行识别率、异常分类准确率、扰动敏感性与随机复现率；
  \item \textbf{性能}：端到端时延、吞吐量、峰值内存、并行效率与规模增长斜率；
  \item \textbf{可追溯性}：输入模式版本、稳定 ID 覆盖率、模型范围/限制字段完整率、告警可定位率；
  \item \textbf{工程有效性}：与标准、外部引擎或现场数据的偏差，以及误判/漏判对业务指标的影响。
\end{itemize}
指标应预先给出验收阈值和失败处置；未达到阈值时只能标记为实验性或受限适用。

\subsection{适用范围、局限性与后续改进}
本手册只评价当前源码覆盖的模型、后端和数据结构，不对未建模物理过程、未安装求解器或未执行的
外部验证作保证。后续改进应按“缺口 $\to$ 可量化指标 $\to$ 源码/测试变更 $\to$ 独立验证”闭环，
优先消除会造成错误结论的语义缺口，其次提升数值鲁棒性与性能，最后扩展模型范围。


\end{document}
